\documentclass{article}
\usepackage{readmasters}
\begin{document}

\origpage{181}

\section*{Theorie der algebraischen Functionen einer Veränderlichen.}

(Von den Herren \emph{R.\ Dedekind} in Braunschweig und \emph{H.\ Weber} in Königsberg.)

\section*{Einleitung.}

Die im Nachstehenden mitgetheilten Untersuchungen verfolgen den Zweck, die Theorie der algebraischen Functionen einer Veränderlichen, welche eines der Hauptergebnisse der \emph{Riemann}schen Schöpfung ist, von einem einfachen und zugleich strengen und völlig allgemeinen Gesichtspunkt aus zu begründen. Bei den bisherigen Untersuchungen über diesen Gegenstand werden in der Regel gewisse beschränkende Voraussetzungen über die Singularitäten der betrachteten Functionen gemacht, und die sogenannten Ausnahmefälle entweder als Grenzfälle beiläufig erwähnt oder auch ganz bei Seite gesetzt. Ebenso werden gewisse Grundsätze über die Stetigkeit und Entwickelbarkeit zugelassen, deren Evidenz sich auf geometrische Anschauung verschiedener Art stützt. Eine sichere Basis für die Grundvorstellungen sowie für eine allgemeine und ausnahmslose Behandlung der Theorie lässt sich gewinnen, wenn man von einer Verallgemeinerung der Theorie der rationalen Functionen einer Veränderlichen, ins Besondere des Satzes, dass jede ganze rationale Function einer Veränderlichen sich in lineare Factoren zerlegen lässt, ausgeht. Diese Verallgemeinerung ist einfach und bekannt in dem ersten Falle, in welchem die von \emph{Riemann} mit $p$ bezeichnete Zahl (das Geschlecht nach \emph{Clebsch}) den Werth Null hat. Für den allgemeinen Fall, welcher sich zu dem eben genannten ähnlich verhält, wie der Fall der allgemeinsten algebraischen Zahlen zu demjenigen der rationalen Zahlen, wiesen die mit bestem Erfolge in der Zahlentheorie angewandten Methoden, die sich an \emph{Kummer}s Schöpfung der idealen Zahlen anschliessen, und der

\origpage{182}
Uebertragung auf die Theorie der Functionen fähig sind, auf den richtigen Weg${}^{*)}$.

Versteht man, analog der Zahlentheorie, unter einem \emph{Körper algebraischer Functionen} ein System solcher Functionen von der Beschaffenheit, dass die Anwendung der vier Species auf Functionen des Systems immer zu Functionen desselben Systems führt, so deckt sich dieser Begriff vollständig mit dem der \emph{Riemann}schen Classe algebraischer Functionen. Unter den Functionen eines solchen Körpers kann eine beliebige als unabhängige Veränderliche und die übrigen als von ihr abhängig betrachtet werden. Für jede dieser „Darstellungsweisen“ ergiebt sich ein System von Functionen des Körpers, die als \emph{ganze Functionen} zu bezeichnen sind, deren Quotienten den ganzen Körper erschöpfen. Unter diesen ganzen Functionen lassen sich nun wieder Gruppen von Functionen aussondern, welchen die charakteristischen Merkmale solcher ganzen rationalen Functionen zukommen, die einen gemeinschaftlichen Theiler haben. Ein solcher Theiler existirt zwar im allgemeinen Falle nicht, wenn man aber die bezüglichen Sätze über rationale Functionen nicht an den Theiler selbst, sondern an das System der durch denselben theilbaren Functionen knüpft, so gestatten sie eine vollkommene Uebertragung auf die allgemeinen algebraischen Functionen. Auf diese Weise gelangt man zu dem Begriff des \emph{Ideals}, ein Name, der aus \emph{Kummer}s zahlentheoretischen Arbeiten stammt, wo die nicht existirenden Theiler als „ideale Theiler“ in die Rechnung eingeführt werden.

Obwohl es sich in der vorliegenden Arbeit keineswegs um „ideale“ Functionen handelt, sondern alle Operationen nur an Systemen wirklich existirender Functionen ausgeführt werden, schien es doch zweckmässig, den Namen „Ideal“, der in der Zahlentheorie bereits gebräuchlich ist, beizubehalten.

Mit diesen Idealen lässt sich nach gehöriger Erklärung der Multiplication ganz nach denselben Regeln rechnen, wie mit rationalen Functionen.

\textbf{*)} Die idealen Zahlen sind von \emph{Kummer} zuerst eingeführt durch die Abhandlung: \emph{Zur Theorie der complexen Zahlen} (\emph{Crelle}s Journal, Bd.~35); eine weitere Fortführung und eine allgemeine Darstellung der Theorie der algebraischen Zahlen findet man in der zweiten und dritten Auflage von \emph{Dirichlet}s Vorlesungen über Zahlentheorie, sowie in der Abhandlung von \emph{Dedekind}: \emph{Sur la théorie des nombres entiers algébriques} (Paris 1877. Abdruck aus dem Bulletin des Sciences math.\ et astron.\ von \emph{Darboux} und \emph{Hoüel}). Die Kenntniss dieser Schriften wird aber in unserer Arbeit nirgends vorausgesetzt.

Aus mündlichen Mittheilungen ist uns jetzt bekannt geworden, dass bereits vor Jahren \emph{Kronecker} mit Beziehung auf die Arbeiten von \emph{Weierstrass} Untersuchungen angestellt hat, die auf derselben Grundlage, wie die unsrigen, beruhen.

\origpage{183}
Ins Besondere ergiebt sich der Satz, dass jedes Ideal auf eine einzige Weise in Factoren zerlegbar ist, welche selbst nicht weiter zerlegt werden können und daher \emph{Primideale} genannt werden. Diese Primideale entsprechen den linearen Factoren in der Theorie der ganzen rationalen Functionen. Auf Grund derselben gelangt man zu einer völlig präcisen und allgemeinen Definition des „Punktes der \emph{Riemann}schen Fläche“, d.\ h. eines vollkommen bestimmten Systems von Zahlwerthen, welche man den Functionen des Körpers widerspruchslos beilegen kann.

Eine darauf gegründete formale Definition des Differentialquotienten führt sodann zu der Geschlechtszahl und zu einer ganz allgemeinen, eleganten Darstellung der Differentiale erster Gattung. Hieran schliesst sich der Beweis des \emph{Riemann-Roch}schen Satzes über die Anzahl der willkürlichen Constanten in einer durch ihre Unendlichkeitspunkte bestimmten Function, und die Theorie der Differentiale zweiter und dritter Gattung. Bis zu diesem Punkte kommt die Stetigkeit und Entwickelbarkeit der untersuchten Functionen in keiner Weise in Betracht. Es würde z.\ B. nirgends eine Lücke bleiben, wenn man das Gebiet der benutzten Zahlen auf das System der algebraischen Zahlen beschränken wollte. Dadurch wird ein wohl abgegrenzter und ziemlich umfassender Theil der Theorie der algebraischen Functionen lediglich durch die seiner eigenen Sphäre angehörigen Mittel behandelt.

Freilich ergeben sich alle diese Resultate durch einen weit geringeren Aufwand von Mitteln und als Specialfälle einer vielumfassenden Allgemeinheit aus \emph{Riemann}s Theorie; allein es ist bekannt, dass diese Theorie bezüglich einer strengen Begründung noch gewisse Schwierigkeiten bietet, und bis es gelungen ist, diese Schwierigkeiten vollständig zu überwinden, dürfte der von uns betretene Weg oder wenigstens ein verwandter, wohl der einzige sein, der für die Theorie der algebraischen Functionen mit befriedigender Strenge und Allgemeinheit zum Ziele führt. So würde sich die Theorie der Ideale selbst ausserordentlich vereinfachen, wenn man den Begriff der \emph{Riemann}schen Fläche und ins Besondere den eines Punktes derselben sammt den auf die Stetigkeit der algebraischen Functionen gegründeten Anschauungen voraussetzen wollte. In unserer Arbeit ist umgekehrt auf einem langen Umwege die Theorie der Ideale algebraisch begründet und aus dieser eine vollkommen präcise und strenge Definition des „Punktes der \emph{Riemann}schen Fläche“ gewonnen, welche auch als Basis für die Untersuchung der Stetigkeit und der damit zusammenhängenden Fragen dienen kann. Diese Fragen,

\origpage{184}
wozu auch die auf die \emph{Abel}schen Integrale und die Periodicitätsmoduln bezüglichen gehören, bleiben von unserer Untersuchung einstweilen ausgeschlossen. Wir hoffen bei einer anderen Gelegenheit darauf zurückzukommen.

Königsberg, den 22. October 1880.

\section*{I. Abtheilung.}

\subsection*{§~1. Körper algebraischer Functionen.}

Eine Variable $\theta$ heisst eine \emph{algebraische Function} einer unabhängigen Veränderlichen $z$, wenn dieselbe einer irreductibeln algebraischen Gleichung
\[ F(\theta, z) = 0 \tag{1.} \]
genügt. $F$ bedeutet hierin einen Ausdruck von der Form
\[ F(\theta, z) = a_{0}\theta^{n} + a_{1}\theta^{n-1} + \cdots + a_{n-1}\theta + a_{n}, \]
worin die Coefficienten $a_{0}, a_{1}, \ldots a_{n}$ ganze rationale Functionen von $z$ ohne gemeinschaftlichen Theiler sind. Die vorausgesetzte Irreductibilität der Gleichung (1.) involvirt, dass $\theta$ nicht einer Gleichung niedrigeren Grades in Bezug auf $\theta$ genügt, und, wie sich aus dem Algorithmus des grössten gemeinschaftlichen Theilers ergiebt, wenn
\[ G(\theta, z) = b_{0}\theta^{m} + b_{1}\theta^{m-1} + \cdots + b_{m-1}\theta + b_{m} = 0 \]
eine zweite Gleichung ist, welcher $\theta$ genügt, dass $G(\theta, z)$ durch $F(\theta, z)$ algebraisch theilbar sein muss. Es lässt sich nun nachweisen, dass $G(\theta, z)$ auch in Bezug auf $z$ nicht von niedrigerem Grade sein kann als $F(\theta, z)$ und nur dann vom selben Grade, wenn sich aus $G(\theta, z)$ ein von $z$ unabhängiger Factor absondern lässt. Nehmen wir an, die Coefficienten $b_{0}, b_{1}, \ldots b_{m}$ seien von gemeinschaftlichen Factoren befreit, und bezeichnen wir mit
\[ H(\theta, z) = c_{0}\theta^{m-n} + c_{1}\theta^{m-n-1} + \cdots + c_{m-n} \]
den vom Nenner befreiten Quotienten von $G$ durch $F$, so ist
\[ kG(\theta, z) = F(\theta, z).H(\theta, z), \]
worin $k$ eine ganze rationale Function von $z$ ist, und die Vergleichung der Coefficienten ergiebt
\[ \begin{aligned} kb_{0} &= a_{0}c_{0}, \\ kb_{1} &= a_{0}c_{1} + a_{1}c_{0}, \\ kb_{2} &= a_{0}c_{2} + a_{1}c_{1} + a_{2}c_{0}, \\ &\cdots\cdots\cdots\cdots \end{aligned} \]

\origpage{185}
worin die $c_{0}, c_{1}, \ldots c_{m-n}$ gleichfalls ohne gemeinschaftlichen Theiler vorausgesetzt werden können.

Hieraus folgt zunächst, dass $k$ constant sein muss, und $=1$ gesetzt werden kann; denn ist durch irgend einen Linearfactor von $k$ $a_{0}, a_{1}, \ldots a_{r-1}, c_{0}, c_{1}, \ldots c_{s-1}$ theilbar, $a_{r}$, $c_{s}$ nicht theilbar, so folgt aus
\[ kb_{r+s} = \cdots a_{r-1}c_{s+1} + a_{r}c_{s} + a_{r+1}c_{s-1} + \cdots \]
der Widerspruch, dass $a_{r}c_{s}$ durch denselben Linearfactor theilbar sein müsste. Hieraus aber folgt weiter, dass der Grad von $G(\theta, z)$ in Bezug auf $z$ gleich ist der Summe der Grade von $F$ und $H$ in Bezug auf $z$; denn sind $a_{r}$, $c_{s}$ die ersten unter den Coefficienten $a$, $c$, deren Grad den Maximalwerth erreicht, so folgt wieder aus
\[ b_{r+s} = \cdots a_{r-1}c_{s+1} + a_{r}c_{s} + a_{r+1}c_{s-1} + \cdots, \]
dass der Grad von $b_{r+s}$ gleich der Summe der Grade von $a_{r}$ und $c_{s}$ ist.

Dividirt man die Gleichung (1.) durch $a_{0}$, so kann dieselbe auch in die Form gesetzt werden
\[ f(\theta, z) = \theta^{n} + b_{1}\theta^{n-1} + \cdots + b_{n-1}\theta + b_{n} = 0, \tag{2.} \]
worin die Coefficienten $b_{1}, b_{2}, \ldots b_{n}$ auch gebrochene rationale Functionen von $z$ sein können.

Das System aller rationalen Functionen von $\theta$ und $z$, $\Phi(\theta, z)$, hat die Eigenschaft, dass seine Individuen sich durch die elementaren Rechenoperationen, Addition, Subtraction, Multiplication und Division reproduciren, und dies System wird daher als ein \emph{Körper algebraischer Functionen $\Omega$ vom Grade $n$} bezeichnet. Ist zunächst $\varphi(\theta)$ eine \emph{ganze} rationale Function von $\theta$, deren Coefficienten rational von $z$ abhängen, so kann man durch algebraische Division zwei eben solche Functionen $q(\theta)$, $r(\theta)$ bestimmen, von denen die zweite den Grad $n-1$ nicht übersteigt, so dass
\[ \varphi(\theta) = q(\theta)f(\theta) + r(\theta) \]
oder wegen (2.)
\[ \varphi(\theta) = r(\theta). \]
Ist $\varphi(\theta)$ durch $f(\theta)$ nicht theilbar, so haben diese beiden Functionen (wegen der vorausgesetzten Irreductibilität von $f(\theta)$) keinen Theiler gemein, und daher lassen sich durch die Methode des grössten gemeinschaftlichen Theilers zwei Functionen $f_{1}(\theta)$, $\varphi_{1}(\theta)$ so bestimmen, dass
\[ f(\theta)f_{1}(\theta) + \varphi(\theta)\varphi_{1}(\theta) = 1, \]

\origpage{186}
also wegen (2.)
\[ \varphi_{1}(\theta) = \frac{1}{\varphi(\theta)}. \]
Aus diesen beiden Bemerkungen, zusammengenommen mit der Voraussetzung der Irreductibilität von $f(\theta)$ ergiebt sich der folgende

\emph{Lehrsatz.} \emph{Jede Function $\zeta$ des Körpers $\Omega$ lässt sich auf eine einzige Weise in die Form setzen:}
\[ \zeta = x_{0} + x_{1}\theta + \cdots + x_{n-1}\theta^{n-1}, \]
\emph{worin die Coefficienten $x_{0}, x_{1}, \ldots x_{n-1}$ rationale Functionen von $z$ sind. Umgekehrt gehört jede Function dieser Form selbstverständlich dem Körper $\Omega$ an.}

Wählt man unter den Functionen des Körpers $\Omega$ $n$ beliebige aus:
\[ \begin{aligned} \eta_{1} &= x_{0}^{(1)} + x_{1}^{(1)}\theta + \cdots + x_{n-1}^{(1)}\theta^{n-1}, \\ \eta_{2} &= x_{0}^{(2)} + x_{1}^{(2)}\theta + \cdots + x_{n-1}^{(2)}\theta^{n-1}, \\ &\cdots\cdots\cdots\cdots \\ \eta_{n} &= x_{0}^{(n)} + x_{1}^{(n)}\theta + \cdots + x_{n-1}^{(n)}\theta^{n-1}, \end{aligned} \]
jedoch so, dass die Determinante
\[ \Sigma \pm x_{0}^{(1)} x_{1}^{(2)} \ldots x_{n-1}^{(n)} \]
nicht identisch Null ist, so ergiebt sich, dass jede Function des Körpers $\Omega$ auch in der Form dargestellt werden kann
\[ \zeta = y_{1}\eta_{1} + y_{2}\eta_{2} + \cdots + y_{n}\eta_{n}, \]
deren Coefficienten $y_{1}, y_{2}, \ldots y_{n}$ rationale Functionen von $z$ sind. Ein solches System von Functionen $\eta_{1}, \eta_{2}, \ldots \eta_{n}$ soll eine \emph{Basis des Körpers $\Omega$} heissen.

Damit ein Functionensystem $\eta_{1}, \eta_{2}, \ldots \eta_{n}$ des Körpers $\Omega$ eine Basis desselben bilde, ist erforderlich und hinreichend, dass zwischen ihnen keine Gleichung (Identität) von der Form
\[ y_{1}\eta_{1} + y_{2}\eta_{2} + \cdots + y_{n}\eta_{n} = 0 \]
bestehe, in welcher die Coefficienten $y_{1}, y_{2}, \ldots y_{n}$ nicht sämmtlich verschwinden. Beispielsweise bilden die Functionen $1, \theta, \theta^{2}, \ldots \theta^{n-1}$ eine Basis von $\Omega$.

\subsection*{§~2. Normen, Spuren und Discriminanten.}

Wählt man zur Darstellung der Functionen von $\Omega$ eine beliebige Basis $\eta_{1}, \eta_{2}, \ldots \eta_{n}$, so kann man, wenn $\zeta$ irgend eine Function in $\Omega$ bedeutet, setzen:

\origpage{187}
\[ \left\{ \begin{aligned} \zeta\eta_{1} &= y_{1,1}\eta_{1} + y_{1,2}\eta_{2} + \cdots + y_{1,n}\eta_{n}, \\ \zeta\eta_{2} &= y_{2,1}\eta_{1} + y_{2,2}\eta_{2} + \cdots + y_{2,n}\eta_{n}, \\ &\cdots\cdots\cdots\cdots \\ \zeta\eta_{n} &= y_{n,1}\eta_{1} + y_{n,2}\eta_{2} + \cdots + y_{n,n}\eta_{n}, \end{aligned} \right. \tag{1.} \]
worin die Coefficienten $y_{\iota,\iota'}$ rationale Functionen von $z$ sind. Hieraus ergiebt sich die Gleichung
\[ \begin{vmatrix} y_{1,1}-\zeta, & y_{1,2}, & \ldots & y_{1,n} \\ y_{2,1}, & y_{2,2}-\zeta, & \ldots & y_{2,n} \\ \cdots & \cdots & \cdots & \cdots \\ y_{n,1}, & y_{n,2}, & \ldots & y_{n,n}-\zeta \end{vmatrix} = 0, \tag{2.} \]
welche, nach Potenzen von $\zeta$ geordnet, die Gestalt hat
\[ \varphi(\zeta) = \zeta^{n} + b_{1}\zeta^{n-1} + \cdots + b_{n-1}\zeta + b_{n} = 0 \tag{3.} \]
und, wie sich aus dem Multiplicationssatz der Determinanten ohne Schwierigkeit ergiebt, von der Wahl der zu Grunde gelegten Basis $\eta_{1}, \eta_{2}, \ldots \eta_{n}$ ganz unabhängig ist. Von den Coefficienten $b_{1}, b_{2}, \ldots b_{n}$ der Function $\varphi$, welche sämmtlich rationale Functionen von $z$ und durch die Function $\zeta$ vollkommen bestimmt sind, sollen zwei, die in der Folge von Wichtigkeit sind, durch besondere Namen ausgezeichnet werden. Die Function
\[ (-1)^{n}b_{n} = \begin{vmatrix} y_{1,1}, & y_{1,2}, & \ldots & y_{1,n} \\ y_{2,1}, & y_{2,2}, & \ldots & y_{2,n} \\ \cdots & \cdots & \cdots & \cdots \\ y_{n,1}, & y_{n,2}, & \ldots & y_{n,n} \end{vmatrix} \tag{4.} \]
heisst die \emph{Norm} der Function $\zeta$, und wird mit $N(\zeta)$ bezeichnet. Von dieser Function gelten folgende Sätze.

1. Die einzige Function, deren Norm identisch verschwindet, ist die Function „Null“; denn macht man in dem System (1.) die Annahme, dass $N(\zeta) = 0$ sei, so folgt, dass sich ein System rationaler Functionen von $z$, die nicht sämmtlich verschwinden, $y_{1}, y_{2}, \ldots y_{n}$ so bestimmen lässt, dass
\[ \zeta(y_{1}\eta_{1} + y_{2}\eta_{2} + \cdots + y_{n}\eta_{n}) = 0, \]
also, da $\eta_{1}, \eta_{2}, \ldots \eta_{n}$ eine Basis von $\Omega$ bilden, $\zeta = 0$.

2. Die Norm einer rationalen Function von $z$ ist die $n^{\text{te}}$ Potenz dieser Function. Denn ist $\zeta$ rational, so reduciren sich die Gleichungen (1.) auf die Identitäten $\zeta\eta_{h} = \zeta\eta_{h}$, woraus $N(\zeta) = \zeta^{n}$ folgt.

\origpage{188}

3. Ist $\zeta'$ irgend eine zweite Function des Körpers $\Omega$ und das dem System (1.) entsprechende Gleichungssystem für diese Function
\[ \zeta'\eta_{h} = \overset{\iota}{\Sigma} y'_{h,\iota}\eta_{\iota}, \]
so folgt:
\[ \zeta\zeta'\eta_{h} = \overset{\iota,\iota'}{\Sigma} y'_{h,\iota}y_{\iota,\iota'}\eta_{\iota'} \]
\ednote{Over the second $y$ the print shows a small malformed mark; it is not clear from the print whether it is a prime or stray ink. Set unprimed, as the composition of the two systems requires.}
und daraus nach dem Multiplicationssatz der Determinanten
\[ N(\zeta\zeta') = N(\zeta)N(\zeta'). \]

4. Aus 2. und 3. folgt:
\[ N(\zeta)N\Bigl(\frac{1}{\zeta}\Bigr) = 1, \]
also:
\[ N\Bigl(\frac{\zeta}{\zeta'}\Bigr) = \frac{N(\zeta)}{N(\zeta')}. \]

5. Endlich ergiebt sich aus der Definition der Function $\varphi$, (2.), (3.) der wichtige Satz: Ist $t$ eine beliebige Constante (oder auch eine rationale Function von $z$), so ist
\[ \varphi(t) = N(t-\zeta). \]

Es soll sodann die Function
\[ -b_{1} = y_{1,1} + y_{2,2} + \cdots + y_{n,n} \tag{5.} \]
die \emph{Spur} von $\zeta$ genannt und mit $S(\zeta)$ bezeichnet werden. Für diese ergeben sich unmittelbar aus der Definition die Sätze:
\[ S(0) = 0, \tag{6.} \]
\[ S(1) = n. \tag{7.} \]

Und wenn $x$ eine rationale Function von $z$, ferner $\zeta$, $\zeta'$ zwei Functionen in $\Omega$ bedeuten:
\[ S(x\zeta) = xS(\zeta), \tag{8.} \]
\[ S(\zeta+\zeta') = S(\zeta) + S(\zeta'). \tag{9.} \]

Es hat sich aus dieser Betrachtung ergeben, \emph{dass jede Function $\zeta$ in $\Omega$ einer Gleichung $n^{\text{ten}}$ Grades, $\varphi(\zeta) = 0$, genügt, deren Coefficienten rational von $z$ abhängen}. Wenn diese Gleichung irreductibel ist, so bilden die Functionen $1, \zeta, \zeta^{2}, \ldots \zeta^{n-1}$ eine Basis von $\Omega$. Im andern Falle sei
\[ \varphi_{1}(\zeta) = \zeta^{e} + b'_{0}\zeta^{e-1} + \cdots + b'_{e-1}\zeta + b'_{e} = 0 \tag{10.} \]
\ednote{Printed $b'_{0}$ as the coefficient of $\zeta^{e-1}$; the monic equation and the remaining coefficients $b'_{e-1}$, $b'_{e}$ require $b'_{1}$. An apparent misprint.}
die Gleichung niedrigsten Grades, deren Coefficienten in $z$ rational sind, welcher die Function $\zeta$ genügt, und mithin $\varphi_{1}(\zeta) = 0$ irreductibel, $e < n$.

\origpage{189}
Da gleichwohl $\varphi(\zeta)$ verschwindet, so muss $\varphi(\zeta)$ durch $\varphi_{1}(\zeta)$ algebraisch theilbar sein, und wie in §~1 ergiebt sich, dass jede rationale Function $\eta$ von $z$ und $\zeta$ in der Form darstellbar ist
\[ \eta = x_{0} + x_{1}\zeta + \cdots + x_{e-1}\zeta^{e-1}, \]
deren Coefficienten $x_{0}, x_{1}, \ldots x_{e-1}$ rational von $z$ abhängen${}^{*)}$. Ist nun
\[ \theta^{f} + \eta_{1}\theta^{f-1} + \cdots + \eta_{f-1}\theta + \eta_{f} = 0 \]
die Gleichung niedrigsten Grades, welcher $\theta$ genügt, deren Coefficienten rational von $z$ \emph{und} $\zeta$ abhängen, so besteht zwischen den $e.f$ Functionen
\[ \zeta^{h}\theta^{k} \qquad {\scriptstyle (h = 0,\, 1,\, \ldots\, e-1;\ k = 0,\, 1,\, \ldots\, f-1)} \tag{11.} \]
keine lineare Gleichung mit rational von $z$ abhängigen Coefficienten, während \emph{jede} Function in $\Omega$ \emph{linear} mit rational von $z$ abhängigen Coefficienten durch diese Functionen darstellbar ist. Es ergiebt sich daraus, dass dieselben eine Basis von $\Omega$ bilden, und dass sonach
\[ e.f = n, \]
also $e$ ein Theiler von $n$ ist.

Wendet man die Basis (11.) zur Aufstellung der Norm von $\zeta$ an, so erkennt man leicht mittels der Gleichung (10.), dass
\[ N(\zeta) = ((-1)^{e}b'_{e})^{f} = (-1)^{n}b_{e}^{\prime f} \]
wird. Da ferner für ein beliebiges constantes $t$ die Function $\zeta-t$ einer Gleichung von demselben Grade genügt wie $\zeta$, so ergiebt sich der Satz:

10.\ednote{Printed 10. although the preceding propositions of this § are numbered 1. to 5.; 6. expected. An apparent misprint.} \emph{Die Function $\varphi(t)$ (3.) ist entweder irreductibel oder eine ganze Potenz einer irreductibeln Function.}

Ist $\eta_{1}, \eta_{2}, \ldots \eta_{n}$ ein beliebiges System von $n$ Functionen in $\Omega$, gleichviel ob dasselbe eine Basis bildet oder nicht, so führen wir eine zu diesem System gehörige rationale Function von $z$ ein, die wir als dessen \emph{Discriminante}, $\Delta(\eta_{1}, \eta_{2}, \ldots \eta_{n})$ bezeichnen und folgendermassen definiren
\[ \Delta(\eta_{1}, \eta_{2}, \ldots \eta_{n}) = \begin{vmatrix} S(\eta_{1}\eta_{1}), & S(\eta_{1}\eta_{2}), & \ldots & S(\eta_{1}\eta_{n}) \\ S(\eta_{2}\eta_{1}), & S(\eta_{2}\eta_{2}), & \ldots & S(\eta_{2}\eta_{n}) \\ \cdots & \cdots & \cdots & \cdots \\ S(\eta_{n}\eta_{1}), & S(\eta_{n}\eta_{2}), & \ldots & S(\eta_{n}\eta_{n}) \end{vmatrix}. \tag{12.} \]
\emph{Die Discriminante ist dann und nur dann nicht identisch Null, wenn die Functionen $\eta_{1}, \eta_{2}, \ldots \eta_{n}$ eine Basis von $\Omega$ bilden.}

\textbf{*)} Aus der Gleichung $\varphi_{1}(\zeta) = 0$ entspringt ein Körper algebraischer Functionen $\Omega_{1}$ vom Grade $e$, dessen Functionen sämmtlich zugleich im Körper $\Omega$ enthalten sind, und der daher als ein Theiler des Körpers $\Omega$ bezeichnet werden kann.

\origpage{190}

Um den ersten Theil dieser Behauptung zu beweisen, nehmen wir an, es sei $\Delta(\eta_{1}, \eta_{2}, \ldots \eta_{n}) = 0$. Es lässt sich unter dieser Voraussetzung ein System rationaler Functionen $y_{1}, y_{2}, \ldots y_{n}$ von $z$, die nicht alle identisch verschwinden, so bestimmen, dass
\[ \begin{gathered} y_{1}S(\eta_{1}\eta_{k}) + y_{2}S(\eta_{2}\eta_{k}) + \cdots + y_{n}S(\eta_{n}\eta_{k}) = S(\eta_{k}(y_{1}\eta_{1} + y_{2}\eta_{2} + \cdots + y_{n}\eta_{n})) = 0. \\ {\scriptstyle (k = 1,\, 2,\, \ldots\, n)} \end{gathered} \]
Wählt man daher ein System rationaler Functionen $x_{1}, x_{2}, \ldots x_{n}$ von $z$, ganz beliebig und setzt:
\[ \begin{aligned} y_{1}\eta_{1} + y_{2}\eta_{2} + \cdots + y_{n}\eta_{n} &= \eta, \\ x_{1}\eta_{1} + x_{2}\eta_{2} + \cdots + x_{n}\eta_{n} &= \xi, \end{aligned} \]
so folgt:
\[ S(\xi\eta) = 0. \]
Wenn aber die Functionen $\eta_{1}, \eta_{2}, \ldots \eta_{n}$ eine Basis von $\Omega$ bilden, so kann $\xi$ jede beliebige Function in $\Omega$, also, da $\eta$ nicht verschwindet, beispielsweise auch $\dfrac{1}{\eta}$ sein. Dann ist aber die letzte Gleichung sicher nicht erfüllt, und es kann also unter dieser Voraussetzung die Discriminante von $\eta_{1}, \eta_{2}, \ldots \eta_{n}$ \emph{nicht} identisch verschwinden.

Halten wir die Annahme fest, dass $\eta_{1}, \eta_{2}, \ldots \eta_{n}$ eine Basis von $\Omega$ sei, und setzen:
\[ \eta'_{k} = x_{1,k}\eta_{1} + x_{2,k}\eta_{2} + \cdots + x_{n,k}\eta_{n}, \qquad {\scriptstyle (k = 1,\, 2,\, \ldots\, n)} \]
so bilden die Functionen $\eta'_{1}, \eta'_{2}, \ldots \eta'_{n}$ eine Basis von $\Omega$ oder nicht, je nachdem die Determinante der rationalen Functionen $x_{h,k}$ von $z$
\[ X = \Sigma \pm x_{1,1}x_{2,2} \ldots x_{n,n} \]
von Null verschieden ist oder nicht. Nun ist aber
\[ S(\eta'_{h}\eta'_{k}) = \overset{\iota,\iota'}{\underset{1,n}{\Sigma}} x_{\iota,h}x_{\iota',k}S(\eta_{\iota}\eta_{\iota'}), \]
und daraus ergiebt sich nach dem Multiplicationssatz der Determinanten der \emph{Hauptsatz über die Discriminanten}
\[ \Delta(\eta'_{1}, \eta'_{2}, \ldots \eta'_{n}) = X^{2}\Delta(\eta_{1}, \eta_{2}, \ldots \eta_{n}), \tag{13.} \]
woraus auch die Richtigkeit des zweiten Theils der obigen Behauptung erhellt, dass die Discriminante eines Functionensystems stets dann identisch verschwindet, wenn dasselbe keine Basis von $\Omega$ bildet.

\origpage{191}

\subsection*{§~3. Das System der ganzen Functionen von $z$ im Körper $\Omega$.}

\emph{Definition.} Eine Function $\omega$ des Körpers $\Omega$ soll eine \emph{ganze} Function von $z$ heissen, wenn in der Gleichung niedrigsten Grades, welcher dieselbe nach §~2 genügt:
\[ \varphi(\omega) = \omega^{e} + b_{1}\omega^{e-1} + \cdots + b_{e-1}\omega + b_{e} = 0, \tag{1.} \]
die Coefficienten $b_{1}, b_{2}, \ldots b_{e}$ \emph{ganze rationale} Functionen von $z$ sind; im entgegengesetzten Fall heisse sie eine \emph{gebrochene} Function. Der Inbegriff aller ganzen Functionen von $z$ in $\Omega$ soll mit $\mathfrak{o}$ bezeichnet werden. Da nach §~2 $N(t-\omega)$ eine ganze Potenz von $\varphi(t)$ ist, so folgt, dass für eine ganze Function $\omega$ auch die sämmtlichen Coefficienten von $N(t-\omega)$ ganze rationale Functionen von $z$ sind, also ins Besondere:

1. Die Norm und die Spur einer ganzen Function sind ganze rationale Functionen von $z$.

Aus der Definition der ganzen Functionen ergiebt sich ferner:

2. Eine rationale Function von $z$ gehört dann und nur dann zu dem System $\mathfrak{o}$, wenn sie eine ganze rationale Function von $z$ ist.

3. Jede Function $\eta$ in $\Omega$ kann durch Multiplication mit einer von Null verschiedenen ganzen rationalen Function von $z$ in eine Function des Systems $\mathfrak{o}$ verwandelt werden. Denn es genügt $\eta$ nach §~2 einer Gleichung niedrigsten Grades von der Form
\[ b_{0}\eta^{e} + b_{1}\eta^{e-1} + \cdots + b_{e-1}\eta + b_{e} = 0, \]
deren Coefficienten ganze rationale Functionen von $z$ sind, und diese geht durch die Substitution $b_{0}\eta = \omega$ in eine Gleichung von der Form (1.) für $\omega$ über.

4. Eine Function $\omega$ des Körpers $\Omega$, welche irgend einer Gleichung von der Form genügt
\[ \psi(\omega) = \omega^{m} + c_{1}\omega^{m-1} + \cdots + c_{m-1}\omega + c_{m} = 0, \]
in welcher die Coefficienten $c_{1}, \ldots c_{m}$ ganze rationale Functionen von $z$ sind, ist eine \emph{ganze Function}. Denn ist
\[ \varphi(\omega) = \omega^{e} + b_{1}\omega^{e-1} + \cdots + b_{e-1}\omega + b_{e} = 0 \]
die Gleichung niedrigsten Grades, welcher $\omega$ genügt, so muss $\psi(\omega)$ durch $\varphi(\omega)$ algebraisch theilbar sein:
\[ \psi(\omega) = \varphi(\omega)\chi(\omega), \]

\origpage{192}
was, wie leicht zu zeigen ist, zur Folge hat, dass auch die Coefficienten von $\varphi(\omega)$ und $\chi(\omega)$ ganze rationale Functionen von $z$ sind (\emph{Gauss}, Disq.\ Ar.\ art.~42). Hieraus ergiebt sich der Hauptsatz über die ganzen Functionen:

5. \emph{Summe, Differenz, Product zweier ganzen Functionen sind wieder ganze Functionen.}

Sind nämlich $\omega'$, $\omega''$ zwei ganze Functionen in $\Omega$, welche resp.\ den Gleichungen genügen
\[ \begin{aligned} \omega'^{n'} + b'_{1}\omega'^{n'-1} + \cdots + b'_{n'-1}\omega' + b'_{n'} &= 0, \\ \omega''^{n''} + b''_{1}\omega''^{n''-1} + \cdots + b''_{n''-1}\omega'' + b''_{n''} &= 0, \end{aligned} \]
so kann man, wenn man unter $\omega_{1}, \omega_{2}, \ldots \omega_{m}$ die $m = n'n''$ Producte
\[ \omega'^{h'}\omega''^{h''} \qquad {\scriptstyle (h' = 0,\, 1,\, \ldots\, n'-1;\ h'' = 0,\, 1,\, \ldots\, n''-1)} \]
und unter $\omega$ eine der drei Functionen $\omega' \pm \omega''$, $\omega'\omega''$ versteht, setzen:
\[ \begin{aligned} \omega\omega_{1} &= x_{1,1}\omega_{1} + \cdots + x_{1,m}\omega_{m}, \\ &\cdots\cdots\cdots\cdots \\ \omega\omega_{m} &= x_{m,1}\omega_{1} + \cdots + x_{m,m}\omega_{m}, \end{aligned} \]
worin die $x_{h,h'}$ ganze rationale Functionen von $z$ sind, und daraus erhält man
\[ \begin{vmatrix} x_{1,1}-\omega, & x_{1,2}, & \ldots & x_{1,m} \\ x_{2,1}, & x_{2,2}-\omega, & \ldots & x_{2,m} \\ \cdots & \cdots & \cdots & \cdots \\ x_{m,1}, & x_{m,2}, & \ldots & x_{m,m}-\omega \end{vmatrix} = 0, \]
also eine Gleichung für $\omega$, deren Coefficienten ganze rationale Functionen von $z$ sind.

Als Corollar ergiebt sich hieraus, dass jede ganze rationale Function von Functionen in $\mathfrak{o}$ selbst eine Function des Systems $\mathfrak{o}$ ist.

6. Eine ganze Function $\omega$ heisst durch eine andere ganze Function $\omega'$ \emph{theilbar}, wenn eine dritte ganze Function $\omega''$ existirt, welche der Bedingung genügt
\[ \omega = \omega'\omega''. \]
Aus dieser Definition ergiebt sich sofort:

Ist $\omega$ theilbar durch $\omega'$, $\omega'$ durch $\omega''$, so ist auch $\omega$ durch $\omega''$ theilbar.

Ist $\omega'$ und $\omega''$ durch $\omega$ theilbar, so ist auch $\omega' \pm \omega''$ durch $\omega$ theilbar, und allgemein: sind $\omega_{1}, \omega_{2}, \omega_{3}, \ldots$ durch $\omega$ theilbar, $\omega'_{1}, \omega'_{2}, \omega'_{3}, \ldots$ beliebige Functionen in $\mathfrak{o}$, so ist auch $\omega'_{1}\omega_{1} + \omega'_{2}\omega_{2} + \omega'_{3}\omega_{3} + \cdots$ durch $\omega$ theilbar.

\origpage{193}

7. Bilden die Functionen $\eta_{1}, \eta_{2}, \ldots \eta_{n}$ eine Basis von $\Omega$, so kann man (nach 3.) $n$ von Null verschiedene ganze rationale Functionen von $z$, $a_{1}, a_{2}, \ldots a_{n}$ der Art bestimmen, dass
\[ \omega_{1} = a_{1}\eta_{1}, \qquad \omega_{2} = a_{2}\eta_{2}, \qquad \ldots \qquad \omega_{n} = a_{n}\eta_{n} \]
ganze Functionen sind, und diese bilden ebenfalls eine Basis von $\Omega$, da
\[ \Delta(\omega_{1}, \omega_{2}, \ldots \omega_{n}) = a_{1}^{2}a_{2}^{2}\ldots a_{n}^{2}\Delta(\eta_{1}, \eta_{2}, \ldots \eta_{n}) \]
von Null verschieden ist. Es giebt also Basen von $\Omega$, $\omega_{1}, \omega_{2}, \ldots \omega_{n}$, welche aus lauter ganzen Functionen bestehen, und die Discriminante einer solchen Basis ist, da $S(\omega_{r}\omega_{s})$ ganze rationale Functionen von $z$ sind, selbst eine von Null verschiedene ganze rationale Function von $z$. Jede Function von der Form
\[ \omega = x_{1}\omega_{1} + x_{2}\omega_{2} + \cdots + x_{n}\omega_{n}, \tag{2.} \]
in welcher die $x_{1}, x_{2}, \ldots x_{n}$ ganze rationale Functionen von $z$ sind, gehört dann zu dem System $\mathfrak{o}$; aber es ist durchaus nicht nothwendig, dass umgekehrt jede Function in $\mathfrak{o}$ in dieser Form darstellbar sei.

Nehmen wir also an, es existiren in $\mathfrak{o}$ noch andere Functionen als die in der Form (2.) enthaltenen, so müssen sich eine lineare Function $z-c$ und gewisse ganze rationale Functionen $x_{1}, x_{2}, \ldots x_{n}$, die nicht alle durch $z-c$ theilbar sind, so wählen lassen, dass
\[ \frac{x_{1}\omega_{1} + x_{2}\omega_{2} + \cdots + x_{n}\omega_{n}}{z-c} \]
eine ganze Function ist. Die Functionen $x_{1}, x_{2}, \ldots x_{n}$ lassen sich nun auf ihre nicht sämmtlich verschwindenden constanten Reste $c_{1}, c_{2}, \ldots c_{n}$ in Bezug auf $z-c$ reduciren, und daraus erhellt, dass auch
\[ \omega = \frac{c_{1}\omega_{1} + c_{2}\omega_{2} + \cdots + c_{n}\omega_{n}}{z-c} \]
eine ganze Function ist. Ist $c_{1}$ von Null verschieden, so bilden auch die $n$ ganzen Functionen
\[ \omega \quad \text{und} \quad \omega_{2}, \omega_{3}, \ldots \omega_{n} \]
eine Basis von $\Omega$ und zugleich ist nach §~2 (13.)
\[ \Delta(\omega, \omega_{2}, \ldots \omega_{n}) = \frac{c_{1}^{2}}{(z-c)^{2}}\Delta(\omega_{1}, \omega_{2}, \ldots \omega_{n}), \]
also von niedrigerem Grade als $\Delta(\omega_{1}, \omega_{2}, \ldots \omega_{n})$. Da nun diese beiden Discriminanten ganze rationale Functionen von $z$ sind, so gelangt man durch wiederholte Anwendung dieses Verfahrens schliesslich zu einer aus ganzen

\origpage{194}
Functionen bestehenden Basis von $\Omega$, $\omega'_{1}, \omega'_{2}, \ldots \omega'_{n}$, deren Discriminante im Grade nicht weiter erniedrigt werden kann, und welche folglich die Eigenschaft hat, \emph{dass jede Function $\omega$ in $\mathfrak{o}$ in der Form enthalten ist}
\[ \omega = x_{1}\omega'_{1} + x_{2}\omega'_{2} + \cdots + x_{n}\omega'_{n} \]
\emph{mit ganzen rationalen Functionen von $z$ als Coefficienten. Ein solches System soll eine Basis von $\mathfrak{o}$ genannt werden.}

Ist $\omega_{1}, \omega_{2}, \ldots \omega_{n}$ eine Basis von $\mathfrak{o}$ und
\[ \omega'_{\iota} = x_{\iota,1}\omega_{1} + x_{\iota,2}\omega_{2} + \cdots + x_{\iota,n}\omega_{n}, \qquad {\scriptstyle (\iota = 1,\, 2,\, \ldots\, n)} \]
so wird das System $\omega'_{1}, \omega'_{2}, \ldots \omega'_{n}$ dann und nur dann ebenfalls eine Basis von $\mathfrak{o}$ bilden, wenn die Determinante der ganzen rationalen Functionen $x_{\iota,\iota'}$
\[ X = \Sigma \pm x_{1,1}x_{2,2}\ldots x_{n,n} \]
eine von Null verschiedene \emph{Constante} ist. Denn nehmen wir an, es habe diese Determinante irgend einen Linearfactor $z-c$, so lassen sich Constanten $c_{1}, c_{2}, \ldots c_{n}$, nicht sämmtlich verschwindend, so bestimmen, dass die $n$ ganzen rationalen Functionen von $z$
\[ c_{1}x_{1,\iota} + c_{2}x_{2,\iota} + \cdots + c_{n}x_{n,\iota} \]
durch $z-c$ theilbar werden (d.\ h.\ für $z = c$ verschwinden); dann aber ist
\[ \frac{c_{1}\omega'_{1} + c_{2}\omega'_{2} + \cdots + c_{n}\omega'_{n}}{z-c} \]
eine ganze Function und mithin $\omega'_{1}, \omega'_{2}, \ldots \omega'_{n}$ keine Basis von $\mathfrak{o}$.

Da nun andererseits
\[ \Delta(\omega'_{1}, \omega'_{2}, \ldots \omega'_{n}) = X^{2}\Delta(\omega_{1}, \omega_{2}, \ldots \omega_{n}) \]
ist, so folgt, dass die Discriminante einer Basis von $\mathfrak{o}$ von einem constanten Factor abgesehen von der Wahl dieser Basis unabhängig ist. Man erhält also eine vollkommen bestimmte ganze rationale Function von $z$, wenn man in der Discriminante einer beliebigen Basis von $\mathfrak{o}$ den Coefficienten der höchsten Potenz von $z$ durch Division $= 1$ macht. \emph{Diese Function soll die Discriminante des Körpers $\Omega$ oder des Systems $\mathfrak{o}$ genannt und mit $\Delta(\Omega)$ oder $\Delta(\mathfrak{o})$ bezeichnet werden.}

\subsection*{§~4. Die Functionenmoduln.}

Wir betrachten im Folgenden Systeme von Functionen, welche wir \emph{Functionenmoduln} oder auch schlechtweg \emph{Moduln} nennen und folgendermassen

\origpage{195}
definiren. \emph{Ein Functionensystem (in $\Omega$) heisst ein Modul, wenn sich die Functionen desselben durch Addition, Subtraction und durch Multiplication mit ganzen rationalen Functionen von $z$ reproduciren.}

Bezeichnet man mit $\alpha_{1}, \alpha_{2}, \ldots \alpha_{m}$ irgend $m$ gegebene Functionen, mit $x_{1}, x_{2}, \ldots x_{m}$ willkürliche ganze rationale Functionen von $z$, so bildet der Inbegriff aller Functionen von der Form
\[ \alpha = x_{1}\alpha_{1} + x_{2}\alpha_{2} + \cdots + x_{m}\alpha_{m} \]
einen Modul. Ein solcher soll ein \emph{endlicher Modul} genannt und mit
\[ \mathfrak{a} = [\alpha_{1}, \alpha_{2}, \ldots \alpha_{m}] \]
bezeichnet werden. Das Functionensystem $\alpha_{1}, \alpha_{2}, \ldots \alpha_{m}$ heisst die Basis dieses Moduls.

Wir wollen ein Functionensystem $\alpha_{1}, \alpha_{2}, \ldots \alpha_{m}$ \emph{rational irreductibel} oder die Functionen $\alpha_{1}, \alpha_{2}, \ldots \alpha_{p}$\ednote{Printed $\alpha_{p}$ instead of $\alpha_{m}$; an apparent misprint.} \emph{rational unabhängig} nennen, wenn eine Gleichung von der Form
\[ x_{1}\alpha_{1} + x_{2}\alpha_{2} + \cdots + x_{m}\alpha_{m} = 0 \]
für rationale $x$ nur dann bestehen kann, wenn $x_{1} = 0$, $x_{2} = 0$, $\ldots$ $x_{m} = 0$ ist. Ein Functionensystem, welches eine Basis des Körpers $\Omega$ bildet, ist daher stets rational irreductibel, und es giebt kein System von mehr als $n$ rational unabhängigen Functionen in $\Omega$.

Wir beweisen nun zunächst den Satz:

1. \emph{Jeder endliche Modul besitzt eine rational irreductible Basis.}

Der Beweis desselben ergiebt sich unmittelbar aus dem folgenden Hülfssatz:

Sind die ganzen rationalen Functionen $y_{1,1}, y_{2,1}, \ldots y_{m,1}$ ohne gemeinschaftlichen Theiler, so lassen sich andere ganze rationale Functionen $y_{1,2}, y_{2,2}, \ldots y_{m,m}$ so bestimmen, dass
\[ \Sigma \pm y_{1,1}y_{2,2}\ldots y_{m,m} = 1{}^{*)}. \]

\textbf{*)} Der Satz ist richtig und bekannt für $m = 2$. Nehmen wir also an, er sei bewiesen für $m-1$, so können wir, wenn $\partial$ den grössten gemeinschaftlichen Theiler von $y_{1,1}, y_{2,1}, \ldots y_{m-1,1}$ bedeutet, der Gleichung genügen
\[ \begin{vmatrix} y_{1,1}, & y_{2,1}, & \ldots & y_{m-1,1} \\ y_{1,3}, & y_{2,3}, & \ldots & y_{m-1,3} \\ \cdots & \cdots & \cdots & \cdots \\ y_{1,m}, & y_{2,m}, & \ldots & y_{m-1,m} \end{vmatrix} = \partial \]
und wenn wir also die ganzen rationalen Functionen $x$, $y$ so bestimmen, dass
\[ xy_{m,1} - y\partial = (-1)^{m-1} \]
ist, so folgt:
\[ \begin{vmatrix} y_{1,1}, & y_{2,1}. & \ldots & y_{m-1,1}, & y_{m,1} \\ \frac{xy_{1,1}}{\partial}, & \frac{xy_{2,1}}{\partial}, & \ldots & \frac{xy_{m-1,1}}{\partial}, & y \\ y_{1,3}, & y_{2,3}, & \ldots & y_{m-1,3}, & 0 \\ \cdots & \cdots & \cdots & \cdots & \cdots \\ y_{1,m}, & y_{2,m}, & \ldots & y_{m-1,m}, & 0 \end{vmatrix} = 1. \]\ednote{Printed a period after $y_{2,1}$ in the first row where a comma is expected; possibly a comma whose tail failed to ink.}

\origpage{196}

Genügen nun die Functionen $\alpha_{1}, \alpha_{2}, \ldots \alpha_{m}$ einer Gleichung
\[ \overset{\iota}{\underset{1,m}{\Sigma}} y_{\iota,1}\alpha_{\iota} = 0, \]
in welcher die ganzen rationalen Functionen $y_{1,1}, \ldots y_{m,1}$ ohne gemeinschaftlichen Theiler angenommen werden können, so setze man
\[ \begin{aligned} \overset{\iota}{\underset{1,m}{\Sigma}} y_{\iota,2}\alpha_{\iota} &= \beta_{2}, \\ &\cdots\cdots\cdots\cdots \\ \overset{\iota}{\underset{1,m}{\Sigma}} y_{\iota,m}\alpha_{\iota} &= \beta_{m}; \end{aligned} \]
dann ist der Modul $[\alpha_{1}, \alpha_{2}, \ldots \alpha_{m}]$ völlig identisch mit dem Modul $[\beta_{2}, \beta_{3}, \ldots \beta_{m}]$, dessen Basis eine Function weniger enthält. Sind die Functionen $\beta_{\iota}$ noch nicht rational unabhängig, so kann man sie in derselben Weise weiter reduciren, und gelangt schliesslich, falls die Functionen $\alpha_{\iota}$ nicht sämmtlich verschwinden (ein Fall, welchen wir von dem Modulbegriff ganz ausschliessen wollen) zu einer irreductibeln Basis. Wir werden in der Folge unter einer Basis schlechtweg stets eine irreductible Basis verstehen.

2. Obwohl man nach dem Vorhergehenden für einen und denselben Modul sehr verschiedene irreductible Basen auffinden kann, so ist doch die Zahl der Functionen, die in einer solchen enthalten sind, stets dieselbe, da im entgegengesetzten Fall dasjenige Functionensystem, welches mehr Functionen enthält, nicht rational irreductibel sein könnte. Sind also $\alpha_{1}, \alpha_{2}, \ldots \alpha_{m}$; $\beta_{1}, \beta_{2}, \ldots \beta_{m}$ zwei irreductible Basen desselben Moduls $\mathfrak{a}$, so ist, da sowohl die $\alpha_{k}$ als die $\beta_{k}$ in $\mathfrak{a}$ enthalten sind:
\[ \alpha_{k} = \overset{\iota}{\underset{1,m}{\Sigma}} p_{k}^{(\iota)}\beta_{\iota}; \qquad \beta_{k} = \overset{\iota}{\underset{1,m}{\Sigma}} q_{\iota}^{(k)}\alpha_{\iota}, \]
worin die Coefficienten $p$, $q$ ganze rationale Functionen von $z$ sind. Hieraus aber folgt:
\[ \overset{\iota}{\underset{1,m}{\Sigma}} q_{\iota}^{(k)}p_{\iota}^{(h)} = 0 \quad \text{oder} \quad 1, \]

\origpage{197}
je nachdem $h$ von $k$ verschieden ist oder nicht, und daraus:
\[ \Sigma \pm p_{1}^{(1)}p_{2}^{(2)}\ldots p_{m}^{(m)}.\Sigma \pm q_{1}^{(1)}q_{2}^{(2)}\ldots q_{m}^{(m)} = 1, \]
und da beide Determinanten ganze rationale Functionen von $z$ sind, so müssen sie beide \emph{constant} sein.

3. \emph{Definition.} Ein Modul $\mathfrak{a}$ heisst durch einen Modul $\mathfrak{b}$ \emph{theilbar}, oder $\mathfrak{b}$ ein \emph{Theiler (Divisor)} von $\mathfrak{a}$, $\mathfrak{a}$ ein \emph{Vielfaches (Multiplum)} von $\mathfrak{b}$ ($\mathfrak{b}$ geht in $\mathfrak{a}$ auf), wenn jede Function in $\mathfrak{a}$ zugleich in $\mathfrak{b}$ enthalten ist. $\mathfrak{b}$ soll ein \emph{echter Theiler} von $\mathfrak{a}$ heissen, wenn $\mathfrak{a}$ durch $\mathfrak{b}$ theilbar, aber nicht mit $\mathfrak{b}$ identisch ist${}^{*)}$.

Aus dieser Definition ergiebt sich sofort:

Ist $\mathfrak{a}$ theilbar durch $\mathfrak{b}$, $\mathfrak{b}$ theilbar durch $\mathfrak{c}$, so ist auch $\mathfrak{a}$ theilbar durch $\mathfrak{c}$.

4. \emph{Definition.} Der Inbegriff $\mathfrak{m}$ aller derjenigen Functionen, welche zugleich in zwei Moduln $\mathfrak{a}$, $\mathfrak{b}$ enthalten sind, bildet, falls er nicht aus der einzigen Function „Null“ besteht, einen Modul (nach der allgemeinen Definition), welcher das \emph{kleinste gemeinschaftliche Vielfache von $\mathfrak{a}$ und $\mathfrak{b}$} heisst, weil jeder Modul, welcher ein Vielfaches zugleich von $\mathfrak{a}$ und von $\mathfrak{b}$ ist, auch ein Vielfaches von $\mathfrak{m}$ ist. Das kleinste gemeinschaftliche Vielfache von einer beliebigen Zahl von Moduln $\mathfrak{a}, \mathfrak{b}, \mathfrak{c}, \ldots$ ist dem entsprechend der Inbegriff aller der Functionen, die zugleich in $\mathfrak{a}, \mathfrak{b}, \mathfrak{c}, \ldots$ enthalten sind. Man kann dasselbe bilden, indem man nach Belieben je zwei der Moduln $\mathfrak{a}, \mathfrak{b}, \mathfrak{c}, \ldots$ durch ihr kleinstes gemeinschaftliches Vielfache ersetzt.

5. \emph{Definition.} Ist $\alpha$ eine beliebige Function in $\mathfrak{a}$, $\beta$ eine beliebige Function in $\mathfrak{b}$, so bildet der Inbegriff aller Functionen von der Form $\alpha + \beta$ einen Modul $\mathfrak{d}$, welcher der \emph{grösste gemeinschaftliche Theiler der beiden Moduln $\mathfrak{a}$ und $\mathfrak{b}$ heisst}. Derselbe ist, wenn $\mathfrak{a}$ und $\mathfrak{b}$ endliche Moduln sind, selbst ein solcher. Ist nämlich
\[ \mathfrak{a} = [\alpha_{1}, \alpha_{2}, \ldots \alpha_{r}], \qquad \mathfrak{b} = [\beta_{1}, \beta_{2}, \ldots \beta_{s}], \]
so ist
\[ \mathfrak{d} = [\alpha_{1}, \alpha_{2}, \ldots \alpha_{r}, \beta_{1}, \beta_{2}, \ldots \beta_{s}]. \]
Nach der Definition der Theilbarkeit ist $\mathfrak{d}$ ein Theiler sowohl von $\mathfrak{a}$ als von $\mathfrak{b}$. Ist umgekehrt $\mathfrak{d}'$ ein Theiler von $\mathfrak{a}$ und von $\mathfrak{b}$, so sind die Functionen

\textbf{*)} Der Begriff der Theilbarkeit der Moduln ist der von den Zahlen her gewohnten Anschauung zuwider gebildet, insofern der Theiler einen grösseren Inhalt an Functionen enthält als das Vielfache.

\origpage{198}
$\alpha$ sowohl als die Functionen $\beta$, mithin auch die Functionen $\alpha + \beta$ in $\mathfrak{d}'$ enthalten; daher ist $\mathfrak{d}$ durch $\mathfrak{d}'$ theilbar.

Die Definition des grössten gemeinschaftlichen Theilers einer beliebigen Anzahl von Moduln ergiebt sich hiernach von selbst.

6. \emph{Definition.} Ist $\mathfrak{a}$ ein Modul, $\alpha$ jede Function in $\mathfrak{a}$ und $\mu$ eine beliebige Function in $\Omega$, so verstehen wir unter dem Product $\mu\mathfrak{a}$ oder $\mathfrak{a}\mu$ den Inbegriff aller Functionen $\mu\alpha$, welcher wieder ein Modul ist. Ist
\[ \mathfrak{a} = [\alpha_{1}, \alpha_{2}, \ldots \alpha_{r}] \]
ein endlicher Modul, so ist
\[ \mu\mathfrak{a} = [\mu\alpha_{1}, \mu\alpha_{2}, \ldots \mu\alpha_{r}], \]
also ebenfalls ein endlicher Modul, und aus $\mu\mathfrak{a} = \mu\mathfrak{b}$ folgt $\mathfrak{a} = \mathfrak{b}$, wenn $\mu$ von Null verschieden ist.

7. \emph{Definition.} Sind $\mathfrak{a}$, $\mathfrak{b}$ zwei Moduln, $\alpha$, $\beta$ sämmtliche Functionen in $\mathfrak{a}$, resp.\ in $\mathfrak{b}$, so verstehen wir unter dem \emph{Product}
\[ \mathfrak{a}\mathfrak{b} = \mathfrak{b}\mathfrak{a} = \mathfrak{c} \]
den Inbegriff aller Producte einer Function $\alpha$ und einer Function $\beta$ und aller Summen solcher Producte, also sämmtlicher Functionen, welche durch das Zeichen
\[ \gamma = \Sigma\alpha\beta \]
bezeichnet werden können.

Dieses Functionensystem bildet jederzeit einen Modul, und zwar einen endlichen, wenn $\mathfrak{a}$ und $\mathfrak{b}$ solche sind. Sind nämlich $\mathfrak{a}$ und $\mathfrak{b}$ so definirt, wie in 5., so bilden die $r.s$ Functionen $\alpha_{\iota}\beta_{\varkappa}$ eine, wenn auch reductible, Basis von $\mathfrak{c}$. Ein Product aus beliebig vielen Moduln $\mathfrak{a}, \mathfrak{b}, \mathfrak{c}, \ldots$ erklärt sich hiernach von selbst, und es gilt für dasselbe der Fundamentalsatz der Multiplication von der Vertauschbarkeit der Factoren. Sind die einzelnen Functionen eines solchen Products, deren Anzahl $m$ sei, einander gleich und $= \mathfrak{a}$, so wird dasselbe mit $\mathfrak{a}^{m}$ bezeichnet, und es ist
\[ \mathfrak{a}^{m+m'} = \mathfrak{a}^{m}\mathfrak{a}^{m'}. \]

Im Allgemeinen ist ein Product $\mathfrak{a}\mathfrak{b}$ \emph{nicht} durch $\mathfrak{a}$ theilbar. Dagegen gilt der Satz, dessen Beweis sich unmittelbar aus der Definition ergiebt:

\emph{Ist $\mathfrak{a}$ theilbar durch $\mathfrak{a}_{1}$, $\mathfrak{b}$ durch $\mathfrak{b}_{1}$, so ist $\mathfrak{a}\mathfrak{b}$ theilbar durch $\mathfrak{a}_{1}\mathfrak{b}_{1}$.}

8. \emph{Definition.} Unter dem Quotienten $\frac{\mathfrak{b}}{\mathfrak{a}}$ zweier Moduln $\mathfrak{a}$, $\mathfrak{b}$ soll der Inbegriff aller derjenigen Functionen $\gamma$ verstanden werden, welche die Eigenschaft

\origpage{199}
haben, dass $\gamma\mathfrak{a}$ durch $\mathfrak{b}$ theilbar ist. Dieser Quotient ist, falls er nicht aus der einzigen Function „Null“ besteht, ein Modul $\mathfrak{c}$, was sofort aus der Definition erhellt. Das Product $\frac{\mathfrak{b}}{\mathfrak{a}}\cdot\mathfrak{a}$ ist jederzeit durch $\mathfrak{b}$ theilbar, wenn auch nicht immer gleich $\mathfrak{b}$.

\subsection*{§~5. Congruenzen.}

Zwei Functionen $\alpha$, $\beta$ heissen \emph{congruent} nach dem Modul $\mathfrak{a}$
\[ \alpha \equiv \beta \ (\text{mod.}\,\mathfrak{a}), \]
wenn die Differenz der beiden Functionen, $\alpha-\beta$, in dem Modul $\mathfrak{a}$ enthalten ist.

Aus dieser Definition ergeben sich unmittelbar die folgenden Sätze:

1. Ist $\alpha \equiv \beta$, $\beta \equiv \gamma \ (\text{mod.}\,\mathfrak{a})$, so ist $\alpha \equiv \gamma \ (\text{mod.}\,\mathfrak{a})$.

2. Ist $\mathfrak{d}$ irgend ein Theiler von $\mathfrak{a}$, so folgt aus $\alpha \equiv \beta \ (\text{mod.}\,\mathfrak{a})$, dass auch $\alpha \equiv \beta \ (\text{mod.}\,\mathfrak{d})$ ist.

3. Ist $\alpha \equiv \beta \ (\text{mod.}\,\mathfrak{a})$, $\mu$ eine beliebige Function in $\Omega$, so folgt $\mu\alpha \equiv \mu\beta \ (\text{mod.}\,\mu\mathfrak{a})$, und umgekehrt folgt aus der letzteren Congruenz die erstere, wenn $\mu$ von Null verschieden.

4. Ist $\alpha \equiv \beta$, $\alpha_{1} \equiv \beta_{1} \ (\text{mod.}\,\mathfrak{a})$, so ist auch $\alpha \pm \alpha_{1} \equiv \beta \pm \beta_{1} \ (\text{mod.}\,\mathfrak{a})$.

Sind $\lambda_{1}, \lambda_{2}, \ldots \lambda_{m}$ beliebig gegebene Functionen in $\Omega$, $c_{1}, c_{2}, \ldots c_{m}$ \emph{willkürliche Constanten}, so heisst der Inbegriff aller Functionen von der Form
\[ c_{1}\lambda_{1} + c_{2}\lambda_{2} + \cdots + c_{m}\lambda_{m} \]
eine \emph{Schaar} und wird mit $(\lambda_{1}, \lambda_{2}, \ldots \lambda_{m})$ bezeichnet. Das Functionensystem $\lambda_{1}, \lambda_{2}, \ldots \lambda_{m}$ heisst \emph{die Basis der Schaar}. Die Functionen $\lambda_{1}, \lambda_{2}, \ldots \lambda_{m}$ heissen \emph{linear unabhängig} oder ihr System \emph{linear irreductibel}, wenn eine Gleichung (Identität) von der Form
\[ c_{1}\lambda_{1} + c_{2}\lambda_{2} + \cdots + c_{m}\lambda_{m} = 0 \]
nicht anders bestehen kann, als wenn die constanten Coefficienten $c_{1}, c_{2}, \ldots c_{m}$ alle verschwinden.

Hiernach gilt der Satz, \emph{dass jede Schaar eine linear irreductible Basis besitzt}. Denn ist $c_{1}\lambda_{1} + c_{2}\lambda_{2} + \cdots + c_{m}\lambda_{m} = 0$ und $c_{1}$ von Null verschieden, so ist die Schaar $(\lambda_{1}, \lambda_{2}, \ldots \lambda_{m})$ identisch mit der Schaar $(\lambda_{2}, \lambda_{3}, \ldots \lambda_{m})$, deren Basis eine Function weniger enthält. Ist diese noch nicht linear irreductibel, so kann man auf die gleiche Weise weiterschliessen. Auch hier soll in

\origpage{200}
der Folge unter einer Basis schlechtweg eine irreductible Basis verstanden sein. Die Anzahl der Functionen, welche in einer irreductiblen Basis einer Schaar enthalten sind, ist stets dieselbe und heisst die \emph{Dimension} der Schaar. Ist $m$ die Dimension, so heisst die Schaar auch eine \emph{$m$-fache}. Irgend $m$ Functionen einer solchen Schaar bilden eine irreductible Basis derselben dann und nur dann, wenn sie linear unabhängig sind.

Die Functionen $\lambda_{1}, \lambda_{2}, \ldots \lambda_{m}$ heissen \emph{linear unabhängig in Bezug auf den Modul} $\mathfrak{a}$, wenn eine Congruenz von der Form
\[ c_{1}\lambda_{1} + c_{2}\lambda_{2} + \cdots + c_{m}\lambda_{m} \equiv 0 \ (\text{mod.}\,\mathfrak{a}) \]
für keine anderen als verschwindende constante Coefficienten $c_{1}, c_{2}, \ldots c_{m}$ besteht. Zwei Summen von der Form $\Sigma c_{\iota}\lambda_{\iota}$ mit verschiedenen Werthen der constanten Coefficienten $c_{\iota}$ sind dann auch stets incongruent nach dem Modul $\mathfrak{a}$.

Es seien nun $\mathfrak{a}$ und $\mathfrak{b}$ zwei Moduln, und es werde \emph{zunächst} angenommen, es existiren in $\mathfrak{b}$ nur eine endliche Anzahl von Functionen $\lambda_{1}, \lambda_{2}, \ldots \lambda_{m}$, welche nach dem Modul $\mathfrak{a}$ linear unabhängig sind. Jede Function $\beta$ in $\mathfrak{b}$ genügt dann einer und nur einer Congruenz von der Form
\[ \beta \equiv c_{1}\lambda_{1} + c_{2}\lambda_{2} + \cdots + c_{m}\lambda_{m} \ (\text{mod.}\,\mathfrak{a}) \]
mit constanten Coefficienten $c_{1}, c_{2}, \ldots c_{m}$. Die Schaar $(\lambda_{1}, \lambda_{2}, \ldots \lambda_{m})$ kann daher ein \emph{vollständiges Restsystem des Moduls $\mathfrak{b}$ nach dem Modul} $\mathfrak{a}$ und $\lambda_{1}, \lambda_{2}, \ldots \lambda_{m}$ eine \emph{Basis} desselben genannt werden, und man kann in symbolischer Bezeichnung setzen:
\[ \mathfrak{b} \equiv (\lambda_{1}, \lambda_{2}, \ldots \lambda_{m}) \ (\text{mod.}\,\mathfrak{a}). \]

Wählt man in $\mathfrak{b}$ irgend ein System von $m$ Functionen $\lambda'_{1}, \lambda'_{2}, \ldots \lambda'_{m}$ aus, so gelten $m$ Congruenzen
\[ \lambda'_{h} \equiv \overset{\iota}{\Sigma} k_{h,\iota}\lambda_{\iota} \ (\text{mod.}\,\mathfrak{a}) \]
mit constanten $k_{h,\iota}$, und dies System bildet dann und nur dann eine Basis eines vollständigen Restsystems von $\mathfrak{b}$ nach $\mathfrak{a}$, wenn die Determinante
\[ \Sigma \pm k_{1,1}k_{2,2}\ldots k_{m,m} \]
von Null verschieden ist.

\subsection*{§~6. Norm eines Moduls in Bezug auf einen andern.}

Ist $(\lambda_{1}, \lambda_{2}, \ldots \lambda_{m})$ ein beliebiges vollständiges Restsystem eines Moduls $\mathfrak{b}$ in Bezug auf einen andern $\mathfrak{a}$, so ergiebt sich, weil $z\mathfrak{b}$ durch $\mathfrak{b}$

\origpage{201}
theilbar ist, ein ganz bestimmtes System von $m^{2}$ Constanten $c_{h,k}$, durch welches die Congruenzen erfüllt werden:
\[ \left. \begin{aligned} z\lambda_{1} &\equiv c_{1,1}\lambda_{1} + c_{2,1}\lambda_{2} + \cdots + c_{m,1}\lambda_{m} \\ z\lambda_{2} &\equiv c_{1,2}\lambda_{1} + c_{2,2}\lambda_{2} + \cdots + c_{m,2}\lambda_{m} \\ &\cdots\cdots\cdots\cdots \\ z\lambda_{m} &\equiv c_{1,m}\lambda_{1} + c_{2,m}\lambda_{2} + \cdots + c_{m,m}\lambda_{m} \end{aligned} \right\} (\text{mod.}\,\mathfrak{a}) \]
und durch Auflösung dieses Systems erkennt man, dass jede Function $\lambda_{\iota}$, und mithin jede Function $\beta$ des Moduls $\mathfrak{b}$ durch Multiplication mit der ganzen rationalen Function $m^{\text{ten}}$ Grades von $z$
\[ (\mathfrak{b}, \mathfrak{a}) = (-1)^{m} \begin{vmatrix} c_{1,1}-z, & c_{2,1}, & \ldots & c_{m,1} \\ c_{1,2}, & c_{2,2}-z, & \ldots & c_{m,2} \\ \cdots & \cdots & \cdots & \cdots \\ c_{1,m}, & c_{2,m}, & \ldots & c_{m,m}-z \end{vmatrix} \]
in eine Function des Moduls $\mathfrak{a}$ verwandelt wird. Diese Function $(\mathfrak{b}, \mathfrak{a})$ ist, wie sich aus dem Multiplicationssatz der Determinanten leicht ergiebt, von der Wahl der Basis $\lambda_{1}, \lambda_{2}, \ldots \lambda_{m}$ unabhängig, also nur von den beiden Moduln $\mathfrak{a}$, $\mathfrak{b}$ abhängig und soll die \emph{Norm von $\mathfrak{a}$ in Bezug auf} $\mathfrak{b}$ genannt werden.

Ist jede Function in $\mathfrak{b}$ zugleich in $\mathfrak{a}$ enthalten, also $\mathfrak{b}$ durch $\mathfrak{a}$ theilbar, so ist $m = 0$ und $(\mathfrak{b}, \mathfrak{a}) = 1$ zu setzen. Wenn dagegen $\mathfrak{b}$ \emph{nicht}, wie oben angenommen, eine endliche Anzahl in Bezug auf $\mathfrak{a}$ linear unabhängiger Functionen enthält, dann soll festgesetzt werden, dass $(\mathfrak{b}, \mathfrak{a}) = 0$ sei.

1. Ist $\mathfrak{m}$ das kleinste gemeinschaftliche Vielfache, $\mathfrak{d}$ der grösste gemeinschaftliche Theiler von $\mathfrak{a}$ und $\mathfrak{b}$, so ist jede Congruenz zwischen Functionen des Moduls $\mathfrak{b}$ in Bezug auf den Modul $\mathfrak{a}$ vollkommen gleichbedeutend mit der Congruenz derselben Functionen nach dem Modul $\mathfrak{m}$; andererseits ist jede Function in $\mathfrak{b}$ einer Function in $\mathfrak{d}$ und umgekehrt jede Function in $\mathfrak{d}$ einer Function in $\mathfrak{b}$ congruent nach dem Modul $\mathfrak{a}$. Aus diesen Bemerkungen ergiebt sich sofort der wichtige Satz:
\[ (\mathfrak{b}, \mathfrak{a}) = (\mathfrak{b}, \mathfrak{m}) = (\mathfrak{d}, \mathfrak{a}), \]
welcher auch richtig bleibt, wenn $(\mathfrak{b}, \mathfrak{a}) = 0$ ist.

2. \emph{Ist der Modul $\mathfrak{a}$ theilbar durch den Modul $\mathfrak{b}$, dieser durch den dritten Modul $\mathfrak{c}$, so ist}
\[ (\mathfrak{c}, \mathfrak{a}) = (\mathfrak{c}, \mathfrak{b})(\mathfrak{b}, \mathfrak{a}). \]
Der Satz ist offenbar richtig, wenn von den beiden Normen $(\mathfrak{c}, \mathfrak{b})$, $(\mathfrak{b}, \mathfrak{a})$

\origpage{202}
eine verschwindet. Ist dies nicht der Fall und ist
\[ \begin{aligned} \mathfrak{c} &\equiv (\varrho_{1}, \varrho_{2}, \ldots \varrho_{r}) \ (\text{mod.}\,\mathfrak{b}), \\ \mathfrak{b} &\equiv (\lambda_{1}, \lambda_{2}, \ldots \lambda_{s}) \ (\text{mod.}\,\mathfrak{a}), \end{aligned} \]
so sind die Functionen $\varrho_{1}, \varrho_{2}, \ldots \varrho_{r}, \lambda_{1}, \lambda_{2}, \ldots \lambda_{s}$ zusammengenommen linear unabhängig nach dem Modul $\mathfrak{a}$; denn ist
\[ \overset{\iota}{\Sigma} c_{\iota}\varrho_{\iota} + \overset{\iota}{\Sigma} c'_{\iota}\lambda_{\iota} \equiv 0 \quad (\text{mod.}\,\mathfrak{a}), \]
so folgt, da $\mathfrak{a}$ durch $\mathfrak{b}$ theilbar ist und die Functionen $\lambda_{\iota}$ in $\mathfrak{b}$ enthalten sind,
\[ \begin{aligned} \overset{\iota}{\Sigma} c_{\iota}\varrho_{\iota} &\equiv 0 \ (\text{mod.}\,\mathfrak{b}), \quad \text{mithin} \quad c_{\iota} = 0, \\ \overset{\iota}{\Sigma} c'_{\iota}\lambda_{\iota} &\equiv 0 \ (\text{mod.}\,\mathfrak{a}), \quad \text{mithin} \quad c'_{\iota} = 0. \end{aligned} \]
Da ferner jede Function $\gamma$ in $\mathfrak{c}$ einer Congruenz von der Form genügt:
\[ \gamma \equiv \Sigma c_{\iota}\varrho_{\iota} + \Sigma c'_{\iota}\lambda_{\iota} \ (\text{mod.}\,\mathfrak{a}), \]
so ist die $(r+s)$-fache Schaar $(\varrho_{1}, \varrho_{2}, \ldots \varrho_{r}, \lambda_{1}, \lambda_{2}, \ldots \lambda_{s})$ ein vollständiges Restsystem von $\mathfrak{c}$ nach $\mathfrak{a}$, oder
\[ \mathfrak{c} \equiv (\varrho_{1}, \varrho_{2}, \ldots \varrho_{r}, \lambda_{1}, \lambda_{2}, \ldots \lambda_{s}) \ (\text{mod.}\,\mathfrak{a}). \]
Ist daher
\[ \begin{aligned} z\varrho_{1} &= e_{1,1}\varrho_{1} + \cdots + e_{r,1}\varrho_{r} + \beta_{1} \\ &\cdots\cdots\cdots\cdots \\ z\varrho_{r} &= e_{1,r}\varrho_{1} + \cdots + e_{r,r}\varrho_{r} + \beta_{r}, \end{aligned} \]
worin die $e_{\iota,\varkappa}$ Constanten, die $\beta_{\iota}$ Functionen in $\mathfrak{b}$ sind, also:
\[ (\mathfrak{c}, \mathfrak{b}) = (-1)^{r} \begin{vmatrix} e_{1,1}-z, & \ldots & e_{r,1} \\ \cdots & \cdots & \cdots \\ e_{1,r}, & \ldots & e_{r,r}-z \end{vmatrix}, \]
ferner:
\[ \left. \begin{aligned} \beta_{1} &\equiv h_{1,1}\lambda_{1} + \cdots + h_{s,1}\lambda_{s} \\ &\cdots\cdots\cdots\cdots \\ \beta_{r} &\equiv h_{1,r}\lambda_{1} + \cdots + h_{s,r}\lambda_{s} \\ z\lambda_{1} &\equiv c_{1,1}\lambda_{1} + \cdots + c_{s,1}\lambda_{s} \\ &\cdots\cdots\cdots\cdots \\ z\lambda_{s} &\equiv c_{1,s}\lambda_{1} + \cdots + c_{s,s}\lambda_{s} \end{aligned} \right\} (\text{mod.}\,\mathfrak{a}) \]
mit constanten Coefficienten $h_{\iota,\varkappa}$, $c_{\iota,\varkappa}$, also
\[ (\mathfrak{b}, \mathfrak{a}) = (-1)^{s} \begin{vmatrix} c_{1,1}-z, & \ldots & c_{s,1} \\ \cdots & \cdots & \cdots \\ c_{1,s}, & \ldots & c_{s,s}-z \end{vmatrix}, \]

\origpage{203}
so folgt:
\[ \left. \begin{aligned} z\varrho_{1} &\equiv e_{1,1}\varrho_{1} + \cdots + e_{r,1}\varrho_{r} + h_{1,1}\lambda_{1} + \cdots + h_{s,1}\lambda_{s} \\ &\cdots\cdots\cdots\cdots \\ z\varrho_{r} &\equiv e_{1,r}\varrho_{1} + \cdots + e_{r,r}\varrho_{r} + h_{1,r}\lambda_{1} + \cdots + h_{s,r}\lambda_{s} \\ z\lambda_{1} &\equiv c_{1,1}\lambda_{1} + \cdots + c_{s,1}\lambda_{s} \\ &\cdots\cdots\cdots\cdots \\ z\lambda_{s} &\equiv c_{1,s}\lambda_{1} + \cdots + c_{s,s}\lambda_{s} \end{aligned} \right\} (\text{mod.}\,\mathfrak{a}) \]
und hieraus
\[ (\mathfrak{c}, \mathfrak{a}) = (-1)^{r+s} \begin{vmatrix} e_{1,1}-z, & \ldots & e_{r,1}, & h_{1,1}, & \ldots & h_{s,1} \\ \cdots & \cdots & \cdots & \cdots & \cdots & \cdots \\ e_{1,r}, & \ldots & e_{r,r}-z, & h_{1,r}, & \ldots & h_{s,r} \\ 0 & \ldots & 0 & c_{1,1}-z, & \ldots & c_{s,1} \\ \cdots & \cdots & \cdots & \cdots & \cdots & \cdots \\ 0 & \ldots & 0 & c_{1,s}, & \ldots & c_{s,s}-z \end{vmatrix} = (\mathfrak{c}, \mathfrak{b})(\mathfrak{b}, \mathfrak{a}). \]

3. Wenn die Basis-Functionen $\beta_{1}, \beta_{2}, \ldots \beta_{s}$ eines endlichen Moduls $\mathfrak{b} = [\beta_{1}, \beta_{2}, \ldots \beta_{s}]$ durch Multiplication mit von Null verschiedenen ganzen rationalen Functionen von $z$ in Functionen eines Moduls $\mathfrak{a}$ verwandelt werden können, so ist die Norm von $\mathfrak{a}$ in Bezug auf $\mathfrak{b}$, $(\mathfrak{b}, \mathfrak{a})$ von Null verschieden. Zugleich ist das kleinste gemeinschaftliche Vielfache von $\mathfrak{a}$ und $\mathfrak{b}$ ein endlicher Modul $\mathfrak{m}$, von dem eine irreductible Basis in der Form angenommen werden kann:
\[ \begin{aligned} \mu_{1} &= a_{1,1}\beta_{1}, \\ \mu_{2} &= a_{1,2}\beta_{1} + a_{2,2}\beta_{2}, \\ &\cdots\cdots\cdots\cdots \\ \mu_{s} &= a_{1,s}\beta_{1} + a_{2,s}\beta_{2} + \cdots + a_{s,s}\beta_{s}, \end{aligned} \]
worin die Coefficienten $a_{\iota,\varkappa}$ ganze rationale Functionen von $z$ sind, und zwar so, dass
\[ (\mathfrak{b}, \mathfrak{a}) = a_{1,1}a_{2,2}\ldots a_{s,s}. \]

Zum Beweise dieses wichtigen Theorems nehmen wir an, es sei $\mathfrak{a}_{1}$ der grösste gemeinschaftliche Theiler von $\mathfrak{a}$ und $[\beta_{1}]$, $\mathfrak{a}_{2}$ der von $\mathfrak{a}_{1}$ und $[\beta_{2}]$ u.\ s.\ f., so dass $\mathfrak{a}_{r}$ der Inbegriff aller Functionen von der Form
\[ \alpha_{r} = \alpha + y_{1}\beta_{1} + \cdots + y_{r}\beta_{r} \]
ist, wo $\alpha$ eine Function in $\mathfrak{a}$, $y_{1}, \ldots y_{r}$ ganze rationale Functionen von $z$ sind. Es ist dann $\mathfrak{a}_{s}$ der grösste gemeinschaftliche Theiler von $\mathfrak{a}$ und $\mathfrak{b}$. Da nun jeder Modul $\mathfrak{a}_{r}$ durch den folgenden $\mathfrak{a}_{r+1}$ theilbar ist, so folgt aus

\origpage{204}
1. und 2.
\[ (\mathfrak{b}, \mathfrak{a}) = (\mathfrak{a}_{s}, \mathfrak{a}) = (\mathfrak{a}_{s}, \mathfrak{a}_{s-1})(\mathfrak{a}_{s-1}, \mathfrak{a}_{s-2})\ldots(\mathfrak{a}_{1}, \mathfrak{a}), \]
und es handelt sich also noch um die Bestimmung von $(\mathfrak{a}_{r}, \mathfrak{a}_{r-1})$. Es ist aber
\[ \alpha_{r} = \alpha_{r-1} + y_{r}\beta_{r} \equiv y_{r}\beta_{r} \quad (\text{mod.}\,\mathfrak{a}_{r-1}), \]
und nach der Voraussetzung giebt es eine von Null verschiedene ganze rationale Function $x_{r}$ von $z$, für welche
\[ x_{r}\beta_{r} \equiv 0 \quad (\text{mod.}\,\mathfrak{a}), \]
also auch
\[ x_{r}\beta_{r} \equiv 0 \quad (\text{mod.}\,\mathfrak{a}_{r-1}). \]
Ist nun $a_{r,r}$ unter allen der letzteren Congruenz genügenden Functionen $x_{r}$ eine von möglichst niedrigem Grade $m_{r}$, die zugleich so angenommen sei, dass der Coefficient der höchsten Potenz von $z = 1$ ist, so sind alle andern dieser Congruenz genügenden Functionen $x_{r}$ durch $a_{r,r}$ theilbar; denn es ist für ein beliebiges ganzes rationales $q$
\[ (x_{r} - qa_{r,r})\beta_{r} \equiv 0 \quad (\text{mod.}\,\mathfrak{a}_{r-1}), \]
und wenn $x_{r}$ nicht durch $a_{r,r}$ theilbar ist, so lässt sich $q$ so wählen, dass $x_{r} - qa_{r,r}$ von niedrigerem Grade wird als $a_{r,r}$, gegen die Voraussetzung.

Setzt man also
\[ y_{r} = qa_{r,r} + b_{r,r} \]
und bestimmt $q$ so, dass der Grad von $b_{r,r}$ kleiner als $m_{r}$ wird, so folgt:
\[ \alpha_{r} \equiv b_{r,r}\beta_{r} \ (\text{mod.}\,\mathfrak{a}_{r-1}) \]
und hieraus
\[ \mathfrak{a}_{r} \equiv (\beta_{r}, z\beta_{r}, \ldots z^{m_{r}-1}\beta_{r}) \ (\text{mod.}\,\mathfrak{a}_{r-1}). \]
Wenn man daher für den Augenblick setzt:
\[ \begin{aligned} a_{r,r} &= c_{0} + c_{1}z + \cdots + c_{m_{r}-1}z^{m_{r}-1} + z^{m_{r}}, \\ \lambda_{k} &= z^{k-1}\beta_{r}, \end{aligned} \]
so folgt:
\[ \begin{gathered} z\lambda_{1} = \lambda_{2}, \quad z\lambda_{2} = \lambda_{3}, \quad \ldots \\ z\lambda_{m_{r}} \equiv -c_{0}\lambda_{1} - c_{1}\lambda_{2} - \cdots - c_{m_{r}-1}\lambda_{m_{r}} \ (\text{mod.}\,\mathfrak{a}_{r-1}) \end{gathered} \]
also:
\[ (\mathfrak{a}_{r}, \mathfrak{a}_{r-1}) = (-1)^{m_{r}} \begin{vmatrix} -z, & 1, & 0, & \ldots & 0 \\ 0, & -z, & 1, & \ldots & 0 \\ \cdots & \cdots & \cdots & \cdots & \cdots \\ 0, & 0, & 0, & \ldots & 1 \\ -c_{0}, & -c_{1}, & -c_{2}, & \ldots & -c_{m_{r}-1}-z \end{vmatrix} = a_{r,r}. \]

\origpage{205}
Hieraus ergiebt sich wie in 2., dass das Functionensystem
\[ \begin{aligned} &\beta_{1}, \quad z\beta_{1}, \quad \ldots \quad z^{m_{1}-1}\beta_{1}, \\ &\beta_{2}, \quad z\beta_{2}, \quad \ldots \quad z^{m_{2}-1}\beta_{2}, \\ &\cdots\cdots\cdots\cdots \\ &\beta_{s}, \quad z\beta_{s}, \quad \ldots \quad z^{m_{s}-1}\beta_{s}, \end{aligned} \]
eine Basis eines vollständigen Restsystems von $\mathfrak{b}$ nach $\mathfrak{a}$ bildet, und dass
\[ (\mathfrak{b}, \mathfrak{a}) = a_{1,1}a_{2,2}\ldots a_{s,s} \]
vom Grade
\[ m = m_{1} + m_{2} + \cdots + m_{s} \]
ist.

Da nun $a_{r,r}\beta_{r} \equiv 0 \ (\text{mod.}\,\mathfrak{a}_{r-1})$, so lässt sich eine Function $\mu_{r}$ in $\mathfrak{a}$ und ganze rationale Functionen $a_{k,r}$ so bestimmen, dass
\[ \mu_{r} = a_{1,r}\beta_{1} + a_{2,r}\beta_{2} + \cdots + a_{r\,r}\beta_{r} \]
\ednote{Printed $a_{r\,r}$ without the comma between the two indices; the notation throughout requires $a_{r,r}$. An apparent dropped sort.}
wird; die auf diese Weise bestimmten Functionen
\[ \begin{aligned} \mu_{1} &= a_{1,1}\beta_{1}, \\ \mu_{2} &= a_{1,2}\beta_{1} + a_{2,2}\beta_{2}, \\ &\cdots\cdots\cdots\cdots \\ \mu_{s} &= a_{1,s}\beta_{1} + a_{2,s}\beta_{2} + \cdots + a_{s,s}\beta_{s} \end{aligned} \]
sind, da keine der Functionen $a_{1,1}, \ldots a_{s,s}$ verschwindet, rational unabhängig und sind sämmtlich zugleich in $\mathfrak{a}$ und in $\mathfrak{b}$, also auch in dem kleinsten gemeinschaftlichen Vielfachen $\mathfrak{m}$ dieser beiden Moduln enthalten. Es ist noch nachzuweisen, dass dieselben eine Basis von $\mathfrak{m}$ bilden.

Es sei $\mathfrak{m}_{r}$ das kleinste gemeinschaftliche Vielfache von $\mathfrak{a}$ und $[\beta_{1}, \beta_{2}, \ldots \beta_{r}]$, $\mathfrak{m}_{s} = \mathfrak{m}$, so dass unter den Moduln $\mathfrak{m}_{1}, \mathfrak{m}_{2}, \ldots \mathfrak{m}_{s}$ jeder durch alle folgenden, also auch durch $\mathfrak{m}$ theilbar ist, und
\[ \nu_{r} = z_{1}\beta_{1} + z_{2}\beta_{2} + \cdots + z_{r}\beta_{r} \]
eine Function in $\mathfrak{m}_{r}$, also auch in $\mathfrak{a}$.

Es ist hiernach
\[ z_{r}\beta_{r} \equiv 0 \quad (\text{mod.}\,\mathfrak{a}_{r-1}), \]
also
\[ z_{r} = x_{r}a_{r,r}, \]
worin $x_{r}$ eine ganze rationale Function bedeutet. Daher ist
\[ \nu_{r} - x_{r}\mu_{r} \equiv 0 \ (\text{mod.}\,\mathfrak{m}_{r-1}), \quad \nu_{1} - x_{1}\mu_{1} = 0, \]
woraus folgt:
\[ \nu_{r} = x_{1}\mu_{1} + x_{2}\mu_{2} + \cdots + x_{r}\mu_{r}, \]

\origpage{206}
also:
\[ \begin{aligned} \mathfrak{m}_{r} &= [\mu_{1}, \mu_{2}, \ldots \mu_{r}], \\ \mathfrak{m} &= [\mu_{1}, \mu_{2}, \ldots \mu_{s}], \end{aligned} \]
w.\ z.\ b.\ w.

Hiernach enthält eine irreductible Basis des Moduls $\mathfrak{m}$ genau ebenso viele Functionen wie eine irreductible Basis von $\mathfrak{b}$. Wählt man statt der Basis $\mu_{1}, \mu_{2}, \ldots \mu_{s}$ eine andere $\mu'_{1}, \mu'_{2}, \ldots \mu'_{s}$, so lässt sich $\mu'_{1}, \mu'_{2}, \ldots \mu'_{s}$ in der Form ausdrücken
\[ \mu'_{k} = a'_{1,k}\beta_{1} + a'_{2,k}\beta_{2} + \cdots + a'_{s,k}\beta_{s} \]
mit ganzen rationalen Coefficienten $a'_{\iota,k}$, und aus §~4, 2. ergiebt sich
\[ (\mathfrak{b}, \mathfrak{a}) = \text{const.}\,\Sigma \pm a'_{1,1}a'_{2,2}\ldots a'_{s,s}. \]

4. Machen wir ins Besondere die Annahme, es sei $\mathfrak{a}$ gleichfalls ein endlicher Modul, der eine irreductible Basis von ebenso vielen Functionen besitzt wie $\mathfrak{b}$, und es sei ausserdem $\mathfrak{a}$ theilbar durch $\mathfrak{b}$, dann lassen sich, wenn
\[ \mathfrak{a} = [\alpha_{1}, \alpha_{2}, \ldots \alpha_{s}] \]
ist, die ganzen rationalen Functionen $b_{\iota,k}$ von $z$ so bestimmen, dass
\[ \alpha_{k} = b_{1,k}\beta_{1} + b_{2,k}\beta_{2} + \cdots + b_{s,k}\beta_{s}, \]
und die Voraussetzung von 3., dass die Functionen $\beta_{\iota}$ durch Multiplication mit ganzen rationalen Functionen von $z$ in Functionen des Moduls $\mathfrak{a}$ verwandelt werden können, ist erfüllt, wie man durch Auflösung dieses Gleichungssystems erkennt. Zugleich ist hier $\mathfrak{a}$ selbst das kleinste gemeinschaftliche Vielfache von $\mathfrak{a}$ und $\mathfrak{b}$, und daraus ergiebt sich
\[ (\mathfrak{b}, \mathfrak{a}) = \text{const.}\,\Sigma \pm b_{1,1}b_{2,2}\ldots b_{n,n}. \]
\ednote{Printed $b_{n,n}$ as the last diagonal term; the bases here have $s$ members and the coefficients run to $b_{s,k}$, so $b_{s,s}$ is expected. An apparent misprint.}

5. Ist $\mathfrak{m}$ das kleinste gemeinschaftliche Vielfache zweier Moduln $\mathfrak{a}$, $\mathfrak{b}$ und $\nu$ eine beliebige Function in $\Omega$, so ist, wie sich aus der Definition ohne Schwierigkeit ergiebt, $\nu\mathfrak{m}$ das kleinste gemeinschaftliche Vielfache von $\nu\mathfrak{a}$ und $\nu\mathfrak{b}$. Ist $(\mathfrak{b}, \mathfrak{a}) = 0$, so ist auch $(\nu\mathfrak{b}, \nu\mathfrak{a}) = 0$. Ist aber $(\mathfrak{b}, \mathfrak{a})$ und $\nu$ von Null verschieden, so ergiebt sich
\[ (\nu\mathfrak{b}, \nu\mathfrak{a}) = (\mathfrak{b}, \mathfrak{a}), \]
wenn man in 3. die Basis-Functionen $\mu_{\iota}$, $\beta_{\iota}$ von $\mathfrak{m}$ und $\mathfrak{b}$ durch $\nu\mu_{\iota}$, $\nu\beta_{\iota}$ ersetzt.

\subsection*{§~7. Die Ideale in $\mathfrak{o}$.}

Ein System $\mathfrak{a}$ von \emph{ganzen} Functionen von $z$ im Körper $\Omega$ heisst ein \emph{Ideal}, wenn es die beiden folgenden Bedingungen erfüllt:

\origpage{207}

I. Summe und Differenz je zweier Functionen in $\mathfrak{a}$ ergeben wieder eine Function in $\mathfrak{a}$.

II. Das Product einer jeden Function in $\mathfrak{a}$ mit einer jeden Function in $\mathfrak{o}$ (§~3) ist wieder eine Function in $\mathfrak{a}$.

Jedes Ideal ist also zugleich ein Modul und alle für die Moduln erklärten Begriffe und Bezeichnungen können auf die Ideale angewandt werden.

Der Modul $\mathfrak{o}$ (das System aller ganzen Functionen von $z$) ist selbst ein Ideal, und \emph{jedes Ideal ist durch $\mathfrak{o}$ theilbar}. Ebenso ist, wenn $\mu$ eine beliebige von Null verschiedene Function in $\mathfrak{o}$ bedeutet, der Modul $\mathfrak{o}\mu$ (das System aller durch $\mu$ theilbaren ganzen Functionen) ein Ideal. Ein solches Ideal soll ein \emph{Hauptideal} genannt werden. Ist $\omega_{1}, \omega_{2}, \ldots \omega_{n}$ eine Basis von $\mathfrak{o}$, so ist
\[ \mathfrak{o}\mu = [\omega_{1}\mu, \omega_{2}\mu, \ldots \omega_{n}\mu] \]
und $\mathfrak{o}\mu$ ist das kleinste gemeinschaftliche Vielfache von $\mathfrak{o}$ und $\mathfrak{o}\mu$. Daher ist nach §~6, 4. und nach der Definition (4.) in §~2:
\[ (\mathfrak{o}, \mathfrak{o}\mu) = \text{const.}\,N(\mu) \tag{1.} \]
und mithin von Null verschieden.

Ist $\mathfrak{a}$ irgend ein Ideal und $\alpha$ eine beliebige Function in $\mathfrak{a}$, so ist (wegen II.) das Hauptideal $\mathfrak{o}\alpha$ theilbar durch $\mathfrak{a}$, und mithin nach §~6, 2.:
\[ (\mathfrak{o}, \mathfrak{o}\alpha) = (\mathfrak{o}, \mathfrak{a})(\mathfrak{a}, \mathfrak{o}\alpha), \tag{2.} \]
mithin auch $(\mathfrak{o}, \mathfrak{a})$ von Null verschieden. Da nun wieder $\mathfrak{a}$ das kleinste gemeinschaftliche Vielfache von $\mathfrak{a}$ und $\mathfrak{o}$ ist, so besitzt $\mathfrak{a}$ nach §~6, 3. eine irreductible Basis, welche aus $n$ ganzen Functionen $\alpha_{1}, \alpha_{2}, \ldots \alpha_{n}$ besteht, die demnach auch eine Basis des Körpers $\Omega$ bilden.

Die Norm von $\mathfrak{a}$ in Bezug auf $\mathfrak{o}$, d.\ h.\ die ganze rationale Function $(\mathfrak{o}, \mathfrak{a})$ von $z$ soll die \emph{Norm des Ideals} $\mathfrak{a}$ genannt und mit $N(\mathfrak{a})$ bezeichnet werden. Der Grad dieser ganzen rationalen Function heisst zugleich der \emph{Grad des Ideals} $\mathfrak{a}$.

Ist
\[ \mathfrak{a} = [\alpha_{1}, \alpha_{2}, \ldots \alpha_{n}], \quad \mathfrak{o} = [\omega_{1}, \omega_{2}, \ldots \omega_{n}] \]
und
\[ \begin{aligned} \alpha_{1} &= a_{1,1}\omega_{1} + a_{2,1}\omega_{2} + \cdots + a_{n,1}\omega_{n}, \\ \alpha_{2} &= a_{1,2}\omega_{1} + a_{2,2}\omega_{2} + \cdots + a_{n,2}\omega_{n}, \\ &\cdots\cdots\cdots\cdots \\ \alpha_{n} &= a_{1,n}\omega_{1} + a_{2,n}\omega_{2} + \cdots + a_{n,n}\omega_{n} \end{aligned} \]

\origpage{208}
mit ganzen rationalen Coefficienten $a_{\iota,\varkappa}$, so ergiebt sich aus §~6, 4.:
\[ N(\mathfrak{a}) = \text{const.}\,\Sigma \pm a_{1,1}a_{2,2}\ldots a_{n,n}. \tag{3.} \]
Da jede Function in $\mathfrak{o}$, also auch die Function „1“ durch Multiplication mit $N(\mathfrak{a})$ in eine Function des Ideals $\mathfrak{a}$ verwandelt wird, so ist $N(\mathfrak{a})$ stets eine Function in $\mathfrak{a}$.

Die Norm des Ideals $\mathfrak{o}$ ist gleich 1 und umgekehrt ist $\mathfrak{o}$ das einzige Ideal, welches diese Eigenschaft hat. Auch ist $\mathfrak{o}$ das einzige Ideal, welches die Function „1“ (oder eine Constante) enthält.

Ist $\alpha$ eine Function in $\mathfrak{a}$, so folgt aus (1.), (2.), (3.):
\[ N(\alpha) = \text{const.}\,N(\mathfrak{a})(\mathfrak{a}, \mathfrak{o}\alpha), \tag{4.} \]
d.\ h.\ die Norm einer jeden in $\mathfrak{a}$ enthaltenen Function ist durch die Norm von $\mathfrak{a}$ theilbar.

Für die Congruenzen in Bezug auf einen Idealmodul gilt der folgende Satz, welcher die Ideale wesentlich von den allgemeinen Moduln unterscheidet.

Sind $\mu$, $\mu_{1}$, $\nu$, $\nu_{1}$ Functionen in $\mathfrak{o}$, welche den Congruenzen genügen
\[ \mu \equiv \mu_{1}, \quad \nu \equiv \nu_{1} \ (\text{mod.}\,\mathfrak{a}), \]
so ist auch
\[ \mu\nu \equiv \mu_{1}\nu_{1} \ (\text{mod.}\,\mathfrak{a}). \]

\subsection*{§~8. Multiplication und Theilung der Ideale.}

Aus den Grundeigenschaften I., II. der Ideale und aus den Begriffsbestimmungen in §~4 ergiebt sich zunächst:

1. Das kleinste gemeinschaftliche Vielfache, der grösste gemeinschaftliche Theiler, das Product von zwei (und also auch von beliebig vielen) Idealen sind selbst Ideale. Ebenso ist, wenn $\nu$ eine Function in $\mathfrak{o}$, $\mathfrak{a}$ ein Ideal ist, das Product $\mathfrak{a}\nu$ ein Ideal.

2. Das Product aus mehreren Idealen ist durch jeden seiner Factoren theilbar, und es ist für jedes Ideal $\mathfrak{a}$
\[ \mathfrak{a}\mathfrak{o} = \mathfrak{a}; \]
denn nach I., II. ist jede Function in $\mathfrak{a}\mathfrak{o}$ zugleich eine Function in $\mathfrak{a}$, und, da $\mathfrak{o}$ die Function „1“ enthält, auch umgekehrt jede Function in $\mathfrak{a}$ zugleich eine Function in $\mathfrak{a}\mathfrak{o}$.

\origpage{209}

3. Ein Hauptideal $\mathfrak{o}\mu$ ist dann und nur dann theilbar durch ein Hauptideal $\mathfrak{o}\nu$, wenn die ganze Function $\mu$ theilbar ist durch die ganze Function $\nu$.

Wir fügen noch folgende Definitionen hinzu:

4. \emph{Definition.} Eine Function $\alpha$ in $\mathfrak{o}$ soll durch das Ideal $\mathfrak{a}$ \emph{theilbar} heissen, wenn das Hauptideal $\mathfrak{o}\alpha$ durch $\mathfrak{a}$ theilbar, oder, was dasselbe sagt, wenn $\alpha$ eine Function in $\mathfrak{a}$ ist.

5. \emph{Definition.} Zwei Ideale $\mathfrak{a}$, $\mathfrak{b}$ heissen \emph{relativ prim}, wenn ihr grösster gemeinschaftlicher Theiler $\mathfrak{o}$ ist. Die nothwendige und hinreichende Bedingung dafür ist, dass in $\mathfrak{a}$ eine Function $\alpha$, in $\mathfrak{b}$ eine Function $\beta$ existirt der Art, dass
\[ \alpha + \beta = 1, \]
oder, anders ausgedrückt, dass in $\mathfrak{a}$ eine der Congruenz $\alpha \equiv 1 \ (\text{mod.}\,\mathfrak{b})$ oder in $\mathfrak{b}$ eine der Congruenz $\beta \equiv 1 \ (\text{mod.}\,\mathfrak{a})$ genügende Function existirt.

6. \emph{Definition.} Ein von $\mathfrak{o}$ verschiedenes Ideal $\mathfrak{p}$ heisst ein \emph{Primideal}, wenn kein anderes Ideal ausser $\mathfrak{p}$ und $\mathfrak{o}$ in $\mathfrak{p}$ aufgeht.

Auf Grund dieser Definitionen ergeben sich nun die folgenden Sätze über die Theilbarkeit der Ideale.

7. Sind $\mathfrak{a}$, $\mathfrak{b}$ zwei Ideale mit dem kleinsten gemeinschaftlichen Vielfachen $\mathfrak{m}$ und dem grössten gemeinschaftlichen Theiler $\mathfrak{d}$, so folgt aus §~6, 1., 2.
\[ \begin{aligned} N(\mathfrak{m}) &= N(\mathfrak{b})(\mathfrak{b}, \mathfrak{m}) = N(\mathfrak{b})(\mathfrak{b}, \mathfrak{a}), \\ N(\mathfrak{a}) &= N(\mathfrak{d})(\mathfrak{d}, \mathfrak{a}) = N(\mathfrak{d})(\mathfrak{b}, \mathfrak{a}), \end{aligned} \]
folglich $(\mathfrak{b}, \mathfrak{a})$ von Null verschieden und
\[ N(\mathfrak{a})N(\mathfrak{b}) = N(\mathfrak{m})N(\mathfrak{d}). \]

8. Ist das Ideal $\mathfrak{a}$ theilbar durch das Ideal $\mathfrak{b}$, so ist, nach §~6, 2.
\[ N(\mathfrak{a}) = (\mathfrak{b}, \mathfrak{a})N(\mathfrak{b}), \]
also $N(\mathfrak{a})$ theilbar durch $N(\mathfrak{b})$.

Ist ins Besondere $(\mathfrak{b}, \mathfrak{a}) = 1$, so ist auch $\mathfrak{b}$ theilbar durch $\mathfrak{a}$, und es folgt:

9. Ist $\mathfrak{a}$ theilbar durch $\mathfrak{b}$ und ist zugleich $N(\mathfrak{a}) = N(\mathfrak{b})$, so ist $\mathfrak{a} = \mathfrak{b}$, d.\ h.\ beide Ideale sind identisch.

10. Ist $\mathfrak{a}$ theilbar durch $\mathfrak{a}_{1}$, $\mathfrak{b}$ durch $\mathfrak{b}_{1}$, so ist $\mathfrak{a}\mathfrak{b}$ theilbar durch $\mathfrak{a}_{1}\mathfrak{b}_{1}$ (§~4, 7.).

\origpage{210}

11. Ist ein Ideal $\mathfrak{a}$ theilbar durch ein Hauptideal $\mathfrak{o}\mu$, so sind alle Functionen in $\mathfrak{a}$ von der Form $\beta\mu$, und der Inbegriff der Functionen $\beta$ ist wieder ein Ideal $\mathfrak{b}$, so dass man setzen kann
\[ \mathfrak{a} = \mu\mathfrak{b}. \]

12. Ist $\mu$ eine beliebige von Null verschiedene Function in $\mathfrak{o}$ und das Ideal $\mathfrak{a}\mu$ theilbar durch das Ideal $\mathfrak{b}\mu$, so ist $\mathfrak{a}$ theilbar durch $\mathfrak{b}$, und aus $\mathfrak{a}\mu = \mathfrak{b}\mu$ folgt $\mathfrak{a} = \mathfrak{b}$.

13. Das kleinste gemeinschaftliche Vielfache zweier Ideale $\mathfrak{a}$, $\mathfrak{o}\nu$, davon eines ein Hauptideal ist, hat nach 11. die Form $\mathfrak{r}\nu$, worin $\mathfrak{r}$ ein Ideal ist. Da andererseits $\mathfrak{a}\nu$ ein gemeinschaftliches Vielfache von $\mathfrak{a}$ und $\mathfrak{o}\nu$, also durch $\mathfrak{r}\nu$ theilbar ist, so ist nach 12. $\mathfrak{r}$ ein Theiler von $\mathfrak{a}$.

14. Ist $\mathfrak{a}$ ein Ideal, $\nu$ eine Function in $\mathfrak{o}$, so ist nach §~6, 2., 5.:
\[ (\mathfrak{o}, \mathfrak{a}\nu) = (\mathfrak{o}, \mathfrak{o}\nu)(\mathfrak{o}\nu, \mathfrak{a}\nu) = (\mathfrak{o}, \mathfrak{o}\nu)(\mathfrak{o}, \mathfrak{a}), \]
also
\[ N(\mathfrak{a}\nu) = \text{const.}\,N(\mathfrak{a})N(\nu). \]
Ist also $\mathfrak{r}\nu$ das kleinste gemeinschaftliche Vielfache, $\mathfrak{d}$ der grösste gemeinschaftliche Theiler der beiden Ideale $\mathfrak{a}$, $\mathfrak{o}\nu$, so ergiebt sich aus 7.
\[ N(\mathfrak{a}) = N(\mathfrak{r})N(\mathfrak{d}). \]

15. Jedes von $\mathfrak{o}$ verschiedene Ideal $\mathfrak{a}$ ist durch ein Primideal $\mathfrak{p}$ theilbar.

Ist nämlich $\mathfrak{a}$ kein Primideal, so hat es mindestens einen von $\mathfrak{o}$ verschiedenen echten Theiler, und von diesen sei $\mathfrak{p}$ ein solcher, dessen Norm von möglichst niedrigem Grade ist. Dieser kann keinen von $\mathfrak{o}$ verschiedenen echten Theiler $\mathfrak{p}'$ haben, denn es wäre auch $\mathfrak{p}'$ ein Theiler von $\mathfrak{a}$ und zugleich (nach 8.) $N(\mathfrak{p}')$ von niedrigerem Grade als $N(\mathfrak{p})$. Dies widerspricht der Voraussetzung über $\mathfrak{p}$, und folglich ist $\mathfrak{p}$ ein Primideal.

16. Ist $\mathfrak{a}$ relativ prim zu $\mathfrak{b}$, so ist $\mathfrak{a}\mathfrak{b}$ das kleinste gemeinschaftliche Vielfache von $\mathfrak{a}$ und $\mathfrak{b}$, und folglich ist jedes durch $\mathfrak{a}$ und durch $\mathfrak{b}$ theilbare Ideal auch durch das Product $\mathfrak{a}\mathfrak{b}$ theilbar.

Denn nach Voraussetzung giebt es in $\mathfrak{a}$, $\mathfrak{b}$ zwei Functionen $\alpha_{1}$, $\beta_{1}$ der Art, dass
\[ \alpha_{1} + \beta_{1} = 1 \]
ist (5.). Ist andererseits $\alpha = \beta$ eine Function des kleinsten gemeinschaftlichen Vielfachen $\mathfrak{m}$ von $\mathfrak{a}$ und $\mathfrak{b}$, so ist hiernach
\[ \alpha = \beta = \alpha_{1}\beta + \alpha\beta_{1}, \]

\origpage{211}
also eine Function in $\mathfrak{a}\mathfrak{b}$. Es ist demnach $\mathfrak{m}$ theilbar durch $\mathfrak{a}\mathfrak{b}$, und da umgekehrt (zufolge 2.) $\mathfrak{a}\mathfrak{b}$ durch $\mathfrak{m}$ theilbar ist, so ist $\mathfrak{m}$ mit $\mathfrak{a}\mathfrak{b}$ identisch, und aus 7. folgt noch für diesen Fall
\[ N(\mathfrak{a}\mathfrak{b}) = N(\mathfrak{a})N(\mathfrak{b}). \]

17. Ist $\mathfrak{a}$ ein beliebiges Ideal, $\mathfrak{p}$ ein Primideal, so ist entweder $\mathfrak{a}$ durch $\mathfrak{p}$ theilbar oder $\mathfrak{a}$ relativ prim zu $\mathfrak{p}$; denn da $\mathfrak{p}$ keinen anderen Theiler hat als $\mathfrak{o}$ und $\mathfrak{p}$, so kann auch der grösste gemeinschaftliche Theiler von $\mathfrak{a}$ und $\mathfrak{p}$ kein anderer sein als $\mathfrak{o}$ oder $\mathfrak{p}$.

18. Ist $\mathfrak{a}$ relativ prim zu $\mathfrak{b}$ und zu $\mathfrak{c}$, so ist $\mathfrak{a}$ auch relativ prim zu $\mathfrak{b}\mathfrak{c}$.

Nach Voraussetzung (5.) giebt es in $\mathfrak{b}$, $\mathfrak{c}$ zwei den Congruenzen
\[ \beta \equiv 1, \quad \gamma \equiv 1 \ (\text{mod.}\,\mathfrak{a}) \]
genügende Functionen, folglich nach §~7
\[ \beta\gamma \equiv 1 \quad (\text{mod.}\,\mathfrak{a}). \]
Da $\beta\gamma$ in $\mathfrak{b}\mathfrak{c}$ enthalten ist, so ist hiermit die Behauptung erwiesen.

Es folgt hieraus noch, dass, falls das Product $\mathfrak{a}\mathfrak{b}$ durch ein Primideal theilbar ist, wenigstens einer der beiden Factoren $\mathfrak{a}$, $\mathfrak{b}$ durch $\mathfrak{p}$ theilbar sein muss, und, auf Hauptideale angewandt, dass, wenn das Product zweier ganzen Functionen, $\mu\nu$, in $\mathfrak{p}$ enthalten ist, wenigstens der eine der beiden Factoren $\mu$, $\nu$ in $\mathfrak{p}$ enthalten sein muss.

19. Ist $\mathfrak{a}$ relativ prim zu $\mathfrak{c}$ und $\mathfrak{a}\mathfrak{b}$ durch $\mathfrak{c}$ theilbar, so ist $\mathfrak{b}$ durch $\mathfrak{c}$ theilbar. Nach Voraussetzung giebt es in $\mathfrak{a}$ eine Function $\alpha$, welche der Congruenz genügt
\[ \alpha \equiv 1 \quad (\text{mod.}\,\mathfrak{c}). \]
Ist nun $\beta$ eine beliebige Function in $\mathfrak{b}$, so ist hiernach
\[ \beta \equiv \alpha\beta \quad \text{und nach Vor.} \quad \equiv 0 \ (\text{mod.}\,\mathfrak{c}), \]
folglich $\beta$ in $\mathfrak{c}$ enthalten, also $\mathfrak{b}$ durch $\mathfrak{c}$ theilbar.

\subsection*{§~9. Gesetze der Theilbarkeit der Ideale.}

Alle diese Sätze, die sich meist unmittelbar aus der Definition der Ideale ergaben, reichen nicht aus, um die vollständige Analogie zu beweisen, die zwischen den Gesetzen der Theilbarkeit der Ideale und denen der ganzen rationalen Functionen herrscht. Wir stützen uns bei diesem Beweis auf folgenden Satz:

\origpage{212}

1. \emph{Ist $\mathfrak{a}$ ein Ideal und $k$ eine beliebige ganze rationale Function von $z$, so lässt sich in $\mathfrak{a}$ eine Function $\alpha$ so auswählen, dass $(\mathfrak{a}, \mathfrak{o}\alpha)$ mit $k$ keinen Theiler gemeinschaftlich hat}${}^{*)}$.

Ist nämlich
\[ \begin{aligned} \mathfrak{a} &= [\alpha_{1}, \alpha_{2}, \ldots \alpha_{n}], \\ \mathfrak{o} &= [\omega_{1}, \omega_{2}, \ldots \omega_{n}], \end{aligned} \]
und $\alpha$ eine beliebige Function in $\mathfrak{a}$, so lassen sich die ganzen rationalen Functionen $x_{h,k}$ so bestimmen, dass
\[ \begin{aligned} \alpha\omega_{1} &= x_{1,1}\alpha_{1} + x_{2,1}\alpha_{2} + \cdots + x_{n,1}\alpha_{n}, \\ \alpha\omega_{2} &= x_{1,2}\alpha_{1} + x_{2,2}\alpha_{2} + \cdots + x_{n,2}\alpha_{n}, \\ &\cdots\cdots\cdots\cdots \\ \alpha\omega_{n} &= x_{1,n}\alpha_{1} + x_{2,n}\alpha_{2} + \cdots + x_{n,n}\alpha_{n}, \end{aligned} \]
und es ist (§~6, 4.)
\[ (\mathfrak{a}, \mathfrak{o}\alpha) = \text{const.}\,\Sigma \pm x_{1,1}x_{2,2}\ldots x_{n,n}. \]
Ist nun $\Sigma \pm x_{1,1}x_{2,2}\ldots x_{n,n}$ durch einen Linearfactor $z-c$ von $k$ theilbar, so lässt sich eine nicht durch $z-c$ theilbare Function $\omega$ in $\mathfrak{o}$ und eine Function $\alpha'$ in $\mathfrak{a}$ so bestimmen, dass
\[ \alpha\omega = (z-c)\alpha'{}^{**)}. \]
Setzt man nun, indem man unter $t$ eine unbestimmte Constante versteht:
\[ t(z-c) - \omega = \omega', \]
so ist
\[ N(\omega') = t^{n}(z-c)^{n} + a_{1}t^{n-1}(z-c)^{n-1} + \cdots + a_{n-1}t(z-c) + a_{n}, \]
worin die von $t$ unabhängigen Coefficienten $a_{1}, a_{2}, \ldots a_{n}$ ganze rationale Functionen von $z$ sind. Es kann nun nicht zugleich $a_{1}$ durch $z-c$, $a_{2}$ durch $(z-c)^{2}$, $\ldots$, $a_{n}$ durch $(z-c)^{n}$ theilbar sein, weil sonst gegen die Voraussetzung $\dfrac{\omega}{z-c}$ eine ganze Function wäre (§~2, 5., §~3, 4.). Daher lassen sich nicht

\textbf{*)} Die Möglichkeit, diesen Satz schon an dieser Stelle zu beweisen, unterscheidet wesentlich die Theorie der algebraischen Functionen von der der algebraischen Zahlen und gestattet bei ersterer eine nicht unerhebliche Vereinfachung im Vergleich mit letzterer.

\textbf{**)} Wenn nämlich die Determinante $\Sigma \pm x_{1,1}x_{2,2}\ldots x_{n,n}$ durch $z-c$ theilbar ist, d.\ h.\ für $z = c$ verschwindet, so kann man ein System von Constanten $c_{1}, c_{2}, \ldots c_{n}$, die nicht sämmtlich verschwinden, so bestimmen, dass die ganzen rationalen Functionen
\[ c_{1}x_{k,1} + c_{2}x_{k,2} + \cdots + c_{k}x_{k,n} \qquad {\scriptstyle (k = 1,\, 2\, \ldots\, n)} \]
\ednote{Printed $c_{k}x_{k,n}$ as the last term; the sum runs over $c_{1}, c_{2}, \ldots c_{n}$, so $c_{n}x_{k,n}$ is expected. An apparent misprint.}
für $z = c$ verschwinden, also durch $z-c$ theilbar sind, und es ist dann
\[ \omega = c_{1}\omega_{1} + c_{2}\omega_{2} + \cdots + c_{n}\omega_{n} \]
zu setzen.

\origpage{213}
alle Glieder von $N(\omega')$ durch $(z-c)^{n}$ theilen, und wenn also $(z-c)^{n-r}$ die höchste Potenz von $z-c$ ist, durch welche dieselben theilbar sind, so ist $r \geqslant 0$ und
\[ \frac{N(\omega')}{(z-c)^{n-r}} = t^{n}(z-c)^{r} + b_{1}t^{n-1} + \cdots + b_{n-1}t + b_{n} = f(t), \]
worin die ganzen rationalen Functionen $b_{1}, b_{2}, \ldots b_{n}$ nicht alle für $z = c$ verschwinden. Es giebt daher nur eine endliche Anzahl von constanten Werthen $t$, für welche $f(t)$ durch $z-c$ theilbar ist. Ist $z-c'$ ein von $z-c$ verschiedener Linearfactor von $k$, so wird $f(t)$ auch nur für eine endliche Anzahl von Werthen $t$ durch $z-c'$ theilbar. Daraus folgt, dass man über $t$ so verfügen kann, dass $N(\omega')$ nicht durch $(z-c)^{n}$ und zugleich durch keinen andern Linearfactor von $k$ theilbar wird${}^{*)}$. Setzt man, wenn dies geschehen,
\[ t\alpha - \alpha' = \alpha'', \]
welches ebenfalls eine Function in $\mathfrak{a}$ ist, so folgt
\[ \begin{aligned} \alpha\omega' &= (z-c)\alpha'', \\ N(\alpha'') &= \frac{N(\alpha)N(\omega')}{(z-c)^{n}} \end{aligned} \]
und mithin, da nach §~7, (4.)
\[ (\mathfrak{a}, \mathfrak{o}\alpha) = \text{const.}\,\frac{N(\alpha)}{N(\mathfrak{a})} \]
ist:
\[ (\mathfrak{a}, \mathfrak{o}\alpha'') = \text{const.}\,\frac{(\mathfrak{a}, \mathfrak{o}\alpha)N(\omega')}{(z-c)^{n}}. \]
Die Function $(\mathfrak{a}, \mathfrak{o}\alpha'')$ enthält daher den Factor $z-c$ mindestens einmal weniger als $(\mathfrak{a}, \mathfrak{o}\alpha)$, während sie zugleich keinen anderen Linearfactor von $k$ öfter enthält als $(\mathfrak{a}, \mathfrak{o}\alpha)$. Durch wiederholte Anwendung dieses Verfahrens ergiebt sich die Richtigkeit des ausgesprochenen Satzes.

2. Jedes Ideal $\mathfrak{a}$ kann als grösster gemeinschaftlicher Theiler zweier Hauptideale $\mathfrak{o}\mu$, $\mathfrak{o}\nu$ dargestellt werden, von denen das eine ganz beliebig, nur theilbar durch $\mathfrak{a}$, angenommen werden kann.

\emph{Beweis.} Man wähle nach Belieben in $\mathfrak{a}$ eine von Null verschiedene Function $\nu$, und eine zweite Function $\mu$ der Art, dass die beiden Functionen $(\mathfrak{a}, \mathfrak{o}\nu)$ und $(\mathfrak{a}, \mathfrak{o}\mu)$ keinen gemeinschaftlichen Theiler haben (nach 1). Ist nun $\alpha$ eine beliebige Function in $\mathfrak{a}$, so ist nach §~6 $(\mathfrak{a}, \mathfrak{o}\mu)\alpha$ in $\mathfrak{o}\mu$, $(\mathfrak{a}, \mathfrak{o}\nu)\alpha$ in $\mathfrak{o}\nu$ enthalten, so dass es zwei Functionen $\omega$, $\omega'$ in $\mathfrak{o}$ giebt, für welche
\[ (\mathfrak{a}, \mathfrak{o}\mu)\alpha = \mu\omega, \quad (\mathfrak{a}, \mathfrak{o}\nu)\alpha = \nu\omega'. \]

\textbf{*)} Diese Schlüsse gelten in der analogen Frage der Zahlentheorie nicht mehr.

\origpage{214}
Wählt man daher, was nach der Voraussetzung über $(\mathfrak{a}, \mathfrak{o}\mu)$, $(\mathfrak{a}, \mathfrak{o}\nu)$ möglich ist, zwei ganze rationale Functionen $g$, $h$ von $z$, welche der Bedingung genügen
\[ g(\mathfrak{a}, \mathfrak{o}\mu) + h(\mathfrak{a}, \mathfrak{o}\nu) = 1, \]
so folgt
\[ \alpha = g\mu\omega + h\nu\omega', \]
d.\ h.\ $\mathfrak{a}$ ist theilbar durch den grössten gemeinschaftlichen Theiler von $\mathfrak{o}\mu$ und $\mathfrak{o}\nu$. Und da letzterer umgekehrt durch $\mathfrak{a}$ theilbar ist (weil $\mathfrak{o}\mu$ und $\mathfrak{o}\nu$ durch $\mathfrak{a}$ theilbar sind), so ist er gleich $\mathfrak{a}$, w.\ z.\ b.\ w.

3. Jedes Ideal $\mathfrak{a}$ kann durch Multiplication mit einem Ideal $\mathfrak{m}$ in ein Hauptideal $\mathfrak{o}\mu = \mathfrak{a}\mathfrak{m}$ verwandelt werden.

\emph{Beweis.} Man wähle (nach 1.) in $\mathfrak{a}$ eine Function $\mu$ so, dass $(\mathfrak{a}, \mathfrak{o}\mu)$ keinen Theiler mit $N(\mathfrak{a})$ gemein hat, hierauf eine zweite Function $\nu$ so, dass $(\mathfrak{a}, \mathfrak{o}\nu)$ mit $(\mathfrak{a}, \mathfrak{o}\mu)$ keinen Theiler gemein hat. Dann ist (nach 2.) $\mathfrak{a}$ der grösste gemeinschaftliche Theiler von $\mathfrak{o}\mu$ und $\mathfrak{o}\nu$. Das kleinste gemeinschaftliche Vielfache von $\mathfrak{o}\mu$ und $\mathfrak{o}\nu$ ist (nach §~8, 13) von der Form $\mathfrak{m}\nu$, worin $\mathfrak{m}$ ein Theiler von $\mathfrak{o}\mu$ ist. Nach §~8, 14 ist alsdann
\[ N(\mathfrak{m}) = \frac{N(\mathfrak{o}\mu)}{N(\mathfrak{a})} = (\mathfrak{a}, \mathfrak{o}\mu), \]
also, nach Voraussetzung, ohne gemeinschaftlichen Theiler mit $N(\mathfrak{a})$. Bestimmt man also wieder zwei ganze rationale Functionen $g$, $h$ von $z$, so dass
\[ gN(\mathfrak{m}) + hN(\mathfrak{a}) = 1, \]
so folgt aus §~8, 5, da $N(\mathfrak{m})$ in $\mathfrak{m}$, $N(\mathfrak{a})$ in $\mathfrak{a}$ enthalten ist, dass $\mathfrak{m}$ und $\mathfrak{a}$ relative Primideale sind, und daraus, nach §~8, 16.
\[ N(\mathfrak{m}\mathfrak{a}) = N(\mathfrak{m})N(\mathfrak{a}) = N(\mathfrak{o}\mu). \]
Da nun $\mathfrak{o}\mu$ durch $\mathfrak{m}$ und durch $\mathfrak{a}$, also auch durch $\mathfrak{m}\mathfrak{a}$ theilbar ist (§~8, 16.), so ist nach §~8, 9.
\[ \mathfrak{m}\mathfrak{a} = \mathfrak{o}\mu, \]
w.\ z.\ b.\ w.${}^{*)}$

4. Ist ein Ideal $\mathfrak{c}$ theilbar durch ein Ideal $\mathfrak{a}$, so giebt es ein und \emph{nur} ein Ideal $\mathfrak{b}$, welches der Bedingung
\[ \mathfrak{a}\mathfrak{b} = \mathfrak{c} \]
genügt, welches der \emph{Quotient von $\mathfrak{c}$ durch $\mathfrak{a}$} heisst.

\textbf{*)} Man kann das Ideal $\mathfrak{m}$ zugleich so wählen, dass es relativ prim zu einem beliebigen Ideal $\mathfrak{b}$ wird. Dies wird erreicht, wenn man die Function $\mu$ so annimmt, dass $(\mathfrak{a}, \mathfrak{o}\mu) = N(\mathfrak{m})$ keinen Theiler mit $N(\mathfrak{a})N(\mathfrak{b})$ gemein hat (§~8, 8.).

\origpage{215}

Ist $\mathfrak{a}\mathfrak{b}$ theilbar durch $\mathfrak{a}\mathfrak{b}'$, so ist $\mathfrak{b}$ theilbar durch $\mathfrak{b}'$, und aus $\mathfrak{a}\mathfrak{b} = \mathfrak{a}\mathfrak{b}'$ folgt $\mathfrak{b} = \mathfrak{b}'$.

\emph{Beweis.} Es sei $\mathfrak{c}$ theilbar durch $\mathfrak{a}$ und (nach 3.) $\mathfrak{a}\mathfrak{m} = \mathfrak{o}\mu$. Es ist alsdann auch $\mathfrak{c}\mathfrak{m}$ theilbar durch $\mathfrak{a}\mathfrak{m} = \mathfrak{o}\mu$ und folglich $\mathfrak{c}\mathfrak{m} = \mathfrak{b}\mu$ (§~8, 10., 11.); also, durch Multiplication der letzten Gleichung mit $\mathfrak{a}$,
\[ \mathfrak{c}\mu = \mathfrak{a}\mathfrak{b}\mu \]
und nach §~8, 12.
\[ \mathfrak{c} = \mathfrak{a}\mathfrak{b}, \]
womit der erste Theil des Satzes bewiesen ist${}^{*)}$.

Ist ferner $\mathfrak{a}\mathfrak{b}$ theilbar durch $\mathfrak{a}\mathfrak{b}'$, so ist (§~8, 10.) $\mu\mathfrak{b}$ theilbar durch $\mu\mathfrak{b}'$, also $\mathfrak{b}$ durch $\mathfrak{b}'$. --- Ist $\mathfrak{a}\mathfrak{b} = \mathfrak{a}\mathfrak{b}'$, so folgt $\mu\mathfrak{b} = \mu\mathfrak{b}'$ und mithin $\mathfrak{b} = \mathfrak{b}'$ (§~8, 12.).

5. Jedes von $\mathfrak{o}$ verschiedene Ideal ist entweder ein Primideal, oder es lässt sich, und nur auf eine Weise, als Product von lauter Primidealen darstellen.

\emph{Beweis.} Ist das Ideal $\mathfrak{a}$ von $\mathfrak{o}$ verschieden, so ist es (§~8, 15.) durch ein Primideal $\mathfrak{p}_{1}$ theilbar, und folglich (nach 4.) $= \mathfrak{p}_{1}\mathfrak{a}_{1}$, worin $\mathfrak{a}_{1}$ ein \emph{echter} Theiler von $\mathfrak{a}$ ist (denn aus $\mathfrak{a}_{1} = \mathfrak{a}$ würde nach 4. folgen $\mathfrak{p}_{1} = \mathfrak{o}$). Es ist also der Grad von $N(\mathfrak{a}_{1})$ niedriger als der von $N(\mathfrak{a})$. Ist $\mathfrak{a}_{1}$ von $\mathfrak{o}$ verschieden, so schliesst man ebenso, dass $\mathfrak{a}_{1} = \mathfrak{p}_{2}\mathfrak{a}_{2}$ sein muss, wobei der Grad von $N(\mathfrak{a}_{2})$ wieder niedriger ist als der von $N(\mathfrak{a}_{1})$. Fährt man auf diese Weise fort, so gelangt man schliesslich nach einer endlichen Anzahl von Zerlegungen zu einem Ideal $\mathfrak{a}_{r-1} = \mathfrak{p}_{r}\mathfrak{a}_{r}$ der Art, dass $N(\mathfrak{a}_{r}) = 1$, also $\mathfrak{a}_{r} = \mathfrak{o}$ ist. Es ist daher
\[ \mathfrak{a} = \mathfrak{p}_{1}\mathfrak{p}_{2}\ldots\mathfrak{p}_{r}. \]
Wäre eine solche Zerlegung auf eine zweite Art möglich, etwa
\[ \mathfrak{p}_{1}\mathfrak{p}_{2}\ldots\mathfrak{p}_{r} = \mathfrak{q}_{1}\mathfrak{q}_{2}\ldots\mathfrak{q}_{s}, \]
so müsste (§~8, 18.) von den Primidealen $\mathfrak{p}_{1}, \mathfrak{p}_{2}, \ldots \mathfrak{p}_{r}$ mindestens eines, etwa $\mathfrak{p}_{1}$ durch $\mathfrak{q}_{1}$ theilbar und also $= \mathfrak{q}_{1}$ sein, also nach 4.
\[ \mathfrak{p}_{2}\mathfrak{p}_{3}\ldots\mathfrak{p}_{r} = \mathfrak{q}_{2}\mathfrak{q}_{3}\ldots\mathfrak{q}_{s}. \]
Hieraus schliesst man ebenso $\mathfrak{p}_{2} = \mathfrak{q}_{2}$ u.\ s.\ f.

Fasst man in der so gewonnenen Zerlegung die einander gleichen Primideale zu Potenzen zusammen, so kann man setzen
\[ \mathfrak{a} = \mathfrak{p}_{1}^{e_{1}}\mathfrak{p}_{2}^{e_{2}}\ldots\mathfrak{p}_{r}^{e_{r}}. \]

\textbf{*)} Diese Definition des Quotienten zweier Ideale stimmt mit der in §~4, 8. gegebenen völlig überein.

\origpage{216}

Irgend ein Theiler $\mathfrak{a}_{\iota}$ von $\mathfrak{a}$ kann dann durch kein von $\mathfrak{p}_{1}, \mathfrak{p}_{2}, \ldots \mathfrak{p}_{r}$ verschiedenes Primideal und durch keines öfter als $\mathfrak{a}$ theilbar sein. Man erhält also die sämmtlichen Divisoren von $\mathfrak{a}$, deren Anzahl endlich und $= (e_{1}+1)(e_{2}+1)\ldots(e_{r}+1)$ ist, wenn man in
\[ \mathfrak{p}_{1}^{h_{1}}\mathfrak{p}_{2}^{h_{2}}\ldots\mathfrak{p}_{r}^{h_{r}} \]
die Exponenten $h_{\iota}$ die Reihe der Zahlen $0, 1, 2, \ldots e_{\iota}$ durchlaufen lässt (wobei unter $\mathfrak{p}^{0}$ das Ideal $\mathfrak{o}$ zu verstehen ist). Sind $\mathfrak{a}$, $\mathfrak{b}$ zwei Ideale
\[ \mathfrak{a} = \mathfrak{p}_{1}^{e_{1}}\mathfrak{p}_{2}^{e_{2}}\ldots\mathfrak{p}_{r}^{e_{r}}; \quad \mathfrak{b} = \mathfrak{p}_{1}^{f_{1}}\mathfrak{p}_{2}^{f_{2}}\ldots\mathfrak{p}_{r}^{f_{r}} \]
(worin die Exponenten $e$, $f$ auch zum Theil Null sein können), so erhält man den grössten gemeinschaftlichen Theiler und das kleinste gemeinschaftliche Vielfache von $\mathfrak{a}$ und $\mathfrak{b}$ in der Form
\[ \mathfrak{p}_{1}^{g_{1}}\mathfrak{p}_{2}^{g_{2}}\ldots\mathfrak{p}_{r}^{g_{r}}, \]
wenn man für $g_{1}, g_{2}, \ldots g_{r}$ für ersteren die kleinsten, für letzteres die grössten unter den Zahlen $e_{1}, f_{1}; e_{2}, f_{2}; \ldots e_{r}, f_{r}$ nimmt.

6. Sind $\mathfrak{a}$, $\mathfrak{b}$ irgend zwei Ideale, so ist allgemein
\[ N(\mathfrak{a}\mathfrak{b}) = N(\mathfrak{a})N(\mathfrak{b}). \]

\emph{Beweis.} Es sei, wie in 5., $\mathfrak{a} = \mathfrak{p}_{1}\mathfrak{a}_{1}$, so giebt es, weil $\mathfrak{a}_{1}$ ein echter Theiler von $\mathfrak{a}$ ist, in $\mathfrak{a}_{1}$ eine durch $\mathfrak{a}$ nicht theilbare Function $\eta$. Das kleinste gemeinschaftliche Vielfache und der grösste gemeinschaftliche Theiler von $\mathfrak{a}$ und $\mathfrak{o}\eta$ sind bezw. $\mathfrak{p}_{1}\eta$ und $\mathfrak{a}_{1}$, wie sich (nach 5.) sofort aus der Zerlegung von $\mathfrak{a}$ und $\mathfrak{o}\eta$ in ihre Primfactoren ergiebt. Hieraus folgt aber nach §~8, 14.
\[ N(\mathfrak{a}) = N(\mathfrak{p}_{1})N(\mathfrak{a}_{1}). \]
Durch Wiederholung desselben Schlusses für $\mathfrak{a}_{1}$ u.\ s.\ f. ergiebt sich, wenn $\mathfrak{a} = \mathfrak{p}_{1}\mathfrak{p}_{2}\ldots\mathfrak{p}_{r}$ ist:
\[ N(\mathfrak{a}) = N(\mathfrak{p}_{1})N(\mathfrak{p}_{2})\ldots N(\mathfrak{p}_{r}) \]
und daraus
\[ N(\mathfrak{a}\mathfrak{b}) = N(\mathfrak{a})N(\mathfrak{b}). \]

7. Jedes Primideal ist ein Ideal \emph{ersten Grades} (§~7) und umgekehrt, jedes Ideal ersten Grades ist ein Primideal${}^{*)}$.

\emph{Beweis.} Ist $\mathfrak{p}$ ein Primideal, so ist $N(\mathfrak{p})$ durch $\mathfrak{p}$ theilbar, und daher wenigstens einer der Linearfactoren von $N(\mathfrak{p})$, etwa $z-c$, durch $\mathfrak{p}$ theilbar

\textbf{*)} Durch diesen Satz unterscheidet sich die Theorie der algebraischen Functionen wesentlich von der analogen Theorie der algebraischen Zahlen.

\origpage{217}
(§~8, 18.). Ist $\omega$ eine beliebige Function in $\mathfrak{o}$, welche der Gleichung genügt:
\[ \omega^{n} + a_{1}\omega^{n-1} + \cdots + a_{n-1}\omega + a_{n} = 0, \]
so erhält man daraus, indem man die ganzen rationalen Functionen $a_{1}, a_{2}, \ldots a_{n}$ auf ihre constanten Reste $a_{1}^{(0)}, a_{2}^{(0)}, \ldots a_{n}^{(0)}$ nach $z-c$ reducirt, und die ganze Function
\[ \omega^{n} + a_{1}^{(0)}\omega^{n-1} + \cdots + a_{n-1}^{(0)}\omega + a_{n}^{(0)} \]
in ihre Linearfactoren $(\omega-b_{1}), (\omega-b_{2}), \ldots (\omega-b_{n})$ zerlegt:
\[ (\omega-b_{1})(\omega-b_{2})\ldots(\omega-b_{n}) = (z-c)\omega' \equiv 0 \ (\text{mod.}\,\mathfrak{p}). \]

Es muss also wenigstens einer der Factoren $\omega-b_{1}, \omega-b_{2}, \ldots$ durch $\mathfrak{p}$ theilbar sein, d.\ h.\ es ist
\[ \omega \equiv b \ (\text{mod.}\,\mathfrak{p}), \]
worin $b$ eine Constante bedeutet. Da hiernach jede Function in $\mathfrak{o}$ congruent einer Constanten ist $(\text{mod.}\,\mathfrak{p})$, so ist nach §~6 $(\mathfrak{o}, \mathfrak{p}) = N(\mathfrak{p}) = z-c$ eine \emph{lineare} Function von $z$, wodurch der erste Theil der Behauptung erwiesen ist.

Umgekehrt: ist $\mathfrak{q}$ ein Ideal ersten Grades, und
\[ N(\mathfrak{q}) = z-c, \]
so ist $\mathfrak{q}$ gewiss durch ein Primideal $\mathfrak{p}$ theilbar, und da $N(\mathfrak{q})$ durch $N(\mathfrak{p})$ theilbar ist, so ist $N(\mathfrak{p}) = N(\mathfrak{q}) = z-c$, also (§~8, 9.)
\[ \mathfrak{p} = \mathfrak{q}. \]

Es ergiebt sich hieraus, dass der Grad eines Ideals gleich ist der Anzahl der Primfactoren, in welche sich dasselbe zerlegen lässt. Ist daher
\[ \mathfrak{o}(z-c) = \mathfrak{p}_{1}^{e_{1}}\mathfrak{p}_{2}^{e_{2}}\mathfrak{p}_{3}^{e_{3}}\ldots, \]
so ist
\[ e_{1} + e_{2} + e_{3} + \cdots = n. \]
Es folgt ferner noch, dass eine ganze \emph{rationale} Function von $z$ dann und nur dann durch ein Primideal $\mathfrak{p}$ theilbar ist, wenn sie durch die Norm von $\mathfrak{p}$ theilbar ist.

\subsection*{§~10. Die complementären Basen des Körpers $\Omega$.}

1. \emph{Definition.} Bilden die Functionen $\alpha_{1}, \alpha_{2}, \ldots \alpha_{n}$ eine Basis von $\Omega$, und setzt man zur Abkürzung
\[ \begin{gathered} S(\alpha_{r}\alpha_{s}) = a_{r,s} = a_{s,r}, \\ \Delta(\alpha_{1}, \alpha_{2}, \ldots \alpha_{n}) = \Sigma \pm a_{1,1}a_{2,2}\ldots a_{n,n} = a \quad (\text{§~2}), \end{gathered} \]

\origpage{218}
so lässt sich, da $a$ von Null verschieden ist, ein ganz bestimmtes Functionensystem $\alpha'_{1}, \alpha'_{2}, \ldots \alpha'_{n}$ aus den linearen Gleichungen bestimmen
\[ \alpha_{r} = \overset{\iota}{\underset{1,n}{\Sigma}} a_{r,\iota}\alpha'_{\iota}, \tag{1.} \]
und da
\[ \Delta(\alpha'_{1}, \alpha'_{2}, \ldots \alpha'_{n}) = \frac{1}{a} \]
von Null verschieden ist, so bilden die Functionen $\alpha'_{1}, \alpha'_{2}, \ldots \alpha'_{n}$ ebenfalls eine Basis von $\Omega$. Diese soll die zu $\alpha_{1}, \alpha_{2}, \ldots \alpha_{n}$ \emph{complementäre Basis} heissen.

2. Bezeichnet man, wenn die Indices $r$, $s$ der Zahlenreihe $1, 2, \ldots n$ angehören, mit $(r, s)$ die Zahl 1 oder 0, je nachdem $r$, $s$ gleich oder verschieden sind, so ist
\[ S(\alpha_{r}\alpha'_{s}) = (r, s), \tag{2.} \]
denn durch Auflösung der Gleichungen (1.) folgt
\[ \begin{gathered} \alpha'_{s} = \overset{\iota}{\Sigma} a'_{\iota,s}\alpha_{\iota}; \\ a'_{r,s} = a'_{s,r}; \quad \overset{\iota}{\Sigma} a_{r,\iota}a'_{s,\iota} = (r, s), \end{gathered} \]
hieraus:
\[ \alpha_{r}\alpha'_{s} = \overset{\iota}{\Sigma} a'_{\iota,s}\alpha_{\iota}\alpha_{r}; \quad S(\alpha_{r}\alpha'_{s}) = \overset{\iota}{\Sigma} a'_{\iota,s}a_{\iota,r} = (r, s). \]
Genügt umgekehrt ein Functionensystem $\beta_{s}$ den Bedingungen $S(\alpha_{r}\beta_{s}) = (r, s)$, so ist $\beta_{s} = \alpha'_{s}$; denn setzt man $\beta_{s} = \overset{\iota}{\Sigma} b_{\iota,s}\alpha'_{\iota}$, so folgt wegen (2.)
\[ b_{r,s} = S(\beta_{s}\alpha_{r}) = (r, s). \]
Daraus folgt unmittelbar, dass die Beziehung der $\alpha_{\iota}$ zu den $\alpha'_{\iota}$ eine gegenseitige ist, d.\ h.\ dass die Basis $\alpha_{1}, \alpha_{2}, \ldots \alpha_{n}$ complementär ist zu $\alpha'_{1}, \alpha'_{2}, \ldots \alpha'_{n}$.

3. Ist $\eta$ eine beliebige Function in $\Omega$, so kann man stets setzen
\[ \eta = \Sigma x_{\iota}\alpha_{\iota} = \Sigma x'_{\iota}\alpha'_{\iota}, \]
und durch Anwendung von (2.) folgt:
\[ x_{\iota} = S(\eta\alpha'_{\iota}), \quad x' = S(\eta\alpha_{\iota}), \]
\ednote{Printed $x'$ without a subscript; $x'_{\iota}$ is meant, as in the line above.}
also:
\[ \eta = \overset{\iota}{\Sigma}\alpha_{\iota}S(\eta\alpha'_{\iota}) = \overset{\iota}{\Sigma}\alpha'_{\iota}S(\eta\alpha_{\iota}). \tag{3.} \]

4. Ist $\eta$ eine beliebige von Null verschiedene Function in $\Omega$, so ist
\[ \frac{\alpha'_{1}}{\eta}, \ \frac{\alpha'_{2}}{\eta}, \ \ldots \ \frac{\alpha'_{n}}{\eta} \]
die zu $\eta\alpha_{1}, \eta\alpha_{2}, \ldots \eta\alpha_{n}$ complementäre Basis. Dies folgt aus 2. wegen
\[ S\Bigl(\eta\alpha_{r}\cdot\frac{\alpha'_{s}}{\eta}\Bigr) = S(\alpha_{r}\alpha'_{s}) = (r, s). \]

\origpage{219}

5. Wenn zwei Basen von $\Omega$, $\alpha_{1}, \alpha_{2}, \ldots \alpha_{n}$ und $\beta_{1}, \beta_{2}, \ldots \beta_{n}$ durch die $n$ Gleichungen zusammenhängen
\[ \beta_{s} = \overset{\iota}{\Sigma} x_{\iota,s}\alpha_{\iota} \]
mit rationalen Coefficienten $x_{\iota,s}$, so hängen die zu ihnen complementären Basen zusammen durch die $n$ Gleichungen
\[ \alpha'_{r} = \overset{\iota}{\Sigma} x_{r,\iota}\beta'_{\iota} \]
(transponirte Substitution). Es ist dies eine unmittelbare Folge aus 3. wegen
\[ x_{r,s} = S(\alpha'_{r}\beta_{s}). \]

6. Es ist
\[ \overset{\iota}{\Sigma}\alpha_{\iota}\alpha'_{\iota} = 1, \]
also:
\[ \overset{\iota,\iota'}{\Sigma} a_{\iota,\iota'}\alpha'_{\iota}\alpha'_{\iota'} = \overset{\iota,\iota'}{\Sigma} a'_{\iota,\iota'}\alpha_{\iota}\alpha_{\iota'} = 1. \]
Setzt man nämlich zunächst
\[ \Sigma\alpha_{\iota}\alpha_{\iota'} = \sigma, \]
\ednote{Printed $\alpha_{\iota}\alpha_{\iota'}$; the argument requires $\alpha_{\iota}\alpha'_{\iota}$, as in the first formula of 6. The prime appears to have been set on the subscript.}
so folgt aus 3. (auf die Functionen $\eta\alpha_{r}$ angewandt),
\[ \eta\alpha_{r} = \overset{\iota}{\Sigma}\alpha_{\iota}S(\eta\alpha_{r}\alpha'_{\iota}), \]
mithin, nach der Definition der Spur in §~2, (5.)
\[ S(\eta) = \overset{\iota}{\Sigma} S(\eta\alpha_{\iota}\alpha'_{\iota}) = S(\Sigma\eta\alpha_{\iota}\alpha'_{\iota}), \]
also:
\[ S(\eta\sigma) = S(\eta), \]
und daraus, wenn man in 3. einmal $\eta = \sigma$, dann $\eta = 1$ setzt:
\[ \sigma = \overset{\iota}{\Sigma}\alpha_{\iota}S(\sigma\alpha'_{\iota}) = \overset{\iota}{\Sigma}\alpha_{\iota}S(\alpha'_{\iota}) = 1. \]

Wir gehen etwas genauer ein auf die Bildung der complementären Basen in zwei besonderen Fällen:

7. Es sei $\omega_{1}, \omega_{2}, \ldots \omega_{n}$ eine Basis von $\mathfrak{o}$ und $\varepsilon_{1}, \varepsilon_{2}, \ldots \varepsilon_{n}$ die complementäre Basis (von $\Omega$). Es sei
\[ e_{r,s} = e_{s,r} = S(\omega_{r}\omega_{s}), \]
welches \emph{ganze} rationale Functionen sind, und
\[ D = \text{const.}\,\Sigma \pm e_{1,1}e_{2,2}\ldots e_{n,n} \]
die Discriminante von $\Omega$; es ist dann nach 2.
\[ \varepsilon_{r} = \frac{1}{D}\overset{\iota}{\Sigma}\frac{\partial D}{\partial e_{\iota,r}}\omega_{\iota}, \]

\origpage{220}
woraus hervorgeht, dass die Functionen $D\varepsilon_{r}$ zu den \emph{ganzen Functionen} gehören. Aus 6. folgt aber ferner
\[ D = \overset{\iota,\iota'}{\Sigma}\frac{\partial D}{\partial e_{\iota,\iota'}}\omega_{\iota}\omega_{\iota'} \]
und hieraus
\[ \begin{aligned} \varepsilon_{r}\varepsilon_{s} &= \frac{1}{D^{2}}\overset{\iota,\iota'}{\Sigma}\frac{\partial D}{\partial e_{r,\iota}}\frac{\partial D}{\partial e_{s,\iota'}}\omega_{\iota}\omega_{\iota'} \\ &= \frac{1}{D}\frac{\partial D}{\partial e_{r,s}} + \frac{1}{D^{2}}\overset{\iota,\iota'}{\Sigma}\Bigl(\frac{\partial D}{\partial e_{r,\iota}}\frac{\partial D}{\partial e_{s,\iota'}} - \frac{\partial D}{\partial e_{r,s}}\frac{\partial D}{\partial e_{\iota,\iota'}}\Bigr)\omega_{\iota}\omega_{\iota'} \end{aligned} \]
und daraus nach einem bekannten Determinantensatz:
\[ \varepsilon_{r}\varepsilon_{s} = \frac{1}{D}\frac{\partial D}{\partial e_{r,s}} + \frac{1}{D}\overset{\iota,\iota'}{\Sigma}\frac{\partial^{2}D}{\partial e_{r,\iota}\partial e_{s,\iota'}}\omega_{\iota}\omega_{\iota'}, \]
woraus das wichtige Resultat folgt, dass auch die Functionen $D\varepsilon_{r}\varepsilon_{s}$ zu den \emph{ganzen Functionen} gehören.

8. Es sei $\theta$ eine solche Function in $\Omega$, dass $1, \theta, \theta^{2}, \ldots \theta^{n-1}$ eine Basis von $\Omega$ bilden, und es sei also
\[ f(\theta) = \theta^{n} + a_{1}\theta^{n-1} + \cdots + a_{n-1}\theta + a_{n} = 0 \tag{4.} \]
mit rationalen Coefficienten $a$ irreductibel. Es soll die complementäre Basis zu $1, \theta, \theta^{2}, \ldots \theta^{n-1}$ aufgesucht werden. Setzen wir, wenn $t$ eine unbestimmte Constante bedeutet,
\[ \frac{f(t)}{t-\theta} = \eta_{0} + \eta_{1}t + \eta_{2}t^{2} + \cdots + \eta_{n-1}t^{n-1}, \]
also:
\[ \left\{ \begin{aligned} \eta_{0} &= a_{n-1} + a_{n-2}\theta + \cdots + a_{1}\theta^{n-2} + \theta^{n-1}, \\ \eta_{1} &= a_{n-2} + a_{n-3}\theta + \cdots + \theta^{n-2}, \\ &\cdots\cdots\cdots\cdots \\ \eta_{n-2} &= a_{1} + \theta, \\ \eta_{n-1} &= 1, \end{aligned} \right. \tag{5.} \]
so bilden die Functionen $\eta_{0}, \eta_{1}, \ldots \eta_{n-1}$ gleichfalls eine Basis von $\Omega$, da die Determinante der Gleichungen (5.) $= (-1)^{\frac{1}{2}n(n-1)}$, also von Null verschieden ist. Es lässt sich daher jede Function $\zeta$ in $\Omega$ in der Form darstellen
\[ \zeta = y_{0}\eta_{0} + y_{1}\eta_{1} + \cdots + y_{n-1}\eta_{n-1}. \]
Die Reihe der rationalen Functionen $y_{0}, y_{1}, \ldots y_{n-1}$ setzen wir nun fort, indem wir die Functionen $y_{n}, y_{n+1}, \ldots$ durch die Recursion bestimmen
\[ a_{n}y_{r} + a_{n-1}y_{r+1} + \cdots + a_{2}y_{r+n-2} + a_{1}y_{r+n-1} + y_{r+n} = 0. \tag{6.} \]

\origpage{221}
Nun ist nach (5.)
\[ \left\{ \begin{aligned} \theta\eta_{0} &= \phantom{\eta_{0}} - a_{n}\eta_{n-1}, \\ \theta\eta_{1} &= \eta_{0} - a_{n-1}\eta_{n-1}, \\ \theta\eta_{2} &= \eta_{1} - a_{n-2}\eta_{n-1}, \\ &\cdots\cdots\cdots\cdots \\ \theta\eta_{n-1} &= \eta_{n-2} - a_{1}\eta_{n-1}, \end{aligned} \right. \tag{7.} \]
also
\[ \zeta\theta = y_{1}\eta_{0} + y_{2}\eta_{1} + \cdots + y_{n-1}\eta_{n-2} + y_{n}\eta_{n-1}, \]
und ebenso allgemein für jedes ganze positive $r$:
\[ \zeta\theta^{r} = y_{r}\eta_{0} + y_{r+1}\eta_{1} + \cdots + y_{r+n-2}\eta_{n-2} + y_{r+n-1}\eta_{n-1}, \]
oder, wenn man $\eta_{0}, \eta_{1}, \ldots \eta_{n-1}$ durch $1, \theta, \theta^{2}, \ldots \theta^{n-1}$ ausdrückt:
\[ \zeta\theta^{r} = x_{0}^{(r)} + x_{1}^{(r)}\theta + x_{2}^{(r)}\theta^{2} + \cdots + x_{n-1}^{(r)}\theta^{n-1}, \]
worin
\[ \begin{aligned} x_{0}^{(r)} &= y_{r}a_{n-1} + y_{r+1}a_{n-2} + \cdots + y_{r+n-2}a_{1} + y_{r+n-1}, \\ x_{1}^{(r)} &= y_{r}a_{n-2} + y_{r+1}a_{n-3} + \cdots + y_{r+n-2}, \\ &\cdots\cdots\cdots\cdots \\ x_{n-2}^{(r)} &= y_{r}a_{1} + y_{r+1}, \\ x_{n-1}^{(r)} &= y_{r}. \end{aligned} \]
Mithin ist (nach der Definition von $S$, §~2 (5.))
\[ \begin{aligned} S(\zeta) &= x_{0}^{(0)} + x_{1}^{(1)} + x_{2}^{(2)} + \cdots + x_{n-1}^{(n-1)} \\ &= y_{0}a_{n-1} + 2y_{1}a_{n-2} + \cdots + (n-1)y_{n-2}a_{1} + ny_{n-1}, \end{aligned} \]
also, auf $\zeta = \eta_{r}$ angewandt:
\[ S(\eta_{r}) = (r+1)a_{n-1-r}; \quad S(\eta_{n-1-r}) = (n-r)a_{r}, \]
worin $a_{0} = 1$ zu setzen ist.

Setzt man daher zur Abkürzung
\[ S(\theta^{r}) = s_{r}, \]
so folgt, so lange $r \leqq n$, mittels (5.)
\[ (n-r)a_{r} = a_{r}s_{0} + a_{r-1}s_{1} + \cdots + a_{1}s_{r-1} + s_{r} \tag{8.} \]
und nach (4.) allgemein
\[ 0 = a_{n}s_{r} + a_{n-1}s_{r+1} + \cdots + a_{1}s_{r+n-1} + s_{r+n}. \tag{9.} \]
Aus diesen Formeln folgt aber ferner:
\[ \left\{ \begin{aligned} f'(\theta) &= n\theta^{n-1} + (n-1)a_{1}\theta^{n-1} + \cdots + 2a_{n-2}\theta + a_{n-1} \\ &= s_{0}\eta_{0} + s_{1}\eta_{1} + \cdots + s_{n-1}\eta_{n-1}, \\ \theta^{r}f'(\theta) &= s_{r}\eta_{0} + s_{r+1}\eta_{1} + \cdots + s_{r+n-2}\eta_{n-2} + s_{r+n-1}\eta_{n-1}. \end{aligned} \right. \tag{10.} \]
\ednote{In the first line of (10.) the print sets $(n-1)a_{1}\theta^{n-1}$; the derivative requires $\theta^{n-2}$. An apparent misprint.}

\origpage{222}
Beachtet man nun den Werth der Determinante des Gleichungssystems (5.), so folgt hieraus mit Rücksicht auf die Definition der Norm und der Discriminante in §~2 (4.) und (12.) die wichtige Formel
\[ Nf'(\theta) = (-1)^{\frac{1}{2}n(n-1)} \begin{vmatrix} s_{0}, & s_{1}, & \ldots & s_{n-1} \\ s_{1}, & s_{2}, & \ldots & s_{n} \\ \cdots & \cdots & \cdots & \cdots \\ s_{n-1}, & s_{n}, & \ldots & s_{2n-2} \end{vmatrix} = (-1)^{\frac{1}{2}n(n-1)}\Delta(1, \theta, \theta^{2}, \ldots \theta^{n-1}). \tag{11.} \]
Die Gleichungen (10.) ergeben aber auch mit Rücksicht auf die Definition 1. die zu $1, \theta, \theta^{2}, \ldots \theta^{n-1}$ complementäre Basis:
\[ \frac{\eta_{0}}{f'(\theta)}, \ \frac{\eta_{1}}{f'(\theta)}, \ \ldots \ \frac{\eta_{n-1}}{f'(\theta)}. \]

9. Bedeutet $\mathfrak{a} = [\alpha_{1}, \alpha_{2}, \ldots \alpha_{n}]$ einen Modul, dessen Basis zugleich eine Basis von $\Omega$ ist, so erhält man aus der zu $\alpha_{1}, \alpha_{2}, \ldots \alpha_{n}$ complementären Basis von $\Omega$, $\alpha'_{1}, \alpha'_{2}, \ldots \alpha'_{n}$ einen anderen Modul $\mathfrak{a}' = [\alpha'_{1}, \alpha'_{2}, \ldots \alpha'_{n}]$, welcher der zu $\mathfrak{a}$ \emph{complementäre Modul} genannt wird. Derselbe ist, wie sich aus 5. in Verbindung mit §~4, 2. sofort ergiebt, von der Wahl der Basis von $\mathfrak{a}$ unabhängig.

10. Wir betrachten ins Besondere den zu $\mathfrak{o} = [\omega_{1}, \omega_{2}, \ldots \omega_{n}]$ complementären Modul $\mathfrak{e} = [\varepsilon_{1}, \varepsilon_{2}, \ldots \varepsilon_{n}]$. Setzen wir
\[ \omega_{r}\omega_{s} = \overset{\iota}{\Sigma} e_{r,s}^{(\iota)}\omega_{\iota}, \]
so ist nach 3.
\[ e_{r,s}^{(\iota)} = e_{s,r}^{(\iota)} = S(\omega_{r}\omega_{s}\varepsilon_{\iota}) \]
eine ganze rationale Function von $z$, und es folgt:
\[ \omega_{r}\varepsilon_{s} = \overset{\iota}{\Sigma} e_{r,\iota}^{(s)}\varepsilon_{\iota}. \]
Hieraus ergiebt sich, dass der Modul $\mathfrak{o}\mathfrak{e}$ (§~4, 7.) theilbar ist durch $\mathfrak{e}$; andererseits ist, weil $\mathfrak{o}$ die Function 1 enthält, $\mathfrak{e}$ theilbar durch $\mathfrak{o}\mathfrak{e}$, also
\[ \mathfrak{o}\mathfrak{e} = \mathfrak{e}, \]
d.\ h.\ der Modul $\mathfrak{e}$, der zwar nicht bloss ganze Functionen enthält, besitzt die charakteristische Eigenschaft II.\ §~7 der Ideale. Dasselbe gilt in Folge dessen auch von dem Modul $\mathfrak{e}^{2}$. Da die beiden Moduln $D\mathfrak{e}$, $D\mathfrak{e}^{2}$ nach 7. nur ganze Functionen enthalten, so sind dieselben Ideale, und aus 7. ergiebt sich noch
\[ N(D\mathfrak{e}) = D^{n-1}. \]

\origpage{223}

11. Ist $\theta$ eine Function in $\mathfrak{o}$ von der Art, dass $1, \theta, \theta^{2}, \ldots \theta^{n-1}$ eine Basis von $\Omega$ bildet, so dass in der irreductibeln Gleichung
\[ f(\theta) = \theta^{n} + a_{1}\theta^{n-1} + \cdots + a_{n-1}\theta + a_{n} = 0 \]
die Coefficienten \emph{ganze} rationale Functionen von $z$ sind, so kann man für $r = 0, 1, 2, \ldots n-1$ die ganzen rationalen Functionen $k_{\iota}^{(r)}$ so bestimmen, dass
\[ \theta^{r} = \overset{\iota}{\underset{1,n}{\Sigma}} k_{\iota}^{(r)}\omega_{\iota} \]
wird. Wendet man hierauf den Satz 5. und 8. an, so folgt:
\[ f'(\theta)\varepsilon_{s} = k_{s}^{(0)}\eta_{0} + k_{s}^{(1)}\eta_{1} + \cdots + k_{s}^{(n-1)}\eta_{n-1}, \]
und hieraus ergiebt sich, dass der Modul
\[ f'(\theta)\mathfrak{e} = \mathfrak{k} \]
nur \emph{ganze} Functionen enthält. Aus 10. schliesst man, dass derselbe ein \emph{Ideal} ist.

\subsection*{§~11. Das Verzweigungsideal.}

1. \emph{Hülfssatz.} Sind je zwei der Ideale $\mathfrak{a}$, $\mathfrak{b}$, $\mathfrak{c}$, $\ldots$ relativ prim, so existirt immer eine Function, welche in Bezug auf jedes von ihnen einer gegebenen Function in $\mathfrak{o}$ congruent ist.

\emph{Beweis.} Man setze
\[ \mathfrak{m} = \mathfrak{a}\mathfrak{b}\mathfrak{c}\cdots = \mathfrak{a}\mathfrak{a}_{1} = \mathfrak{b}\mathfrak{b}_{1} = \mathfrak{c}\mathfrak{c}_{1} = \cdots; \]
der grösste gemeinschaftliche Theiler von $\mathfrak{a}_{1} = \mathfrak{b}\mathfrak{c}\ldots$, $\mathfrak{b}_{1} = \mathfrak{a}\mathfrak{c}\ldots$, $\mathfrak{c}_{1} = \mathfrak{a}\mathfrak{b}\ldots$ ist alsdann gleich $\mathfrak{o}$, da kein Primideal zugleich in $\mathfrak{a}_{1}, \mathfrak{b}_{1}, \mathfrak{c}_{1}, \ldots$ aufgehen kann. Folglich kann man (§~4, 5.) $\alpha_{1}$ aus $\mathfrak{a}_{1}$, $\beta_{1}$ aus $\mathfrak{b}_{1}$, $\gamma_{1}$ aus $\mathfrak{c}_{1}$, $\ldots$ so auswählen, dass
\[ \alpha_{1} + \beta_{1} + \gamma_{1} + \cdots = 1, \]
also:
\[ \begin{aligned} \alpha_{1} &\equiv 1, & \beta_{1} &\equiv 0, & \gamma_{1} &\equiv 0, & &\ldots \quad (\text{mod.}\,\mathfrak{a}), \\ \alpha_{1} &\equiv 0, & \beta_{1} &\equiv 1, & \gamma_{1} &\equiv 0, & &\ldots \quad (\text{mod.}\,\mathfrak{b}), \\ \alpha_{1} &\equiv 0, & \beta_{1} &\equiv 0, & \gamma_{1} &\equiv 1, & &\ldots \quad (\text{mod.}\,\mathfrak{c}), \\ &\cdots\cdots\cdots\cdots \end{aligned} \]
Sind daher $\lambda, \mu, \nu, \ldots$ gegebene Functionen in $\mathfrak{o}$, so genügt
\[ \omega \equiv \lambda\alpha_{1} + \mu\beta_{1} + \nu\gamma_{1} + \cdots \ (\text{mod.}\,\mathfrak{m}) \]
den Bedingungen
\[ \omega \equiv \lambda \ (\text{mod.}\,\mathfrak{a}), \quad \omega \equiv \mu \ (\text{mod.}\,\mathfrak{b}), \quad \omega \equiv \nu \ (\text{mod.}\,\mathfrak{c}), \quad \ldots \]

\origpage{224}

2. Es seien $\mathfrak{p}, \mathfrak{p}_{1}, \mathfrak{p}_{2}, \ldots$ die sämmtlichen von einander verschiedenen in einer beliebigen linearen Function $z-c$ aufgehenden Primideale und
\[ \mathfrak{o}(z-c) = \mathfrak{p}^{e}\mathfrak{p}_{1}^{e_{1}}\mathfrak{p}_{2}^{e_{2}}\ldots, \quad e + e_{1} + e_{2} + \cdots = n \quad (\text{§~9, 7.}). \]
Man wähle die Functionen $\lambda, \lambda_{1}, \lambda_{2}, \ldots$ theilbar resp.\ durch $\mathfrak{p}, \mathfrak{p}_{1}, \mathfrak{p}_{2}, \ldots$, aber nicht durch $\mathfrak{p}^{2}, \mathfrak{p}_{1}^{2}, \mathfrak{p}_{2}^{2}, \ldots$ und lasse $b, b_{1}, b_{2}, \ldots$ beliebige, jedoch von einander verschiedene Constanten bedeuten. Nach 1. lässt sich dann eine Function $\zeta$ bestimmen, welche den Congruenzen genügt
\[ \zeta \equiv b + \lambda \ (\text{mod.}\,\mathfrak{p}^{2}), \quad \zeta \equiv b_{1} + \lambda_{1} \ (\text{mod.}\,\mathfrak{p}_{1}^{2}), \quad \zeta \equiv b_{2} + \lambda_{2} \ (\text{mod.}\,\mathfrak{p}_{2}^{2}), \quad \ldots, \]
also
\[ \zeta \equiv b \ (\text{mod.}\,\mathfrak{p}), \quad \zeta \equiv b_{1} \ (\text{mod.}\,\mathfrak{p}_{1}), \quad \zeta \equiv b_{2} \ (\text{mod.}\,\mathfrak{p}_{2}), \quad \ldots, \]
so dass, wenn $a$ irgend eine Constante bedeutet, $\zeta-a$ höchstens durch eines der Primideale $\mathfrak{p}, \mathfrak{p}_{1}, \mathfrak{p}_{2}, \ldots$ und niemals durch eines ihrer Quadrate theilbar ist. Ist daher $\varphi(t) = \Pi(t-a)$ eine ganze Function der Variablen $t$ mit constanten Coefficienten, so ist $\varphi(\zeta) = \Pi(\zeta-a)$ stets und nur dann durch $\mathfrak{p}^{m}$ theilbar, wenn $\varphi(t)$ algebraisch durch $(t-b)^{m}$ theilbar ist, und wenn $\mathfrak{p}^{m}$ die höchste in $\varphi(\zeta)$ aufgehende Potenz von $\mathfrak{p}$ ist, so ist folglich $\mathfrak{p}^{m-1}$ die höchste in $\varphi'(\zeta)$ aufgehende Potenz von $\mathfrak{p}$. Soll daher $\varphi(\zeta)$ durch $z-c$ theilbar sein, so muss $\varphi(t)$ durch die Function $n^{\text{ten}}$ Grades
\[ \psi(t) = (t-b)^{e}(t-b_{1})^{e_{1}}(t-b_{2})^{e_{2}}\ldots \]
theilbar sein. Mithin ist die Congruenz
\[ x_{0} + x_{1}\zeta + x_{2}\zeta^{2} + \cdots + x_{n-1}\zeta^{n-1} \equiv 0 \quad (\text{mod.}\,z-c) \]
nur durch solche ganze rationale $x$ zu befriedigen, die alle durch $z-c$ theilbar sind. Setzt man also, indem man mit $k_{1}^{(0)}, k_{1}^{(1)}, \ldots$ ganze rationale Functionen von $z$ und mit $\omega_{1}, \omega_{2}, \ldots \omega_{n}$ eine Basis von $\mathfrak{o}$ bezeichnet:
\[ \begin{aligned} 1 &= k_{1}^{(0)}\omega_{1} + k_{2}^{(0)}\omega_{2} + \cdots + k_{n}^{(0)}\omega_{n}, \\ \zeta &= k_{1}^{(1)}\omega_{1} + k_{2}^{(1)}\omega_{2} + \cdots + k_{n}^{(1)}\omega_{n}, \\ \zeta^{2} &= k_{1}^{(2)}\omega_{1} + k_{2}^{(2)}\omega_{2} + \cdots + k_{n}^{(2)}\omega_{n}, \\ &\cdots\cdots\cdots\cdots \\ \zeta^{n-1} &= k_{1}^{(n-1)}\omega_{1} + k_{2}^{(n-1)}\omega_{2} + \cdots + k_{n}^{(n-1)}\omega_{n}, \end{aligned} \]
so kann die Determinante
\[ k = \Sigma \pm k_{1}^{(0)}k_{2}^{(1)}\ldots k_{n}^{(n-1)} \]
weder identisch verschwinden, noch durch $z-c$ theilbar sein (vgl.\ die Note zu §~9, 1.).

\origpage{225}

Es folgt also, dass
\[ N(t-\zeta) = f(t, z) \]
irreductibel ist. Da nun $f(\zeta, z) = 0$, also $f(\zeta, c)$ durch $z-c$ theilbar ist, so muss $f(t, c)$ durch $\psi(t)$ theilbar, also, da beide Functionen von gleichem Grade sind,
\[ f(t, c) = \psi(t) \]
sein, woraus man noch für eine folgende Anwendung schliesst:
\[ S(\zeta) \equiv eb + e_{1}b_{1} + e_{2}b_{2} + \cdots \quad (\text{mod.}\,z-c), \]
und, indem man dieselbe Betrachtung auf die Functionen $\zeta^{2}, \zeta^{3}, \ldots$ anwendet, was falls keine der Constanten $b$ verschwindet, sicher gestattet ist:
\[ \begin{aligned} S(\zeta^{2}) &\equiv eb^{2} + e_{1}b_{1}^{2} + e_{2}b_{2}^{2} + \cdots \quad (\text{mod.}\,z-c), \\ S(\zeta^{3}) &\equiv eb^{3} + e_{1}b_{1}^{3} + e_{2}b_{2}^{3} + \cdots \quad (\text{mod.}\,z-c). \\ &\cdots\cdots\cdots\cdots \end{aligned} \]

Es ist also $\mathfrak{p}^{e}$ die höchste in $f(\zeta, c)$, also $\mathfrak{p}^{e-1}$ die höchste in $f'(\zeta, c)$ aufgehende Potenz von $\mathfrak{p}$, und da
\[ f'(\zeta, c) \equiv f'(\zeta, z) \quad (\text{mod.}\,\mathfrak{p}^{e}), \]
so ist $\mathfrak{p}^{e-1}$ auch die höchste in $f'(\zeta, z)$ aufgehende Potenz von $\mathfrak{p}$. Hieraus ergiebt sich
\[ \mathfrak{o}f'(\zeta, z) = \mathfrak{m}\mathfrak{p}^{e-1}\mathfrak{p}_{1}^{e_{1}-1}\ldots, \]
worin $\mathfrak{m}$ und folglich auch $N(\mathfrak{m})$ relativ prim zu $z-c$ ist.

Ist nun $D$ die Discriminante von $\Omega$, so ist hiernach und nach §~10, (11.) und §~2, (13.) (von constanten Factoren abgesehen)
\[ Nf'(\zeta, z) = \Delta(1, \zeta, \zeta^{2}, \ldots \zeta^{n-1}) = Dk^{2} = (z-c)^{n-s}N(\mathfrak{m}), \]
wenn $s$ die Anzahl der verschiedenen in $z-c$ aufgehenden Primideale $\mathfrak{p}, \mathfrak{p}_{1}, \mathfrak{p}_{2}, \ldots$ bedeutet; und da $k$ und $N(\mathfrak{m})$ durch $z-c$ nicht theilbar sind, so ist $(z-c)^{n-s}$ die höchste in $D$ aufgehende Potenz von $z-c$. Folglich:
\[ D = \Pi(z-c)^{n-s}, \tag{1.} \]
worin das Productzeichen $\Pi$ sich auf alle solche linearen Ausdrücke $z-c$ bezieht, in denen weniger als \emph{$n$ verschiedene} Primfactoren aufgehen, die also durch die zweite oder eine höhere Potenz eines Primideals theilbar sind.

\emph{Es giebt also nur eine endliche Anzahl linearer Functionen $z-c$, die durch das Quadrat eines Primideals theilbar sind.}

Wir setzen nun
\[ \mathfrak{z} = \Pi\mathfrak{p}^{e-1}, \tag{2.} \]
worin sich das Productzeichen $\Pi$ auf alle diejenigen Primideale $\mathfrak{p}$ bezieht,

\origpage{226}
von denen eine höhere als die erste, nämlich die $e^{\text{te}}$ Potenz in ihrer Norm aufgeht, und nennen dieses Ideal $\mathfrak{z}$ das \emph{Verzweigungsideal}. Aus (1.) und (2.) folgt sofort
\[ N(\mathfrak{z}) = D. \tag{3.} \]
Da ferner $n-s \geqq e-1$, also $e(n-s) - 2(e-1) \geqq (e-1)(e-2) \geqq 0$ ist, so ist $D$ theilbar durch $\mathfrak{p}^{2(e-1)}$, also auch durch $\mathfrak{z}^{2}$, und man kann, wenn man mit $\mathfrak{d}$ gleichfalls ein Ideal bezeichnet, setzen:
\[ \mathfrak{o}D = \mathfrak{d}\mathfrak{z}^{2}, \quad N(\mathfrak{d}) = D^{n-2}. \tag{4.} \]

3. Ist eine Function $\varrho$ in $\mathfrak{o}$ durch jedes in $z-c$ aufgehende Primideal theilbar, so ist $S(\varrho)$ durch $z-c$ theilbar.

\emph{Beweis.} Es sei $\zeta$ dieselbe Function wie in 2., so dass man setzen kann:
\[ x\varrho = x_{0} + x_{1}\zeta + x_{2}\zeta^{2} + \cdots + x_{n-1}\zeta^{n-1}, \]
worin die Coefficienten $x, x_{0}, x_{1}, \ldots x_{n-1}$ ganze rationale Functionen von $z$ ohne gemeinsamen Theiler sind, von denen die erste durch $z-c$ nicht theilbar ist (vgl.\ 2.). Aus unserer Voraussetzung über die Function $\varrho$ folgt, wenn die Constanten $b$ dieselbe Bedeutung wie in 2. haben
\[ \begin{aligned} x_{0} + x_{1}b + x_{2}b^{2} + \cdots + x_{n-1}b^{n-1} &\equiv 0 \quad (\text{mod.}\,z-c), \\ x_{0} + x_{1}b_{1} + x_{2}b_{1}^{2} + \cdots + x_{n-1}b_{1}^{n-1} &\equiv 0 \quad (\text{mod.}\,z-c), \\ x_{0} + x_{1}b_{2} + x_{2}b_{2}^{2} + \cdots + x_{n-1}b_{2}^{n-1} &\equiv 0 \quad (\text{mod.}\,z-c), \\ \cdots\cdots\cdots\cdots& \end{aligned} \]
und hieraus, indem man die Congruenzen mit $e, e_{1}, e_{2}, \ldots$ multiplicirt und addirt:
\[ x_{0}n + x_{1}S(\zeta) + x_{2}S(\zeta^{2}) + \cdots + x_{n-1}S(\zeta^{n-1}) = xS(\varrho) \equiv 0 \quad (\text{mod.}\,z-c), \]
also, da $x$ durch $z-c$ nicht theilbar ist,
\[ S(\varrho) \equiv 0 \quad (\text{mod.}\,z-c) \]
w.\ z.\ b.\ w.

4. Es sei jetzt
\[ r = (z-c)(z-c_{1})(z-c_{2})\ldots \]
das Product sämmtlicher von einander verschiedenen Linearfactoren von $D$ und
\[ \mathfrak{r} = \mathfrak{p}\mathfrak{p}_{1}\mathfrak{p}_{2}\ldots \]
das Product der sämmtlichen von einander verschiedenen in $r$ aufgehenden Primideale. Ist wie oben $\mathfrak{z}$ das Verzweigungsideal, so ist
\[ \mathfrak{r}\mathfrak{z} = \Pi\mathfrak{p}^{e} = \mathfrak{o}r \tag{5.} \]
und mithin
\[ N(\mathfrak{r}) = \frac{r^{n}}{D}. \]

\origpage{227}
Jede Function $\varrho$ in $\mathfrak{r}$ hat nach 3. die Eigenschaft, dass $S(\varrho)$ durch $r$ theilbar ist. Ist nun, wie in §~10
\[ \mathfrak{e} = [\varepsilon_{1}, \varepsilon_{2}, \ldots \varepsilon_{n}] \]
der zu $\mathfrak{o}$ complementäre Modul, $\varrho$ eine beliebige Function in $\mathfrak{r}$, so kann man setzen
\[ \varrho = x_{1}\varepsilon_{1} + x_{2}\varepsilon_{2} + \cdots + x_{n}\varepsilon_{n}, \]
worin nach §~10, 3.
\[ x_{\iota} = S(\varrho\omega_{\iota}), \]
also, da $\varrho\omega_{\iota}$ eine Function in $\mathfrak{r}$ ist, $x_{\iota}$ eine durch $r$ theilbare ganze rationale Function von $z$. Hieraus folgt, dass das Ideal $\mathfrak{r}$ durch den Modul $r\mathfrak{e}$ theilbar ist. Es ist also auch das Ideal $D\mathfrak{r}$ theilbar durch das Ideal $rD\mathfrak{e}$. Zugleich ist
\[ N(D\mathfrak{r}) = r^{n}D^{n-1}, \quad N(rD\mathfrak{e}) = r^{n}D^{n-1} \qquad (\text{§~10, 10.}); \]
mithin nach §~8, 9.
\[ D\mathfrak{r} = rD\mathfrak{e} \]
oder
\[ \mathfrak{r} = r\mathfrak{e}. \tag{6.} \]
Woraus mittelst der obigen Bemerkung über $\varrho$ folgt, dass, wenn $\varepsilon$ eine beliebige Function in $\mathfrak{e}$ bedeutet, $S(\varepsilon)$ eine \emph{ganze} rationale Function von $z$ ist. Aus der Formel (6.) folgt durch Multiplication mit $\mathfrak{z}$ nach (5.)
\[ \mathfrak{r}\mathfrak{z} = r\mathfrak{e}\mathfrak{z} = \mathfrak{o}r \]
und folglich
\[ \mathfrak{e}\mathfrak{z} = \mathfrak{o}. \tag{7.} \]
Multiplicirt man die letzte Gleichung mit $D$, so ergiebt sich aus (4.)
\[ \mathfrak{e}D\mathfrak{z} = \mathfrak{z}^{2}\mathfrak{d}, \]
folglich
\[ D\mathfrak{e} = \mathfrak{z}\mathfrak{d} \tag{8.} \]
und durch Multiplication dieser Gleichung mit $\mathfrak{e}$ nach (7.)
\[ D\mathfrak{e}^{2} = \mathfrak{d}. \tag{9.} \]

4.\ednote{The print numbers this item 4. a second time, after item 4. on p.~226; 5. would be expected.} Ist $\theta$ eine ganze Function von $z$ in $\Omega$ und $N(t-\theta) = f(t)$, so ist $f'(\theta)$ theilbar durch das Verzweigungsideal $\mathfrak{z}$.

\emph{Beweis.} Ist $f(t)$ reductibel, so ist $f'(\theta) = 0$, also sicher theilbar durch $\mathfrak{z}$. Im anderen Fall ist nach §~10, 11.
\[ \mathfrak{e}f'(\theta) = \mathfrak{k} \]

\origpage{228}
ein Ideal, folglich durch Multiplication mit $\mathfrak{z}$ nach (7.)
\[ \mathfrak{o}f'(\theta) = \mathfrak{k}\mathfrak{z}. \tag{10.} \]
Zugleich folgt, wenn wir wie in §~10, 11.
\[ \begin{aligned} \theta^{r} &= \overset{\iota}{\Sigma} k_{\iota}^{(r)}\omega_{\iota}, \\ k &= \Sigma \pm k_{1}^{(0)}k_{2}^{(1)}\ldots k_{n}^{(n-1)} \end{aligned} \]
setzen,
\[ \begin{aligned} Nf'(\theta) &= N(\mathfrak{k})N(\mathfrak{z}) = DN(\mathfrak{k}) \\ &= \text{const.}\,k^{2}D \quad (\text{§~10, (11.) und §~2, (13.)}), \end{aligned} \]
also:
\[ N(\mathfrak{k}) = \text{const.}\,k^{2} \tag{11.} \]
ein vollständiges Quadrat.

\subsection*{§~12. Die gebrochenen Functionen von $z$ im Körper $\Omega$.}

1. Jede beliebige Function $\eta$ in $\Omega$ lässt sich nach §~3, 3. auf unendlich viele Arten als Quotient zweier ganzen Functionen von $z$ darstellen (der Nenner kann sogar eine ganze rationale Function von $z$ sein). Es sei also
\[ \eta = \frac{\nu}{\mu} \]
und $\mu$, $\nu$ ganze Functionen von $z$ (Functionen in $\mathfrak{o}$). Ist nun $\mathfrak{m}$ der grösste gemeinschaftliche Theiler der beiden Hauptideale $\mathfrak{o}\mu$, $\mathfrak{o}\nu$, also, wenn $\mathfrak{a}$, $\mathfrak{b}$ relative Primideale sind,
\[ \mathfrak{o}\mu = \mathfrak{a}\mathfrak{m}, \quad \mathfrak{o}\nu = \mathfrak{b}\mathfrak{m}, \tag{1.} \]
so folgt (§~4, 6.)
\[ \mathfrak{a}\nu = \mathfrak{b}\mu \quad \text{oder} \quad \mathfrak{a}\eta = \mathfrak{b}. \tag{2.} \]

Ist also $\alpha$ eine beliebige Function in $\mathfrak{a}$, so ist $\alpha\eta$ in $\mathfrak{b}$ enthalten, also jedenfalls eine ganze Function von $z$. Ist umgekehrt $\alpha$ eine ganze Function von $z$, welche die Eigenschaft hat, dass $\alpha\eta = \beta$ eine ganze Function ist, so folgt
\[ \alpha\nu = \beta\mu, \]
also nach (1.)
\[ \alpha\mathfrak{b} = \beta\mathfrak{a}; \]
da nun $\mathfrak{a}$, $\mathfrak{b}$ relativ prim sind, so muss $\alpha$ durch $\mathfrak{a}$, $\beta$ durch $\mathfrak{b}$ theilbar sein, und daraus folgt:

\origpage{229}

Es ist $\mathfrak{a}$ der Inbegriff aller derjenigen ganzen Functionen $\alpha$, welche die Eigenschaft haben, dass $\alpha\eta$ eine ganze Function ist, und der Inbegriff aller dieser ganzen Functionen $\alpha\eta$ ist das Ideal $\mathfrak{b}$; oder anders ausgedrückt:

\emph{Es ist $\mathfrak{b}$ das kleinste gemeinschaftliche Vielfache von $\mathfrak{o}\eta$ und $\mathfrak{o}$, ebenso $\mathfrak{a}$ das kleinste gemeinschaftliche Vielfache von $\frac{\mathfrak{o}}{\eta}$ und $\mathfrak{o}$.} Hiernach muss, wenn $\mathfrak{a}'$, $\mathfrak{b}'$ zwei der Bedingung
\[ \mathfrak{a}'\eta = \mathfrak{b}' \]
genügende Ideale sind, $\mathfrak{a}'$ durch $\mathfrak{a}$ theilbar sein. Sei also
\[ \mathfrak{a}' = \mathfrak{n}\mathfrak{a}, \]
so folgt:
\[ \mathfrak{b}' = \mathfrak{n}\mathfrak{a}\eta = \mathfrak{n}\mathfrak{b}. \]
Umgekehrt ist auch für ein beliebiges Ideal $\mathfrak{n}$
\[ \mathfrak{n}\mathfrak{a}\eta = \mathfrak{n}\mathfrak{b}. \]

2. Es seien jetzt $\mathfrak{a}$, $\mathfrak{b}$ zwei der Bedingung
\[ \mathfrak{a}\eta = \mathfrak{b} \]
genügende Ideale, gleichviel ob relativ prim oder nicht. Der Quotient $\frac{\mathfrak{b}}{\mathfrak{a}}$ ist nach §~4, 8. der Inbegriff aller derjenigen Functionen $\gamma$, welche die Eigenschaft haben, dass $\mathfrak{a}\gamma$ durch $\mathfrak{b}$ theilbar ist. Zu diesen Functionen gehören gewiss alle Functionen von der Form $\omega\eta$, wenn $\omega$ eine beliebige Function in $\mathfrak{o}$ bedeutet. Aber auch umgekehrt ist jede Function $\gamma$ von dieser Form; denn da $\mathfrak{a}\gamma$ durch $\mathfrak{b}$, also auch durch $\mathfrak{o}$ theilbar ist, so ist es ein Ideal (da es die Eigenschaften I., II. §~7 besitzt), also wenn $\mathfrak{c}$ gleichfalls ein Ideal ist:
\[ \mathfrak{a}\gamma = \mathfrak{c}\mathfrak{b} \]
und durch Multiplication mit $\eta$
\[ \mathfrak{b}\gamma = \mathfrak{c}\mathfrak{b}\eta. \]
Ist nun wie oben $\eta = \frac{\nu}{\mu}$, und, wenn $\varrho$, $\sigma$ ganze Functionen sind, $\gamma = \frac{\varrho}{\sigma}$, so folgt hieraus
\[ \mathfrak{b}\varrho\mu = \mathfrak{c}\mathfrak{b}\nu\sigma, \]
also:
\[ \mathfrak{o}\varrho\mu = \mathfrak{c}\nu\sigma, \quad \mathfrak{o}\gamma = \mathfrak{c}\eta. \]
Beides zusammen liefert den Satz
\[ \mathfrak{o}\eta = \frac{\mathfrak{b}}{\mathfrak{a}}. \tag{3.} \]

\origpage{230}
Sind in dieser Darstellung $\mathfrak{b}$, $\mathfrak{a}$ relativ prim, was nach 1. stets und nur auf eine Weise angenommen werden kann, so soll $\mathfrak{b}$ das \emph{Oberideal}, $\mathfrak{a}$ das \emph{Unterideal} der Function $\eta$ heissen.

3. Ist wieder allgemein
\[ \mathfrak{a}\eta = \mathfrak{b}, \quad \text{also} \quad \mathfrak{o}\eta = \frac{\mathfrak{b}}{\mathfrak{a}} \]
und $\alpha$ eine beliebige Function in $\mathfrak{a}$, $\beta$ eine zugehörige Function in $\mathfrak{b}$, so ist
\[ \eta = \frac{\beta}{\alpha} \quad \text{und} \quad \mathfrak{a}\beta = \mathfrak{b}\alpha. \]
Hieraus folgt durch Bildung der Normen
\[ N(\eta) = \text{const.}\,\frac{N(\mathfrak{a})}{N(\mathfrak{b})}. \]
\ednote{As printed, with $N(\mathfrak{a})$ in the numerator. Since $\mathfrak{a}\eta = \mathfrak{b}$, the quotient $N(\mathfrak{b})$ over $N(\mathfrak{a})$ would be expected; this appears to be a misprint.}

4. Sind $\eta$, $\eta'$ zwei Functionen in $\Omega$ und ist wie in 1.
\[ \mathfrak{a}\eta = \mathfrak{b}; \quad \mathfrak{a}'\eta' = \mathfrak{b}', \]
gleichviel ob $\mathfrak{a}$, $\mathfrak{b}$; $\mathfrak{a}'$, $\mathfrak{b}'$ relativ prim sind oder nicht, so folgt
\[ \mathfrak{a}\mathfrak{a}'\eta\eta' = \mathfrak{b}\mathfrak{b}'. \]
Es folgen also aus
\[ \mathfrak{o}\eta = \frac{\mathfrak{b}}{\mathfrak{a}}, \quad \mathfrak{o}\eta' = \frac{\mathfrak{b}'}{\mathfrak{a}'} \]
die Gleichungen
\[ \mathfrak{o}\eta\eta' = \frac{\mathfrak{b}\mathfrak{b}'}{\mathfrak{a}\mathfrak{a}'}, \quad \mathfrak{o}\frac{1}{\eta} = \frac{\mathfrak{a}}{\mathfrak{b}}, \quad \mathfrak{o}\frac{\eta}{\eta'} = \frac{\mathfrak{b}\mathfrak{a}'}{\mathfrak{a}\mathfrak{b}'}. \]

5. Ist $\mathfrak{a}\eta = \mathfrak{b}$, $\mathfrak{a}\eta' = \mathfrak{b}'$, so wird auch
\[ \mathfrak{a}(\eta \pm \eta') = \mathfrak{b}'' \]
ein Ideal sein, weil, wenn $\alpha\eta$, $\alpha\eta'$ ganze Functionen sind, stets auch $\alpha(\eta \pm \eta')$ eine solche ist. Ist also
\[ \mathfrak{o}\eta = \frac{\mathfrak{b}}{\mathfrak{a}}, \quad \mathfrak{o}\eta' = \frac{\mathfrak{b}'}{\mathfrak{a}}, \]
so folgt
\[ \mathfrak{o}(\eta \pm \eta') = \frac{\mathfrak{b}''}{\mathfrak{a}}; \]
haben die beiden Ideale $\mathfrak{b}$, $\mathfrak{b}'$ einen gemeinsamen Theiler, so ist derselbe auch Theiler von $\mathfrak{b}''$.

6. Es sei jetzt $\varrho$ eine Function in $\Omega$, deren Oberideal durch ein beliebig gegebenes Primideal $\mathfrak{p}$, aber nicht durch $\mathfrak{p}^{2}$ theilbar ist (solche Functionen

\origpage{231}
existiren stets; dieselben können sogar ganze Functionen von $z$ sein), also
\[ \mathfrak{o}\varrho = \frac{\mathfrak{m}\mathfrak{p}}{\mathfrak{n}}, \]
worin $\mathfrak{m}$, $\mathfrak{n}$ durch $\mathfrak{p}$ nicht theilbare Ideale sind. Sei ferner $\eta$ eine beliebige Function in $\Omega$, deren Unterideal durch $\mathfrak{p}$ nicht theilbar ist, also
\[ \mathfrak{o}\eta = \frac{\mathfrak{b}}{\mathfrak{a}} \]
und $\mathfrak{a}$ nicht theilbar durch $\mathfrak{p}$. Man wähle eine beliebige Function $\alpha$ in $\mathfrak{a}$, die nicht durch $\mathfrak{p}$ theilbar ist, und eine entsprechende Function $\beta$ in $\mathfrak{b}$, so dass
\[ \eta = \frac{\beta}{\alpha} \]
wird. Sei
\[ \alpha \equiv \alpha_{0}, \quad \beta \equiv \beta_{0} \ (\text{mod.}\,\mathfrak{p}), \quad c_{0} = \frac{\beta_{0}}{\alpha_{0}}, \]
worin $\alpha_{0}$, $\beta_{0}$, $c_{0}$ Constanten sind, deren erste von Null verschieden ist. Nach 5. ist
\[ \mathfrak{o}(\eta - c_{0}) = \mathfrak{o}\frac{\beta - c_{0}\alpha}{\alpha} = \frac{\mathfrak{b}_{1}}{\mathfrak{a}}, \]
und aus
\[ \mathfrak{a}(\beta - c_{0}\alpha) = \mathfrak{b}_{1}\alpha, \quad \beta - c_{0}\alpha \equiv 0 \ (\text{mod.}\,\mathfrak{p}) \]
folgt, da $\alpha$ durch $\mathfrak{p}$ nicht theilbar ist, dass $\mathfrak{b}_{1}$ durch $\mathfrak{p}$ theilbar sein muss. Setzt man also
\[ \eta - c_{0} = \varrho\eta_{1}, \]
so ist auch das Unterideal von $\eta_{1}$ durch $\mathfrak{p}$ nicht theilbar. Auf diese Weise lässt sich eine ganz bestimmte Reihe von Constanten $c_{0}, c_{1}, \ldots c_{r-1}, \ldots$ der Art ermitteln, dass
\[ \begin{aligned} \eta &= c_{0} + \varrho\eta_{1}, \\ \eta_{1} &= c_{1} + \varrho\eta_{2}, \\ &\cdots\cdots\cdots\cdots \\ \eta_{r-1} &= c_{r-1} + \varrho\eta_{r}, \quad \ldots, \end{aligned} \]
worin die $\eta_{1}, \eta_{2}, \ldots \eta_{r}, \ldots$ Functionen bedeuten, deren Unterideale keine andern Primfactoren haben können als das Unterideal von $\eta$ und das Oberideal von $\varrho$ mit Ausschluss von $\mathfrak{p}$. Demnach ist für jedes ganze positive $r$
\[ \eta = c_{0} + c_{1}\varrho + \cdots + c_{r-1}\varrho^{r-1} + \eta_{r}\varrho^{r}. \]
Ist das Unterideal von $\zeta$ durch $\mathfrak{p}^{s}$, nicht durch $\mathfrak{p}^{s+1}$ theilbar, so kann man dieselbe Betrachtung auf die Function $\eta = \zeta\varrho^{s}$ anwenden und erhält
\[ \zeta = c_{0}\varrho^{-s} + c_{1}\varrho^{-s+1} + \cdots + c_{r-1}\varrho^{-s+r} + \eta_{r}\varrho^{-s+r}. \]
\ednote{Printed $c_{r-1}\varrho^{-s+r}$; by the preceding expansion $c_{r-1}\varrho^{-s+r-1}$ is expected, an apparent misprint.}

\origpage{232}

\subsection*{§~13. Die rationalen Transformationen der Functionen des Körpers $\Omega$.}

Ist $z_{1}$ eine beliebige, nicht constante, Function des Körpers $\Omega$ (eine \emph{Variable} in $\Omega$), so besteht, wie in §~2 nachgewiesen, zwischen $z_{1}$ und $z$ eine irreductible algebraische Gleichung, welche, von Nennern befreit, in Bezug auf $z_{1}$ vom Grade $e$, in Bezug auf $z$ vom Grade $e_{1}$ sei. Es ist, wie eben dort gezeigt, $e$ ein Divisor von $n$, $n = ef$. Es sei diese Gleichung
\[ G(\overset{e}{z_{1}}, \overset{e_{1}}{z}) = 0. \tag{1.} \]
Jede rationale Function $\zeta$ von $z$ und $z_{1}$ lässt sich (§~1) mit Hilfe dieser Gleichung auf die beiden Formen bringen
\[ \left\{ \begin{aligned} \zeta &= x_{0} + x_{1}z_{1} + \cdots + x_{e-1}z_{1}^{e-1}, \\ \zeta &= x_{0}^{(1)} + x_{1}^{(1)}z + \cdots + x_{e_{1}-1}^{(1)}z^{e_{1}-1} \end{aligned} \right. \tag{2.} \]
und zwar nur auf eine Weise so, dass $x_{0}, x_{1}, \ldots x_{e-1}$ rationale Functionen von $z$, $x_{0}^{(1)}, x_{1}^{(1)}, \ldots x_{e_{1}-1}^{(1)}$ rationale Functionen von $z_{1}$ sind.

Ist nun $\theta$ eine solche Function, dass $1, \theta, \theta^{2}, \ldots \theta^{n-1}$ eine Basis${}^{*)}$ von $\Omega$ (in Bezug auf $z$) bilden, so bilden nach §~2 die $n$ Functionen
\[ \left\{ \begin{matrix} 1, & z_{1}, & z_{1}^{2}, & \ldots & z_{1}^{e-1}, \\ \theta, & \theta z_{1}, & \theta z_{1}^{2}, & \ldots & \theta z_{1}^{e-1}, \\ \cdots & \cdots & \cdots & \cdots & \cdots \\ \theta^{f-1}, & \theta^{f-1}z_{1}, & \theta^{f-1}z_{1}^{2}, & \ldots & \theta^{f-1}z_{1}^{e-1} \end{matrix} \right. \tag{3.} \]
gleichfalls eine solche Basis, und daraus ergiebt sich nach (2.), dass zwischen den $e_{1}f = n_{1}$ Functionen
\[ \left\{ \begin{matrix} 1, & z, & z^{2}, & \ldots & z^{e_{1}-1}, \\ \theta, & \theta z, & \theta z^{2}, & \ldots & \theta z^{e_{1}-1} \\ \cdots & \cdots & \cdots & \cdots & \cdots \\ \theta^{f-1}, & \theta^{f-1}z, & \theta^{f-1}z^{2}, & \ldots & \theta^{f-1}z^{e_{1}-1}, \end{matrix} \right. \tag{4.} \]
die zur Abkürzung mit
\[ \eta_{1}^{(1)}, \quad \eta_{2}^{(1)}, \quad \ldots \quad \eta_{n_{1}}^{(1)} \]
bezeichnet sein mögen, eine Gleichung von der Form
\[ x_{1}^{(1)}\eta_{1}^{(1)} + x_{2}^{(1)}\eta_{2}^{(1)} + \cdots + x_{n_{1}}^{(1)}\eta_{n_{1}}^{(1)} = 0 \]

\textbf{*)} Man könnte statt der Basis $1, \theta, \ldots \theta^{n-1}$ auch eine beliebige andere Basis von $\Omega$ dieser Betrachtung zu Grunde legen. Es genügt aber für unsern Zweck, wenn wir gerade diese wählen.

\origpage{233}
nur dann besteht, wenn die rationalen Functionen $x_{1}^{(1)}, x_{2}^{(1)}, \ldots x_{n_{1}}^{(1)}$ von $z_{1}$ sämmtlich verschwinden. Daraus folgt nach (2.), dass jede Function $\eta$ in $\Omega$, und zwar nur auf eine einzige Art, darstellbar ist in der Form:
\[ \eta = x_{1}^{(1)}\eta_{1}^{(1)} + x_{2}^{(1)}\eta_{2}^{(1)} + \cdots + x_{n_{1}}^{(1)}\eta_{n_{1}}^{(1)}, \]
worin die $x^{(1)}$ rationale Functionen von $z_{1}$ sind.

Jede solche Function $\eta$ genügt einer algebraischen Gleichung vom Grade $n_{1}$, deren Coefficienten rational von $z_{1}$ abhängen, denn es ist
\[ \begin{aligned} \eta\eta_{1}^{(1)} &= x_{1,1}^{(1)}\eta_{1}^{(1)} + x_{1,2}^{(1)}\eta_{2}^{(1)} + \cdots + x_{1,n_{1}}^{(1)}\eta_{n_{1}}^{(1)}, \\ \eta\eta_{2}^{(1)} &= x_{2,1}^{(1)}\eta_{1}^{(1)} + x_{2,2}^{(1)}\eta_{2}^{(1)} + \cdots + x_{2,n_{1}}^{(1)}\eta_{n_{1}}^{(1)}, \\ &\cdots\cdots\cdots\cdots \\ \eta\eta_{n_{1}}^{(1)} &= x_{n_{1},1}^{(1)}\eta_{1}^{(1)} + x_{n_{1},2}^{(1)}\eta_{2}^{(1)} + \cdots + x_{n_{1},n_{1}}^{(1)}\eta_{n_{1}}^{(1)}, \end{aligned} \]
und mithin
\[ \begin{vmatrix} x_{1,1}^{(1)} - \eta, & x_{1,2}^{(1)}, & \ldots & x_{1,n_{1}}^{(1)} \\ x_{2,1}^{(1)}, & x_{2,2}^{(1)} - \eta, & \ldots & x_{2,n_{1}}^{(1)} \\ \cdots & \cdots & \cdots & \cdots \\ x_{n_{1},1}^{(1)}, & x_{n_{1},2}^{(1)}. & \ldots & x_{n_{1},n_{1}}^{(1)} - \eta \end{vmatrix} = 0. \]
\ednote{The entry $x_{n_{1},2}^{(1)}$ is printed with a period where every other entry carries a comma; probably a comma that printed incompletely.}

Es lässt sich nun zeigen, dass man eine Function $\eta = \theta_{1}$ so auswählen kann, dass $\theta_{1}$ nicht zugleich einer Gleichung niedrigeren Grades, deren Coefficienten rational von $z_{1}$ abhängen, genügt.

Wir stützen uns zum Beweis dieser Behauptung auf den folgenden Satz, dessen Beweis sich leicht durch den Schluss von $m-1$ auf $m$ ergiebt. Ist
\[ F(x_{1}, x_{2}, \ldots x_{m}) \]
eine ganze rationale Function von $x_{1}, x_{2}, \ldots x_{m}$, deren Coefficienten Functionen in $\Omega$ sind, die nicht alle verschwinden, so kann man für die $x_{1}, x_{2}, \ldots x_{m}$ solche constanten oder rational von $z_{1}$ abhängigen Grössen setzen, dass $F$ in eine nicht verschwindende Function in $\Omega$ übergeht. Ist also $F(x_{1}, x_{2}, \ldots x_{m})$ für \emph{alle} solche $x_{1}, x_{2}, \ldots x_{m}$ gleich Null, so folgt auch für beliebige constante oder rational von $z_{1}$ abhängige $dx_{1}, dx_{2}, \ldots dx_{m}$
\[ dF = F'(x_{1})\,dx_{1} + F'(x_{2})\,dx_{2} + \cdots + F'(x_{m})\,dx_{m} = 0. \]
Ist nun
\[ \theta_{1} = x_{1}\eta_{1}^{(1)} + x_{2}\eta_{2}^{(1)} + \cdots + x_{n_{1}}\eta_{n_{1}}^{(1)} \]
und
\[ \left\{ \begin{aligned} 1 &= x_{1,0}\eta_{1}^{(1)} + x_{2,0}\eta_{2}^{(1)} + \cdots + x_{n_{1},0}\eta_{n_{1}}^{(1)}, \\ \theta_{1} &= x_{1,1}\eta_{1}^{(1)} + x_{2,1}\eta_{2}^{(1)} + \cdots + x_{n_{1},1}\eta_{n_{1}}^{(1)}, \\ &\cdots\cdots\cdots\cdots \\ \theta_{1}^{m} &= x_{1,m}\eta_{1}^{(1)} + x_{2,m}\eta_{2}^{(1)} + \cdots + x_{n_{1},m}\eta_{n_{1}}^{(1)}, \end{aligned} \right. \tag{5.} \]

\origpage{234}
so sind die $x_{k,h}$ ganze rationale und homogene Functionen vom Grade $h$ von $x_{1}, x_{2}, \ldots x_{n_{1}}$ und hängen ausserdem rational von $z_{1}$ ab.

Ist also
\[ \varphi(\theta_{1}) = a_{m}\theta_{1}^{m} + a_{m-1}\theta_{1}^{m-1} + \cdots + a_{1}\theta_{1} + a_{0} = 0 \]
die Gleichung niedrigsten Grades, welcher $\theta_{1}$ genügt, deren Coefficienten rational von $z_{1}$ abhängen, so genügen die Functionen $a_{0}, a_{1}, \ldots a_{m}$ der Bedingung
\[ a_{0}x_{\iota,0} + a_{1}x_{\iota,1} + \cdots + a_{m}x_{\iota,m} = 0, \qquad {\scriptstyle (\iota = 1,\, 2,\, \ldots\, n_{1})} \]
und zugleich ist $m \leqq n_{1}$. Da nun nicht alle aus den Coefficienten $x_{h,k}$ zu bildenden $m$-reihigen Determinanten verschwinden können, weil sonst $\theta_{1}$ einer Gleichung von niedrigerem als dem $m^{\text{ten}}$ Grade genügen würde, so folgt aus letzteren Gleichungen, dass man die $a_{0}, a_{1}, \ldots a_{m}$ als ganze homogene Functionen von $x_{1}, x_{2}, \ldots x_{n_{1}}$ voraussetzen kann.

Wenn nun die Gleichung $\varphi(\theta_{1}) = 0$ für alle rational von $z_{1}$ abhängigen $x_{1}, x_{2}, \ldots x_{n_{1}}$ bestehen soll, so muss nach obigem Satze auch
\[ d\varphi = \varphi'(\theta_{1})\,d\theta_{1} + da_{m}\theta_{1}^{m} + \cdots + da_{1}\theta_{1} + da_{0} = 0 \]
sein, und wenn $m < n_{1}$ ist, so lassen sich die $dx_{1}, dx_{2}, \ldots dx_{n_{1}}$, ohne dass sie alle verschwinden, so bestimmen, dass
\[ da_{m} : da_{m-1} : \ldots : da_{1} : da_{0} = a_{m} : a_{m-1} : \ldots : a_{1} : a_{0} \]
und folglich
\[ \varphi'(\theta_{1})\,d\theta_{1} = 0 \]
ist. Da aber $\varphi'(\theta_{1})$ vom Grade $m-1$ ist, so muss $d\theta_{1} = 0$, also $dx_{1} = 0$, $dx_{2} = 0, \ldots dx_{n_{1}} = 0$ sein. Daher kann nur $m = n_{1}$ sein.

Ist also $\theta_{1}$ so bestimmt, dass die Gleichung niedrigsten Grades
\[ F_{1}(\overset{n_{1}}{\theta_{1}}, z_{1}) = 0 \]
den Grad $n_{1}$ wirklich erreicht, so lassen sich alle Functionen in $\Omega$, und zwar nur auf eine Weise in der Form darstellen
\[ \eta = x_{0}^{(1)} + x_{1}^{(1)}\theta_{1} + \cdots + x_{n_{1}-1}^{(1)}\theta_{1}^{n_{1}-1}, \]
worin die Coefficienten $x_{0}^{(1)}, x_{1}^{(1)}, \ldots x_{n_{1}-1}^{(1)}$ rational von $z_{1}$ abhängen; denn man kann unter dieser Voraussetzung $\eta_{1}^{(1)}, \eta_{2}^{(1)}, \ldots \eta_{n_{1}}^{(1)}$ vermittelst der Gleichungen (5.) in der angegebenen Weise darstellen.

\emph{Es lassen sich also sowohl $z_{1}$, $\theta_{1}$ rational durch $z$, $\theta$, als auch umgekehrt $z$, $\theta$ rational durch $z_{1}$, $\theta_{1}$ darstellen.}

Die Variable $z$, die wir bisher als die unabhängige bezeichnet haben, kann daher jede beliebige (nicht constante) Function des Körpers $\Omega$ sein.

\origpage{235}
Während aber die Gesammtheit aller Functionen des Körpers $\Omega$ gänzlich ungeändert bleibt, sind die Begriffe: \emph{Basis, Norm, Spur, Discriminante, ganze Function, Modul, Ideal}, wesentlich abhängig von der Wahl der \emph{unabhängigen} Veränderlichen $z$.

In dem besonderen Falle nur, wenn zwei Variable $z$, $z_{1}$ linear von einander abhängen, ist eine Basis von $\Omega$ in Bezug auf $z$ zugleich eine solche in Bezug auf $z_{1}$; ebenso sind Normen, Spuren und Discriminanten in diesem Fall für $z$ und $z_{1}$ identisch.

Sind $\alpha$, $\beta$ irgend zwei Functionen in $\Omega$, so bestehen zwischen denselben Gleichungen, deren linker Theil eine ganze rationale Function von $\alpha$ und $\beta$ ist.

Unter diesen ist eine (nach §~1)
\[ F(\alpha, \beta) = 0, \]
welche sowohl in Bezug auf $\alpha$ als in Bezug auf $\beta$ von möglichst niedrigem Grade ist, und diese soll die \emph{zwischen $\alpha$ und $\beta$ bestehende irreductible Gleichung} heissen. Diese ist, von einem constanten Factor abgesehen, völlig bestimmt.

\section*{II. Abtheilung.}

\subsection*{§~14. Die Punkte der Riemannschen Fläche.}

Die bisherigen Betrachtungen über die Functionen des Körpers $\Omega$ waren rein formaler Natur. Alle Resultate waren rationale, d.\ h.\ nach den Regeln der Buchstabenrechnung mittelst der vier Species abgeleitete Folgerungen aus der zwischen zwei Functionen in $\Omega$ bestehenden irreductibeln Gleichung. Die numerischen Werthe dieser Functionen kamen nirgends in Betracht. Man würde sogar, ohne andere Principien anzuwenden, die formelle Behandlung noch wesentlich weiter treiben können, indem man zwei Functionen des Körpers $\Omega$ nicht als durch eine Gleichung verbunden, sondern als unabhängige Veränderliche auffasst, wobei dann alles auf algebraische Theilbarkeit von rationalen Functionen zweier Veränderlichen hinausläuft. Wir haben auch diesen Weg durchgeführt, der jedoch in Darstellung und Ausdrucksweise sehr schwerfällig ist und bezüglich der Strenge

\origpage{236}
nichts mehr leistet als der im Vorhergehenden benutzte Gang. Nachdem nun aber der formale Theil der Untersuchung so weit geführt ist, drängt sich die Frage auf, \emph{in welchem Umfange es möglich ist, den Functionen in $\Omega$ solche bestimmten Zahlenwerthe beizulegen, dass alle zwischen diesen Functionen bestehenden rationalen Relationen (Identitäten) in richtige Zahlengleichungen übergehen.} Es erweist sich bei dieser Untersuchung als zweckmässig, auch das Unendlichgrosse als \emph{eine} bestimmte Zahl $\infty$ (Constante) zu betrachten, mit welcher nach bestimmten Regeln gerechnet wird${}^{*)}$. Die mittelst der rationalen Operationen in dem so erweiterten Zahlengebiet ausgeführten Rechnungen führen stets zu einem ganz bestimmten Zahlenresultat, wenn nicht im Verlaufe der Rechnung eines der Zeichen $\infty \pm \infty$, $0.\infty$, $\frac{0}{0}$, $\frac{\infty}{\infty}$ auftritt, Zeichen, welchen kein bestimmter Werth zukommt. Das Auftreten einer solchen Unbestimmtheit in einer Gleichung ist nicht als ein Widerspruch aufzufassen, da in diesem Fall die Gleichung gar keine bestimmte Behauptung mehr enthält, also von der Wahrheit oder Unwahrheit derselben auch keine Rede sein kann. Unter den Functionen des Körpers $\Omega$ finden sich ausser unendlich vielen Veränderlichen auch sämmtliche Constanten, d.\ h.\ Zahlen. Hiernach gelangt man durch die oben gestellte Forderung zu folgendem Begriff.

1. \emph{Definition.} Wenn \emph{alle} Individuen $\alpha, \beta, \gamma, \ldots$ des Körpers $\Omega$ durch \emph{bestimmte} Zahlwerthe $\alpha_{0}, \beta_{0}, \gamma_{0}, \ldots$ so ersetzt werden, dass
\[ \begin{aligned} &(\text{I.}) \quad \alpha_{0} = \alpha, \quad \text{falls } \alpha \text{ constant ist, und allgemein} \\ &(\text{II.}) \quad (\alpha + \beta)_{0} = \alpha_{0} + \beta_{0}, \qquad (\text{IV.}) \quad (\alpha\beta)_{0} = \alpha_{0}\beta_{0}, \\ &(\text{III.}) \quad (\alpha - \beta)_{0} = \alpha_{0} - \beta_{0}, \qquad (\text{V.}) \quad \left(\frac{\alpha}{\beta}\right)_{0} = \frac{\alpha_{0}}{\beta_{0}} \end{aligned} \]
wird, so soll einem solchen Zusammentreffen bestimmter Werthe ein \emph{Punkt} $\mathfrak{P}$ zugeordnet werden (den man sich zur Versinnlichung irgend wie im Raume gelegen vorstellen mag${}^{**)}$), und wir sagen, \emph{in} $\mathfrak{P}$ sei $\alpha = \alpha_{0}$, oder $\alpha$ habe \emph{in} $\mathfrak{P}$ den Werth $\alpha_{0}$. Zwei Punkte heissen stets und nur dann verschieden,

\textbf{*)} Das Unendliche als einen bestimmten Werth zu betrachten ist in der Functionentheorie vielfach üblich und nützlich. Es spricht sich dies bei \emph{Riemann} z.\ B.\ darin aus, dass er seine die algebraischen Functionen darstellenden Flächen als geschlossen betrachtet.

\textbf{**)} Eine geometrische Versinnlichung des „Punktes“ ist übrigens keineswegs nothwendig, und trägt zu einer leichteren Auffassung nicht einmal viel bei. Es genügt, das Wort „Punkt“ als einen kurzen und bequemen Ausdruck für die beschriebene Werth-Coexistenz zu betrachten.

\origpage{237}
wenn eine Function $\alpha$ in $\Omega$ existirt, die in beiden Punkten verschiedene Werthe hat.

Aus dieser Definition des Punktes soll nun die Existenz desselben, sowie der Umfang des Begriffes deducirt werden. Zunächst ist aber hervorzuheben, dass nach dieser Definition der „Punkt“ ein zum Körper $\Omega$ gehöriger \emph{invarianter} Begriff ist, der in keiner Weise abhängt von der Wahl der unabhängigen Veränderlichen, durch welche man die Functionen des Körpers darstellt.

2. \emph{Satz.} Ist ein Punkt $\mathfrak{P}$ gegeben, und $z$ eine in $\mathfrak{P}$ endliche Variable in $\Omega$ (eine solche existirt für jeden Punkt; denn ist $z_{0} = \infty$, so ist $\left(\frac{1}{z}\right)_{0} = 0$, also endlich), so hat auch jede ganze Function $\omega$ von $z$ in $\mathfrak{P}$ einen endlichen Werth $\omega_{0}$ --- denn zwischen $\omega$ und $z$ besteht eine Relation von der Form
\[ 1 = a\frac{1}{\omega} + b\frac{1}{\omega^{2}} + \cdots + k\frac{1}{\omega^{m}}, \]
worin $a, b, \ldots k$ als \emph{ganze rationale} Functionen von $z$ nach (II.), (III.), (IV.) in $\mathfrak{P}$ endliche Werthe haben. Mithin kann $\left(\frac{1}{\omega}\right)_{0}$ nicht gleich 0, also $\omega_{0}$ nicht gleich $\infty$ sein.

3. \emph{Satz.} Ist $z$ irgend eine in $\mathfrak{P}$ endliche Variable, so ist der Inbegriff $\mathfrak{p}$ aller derjenigen ganzen Functionen $\pi$ von $z$, welche in $\mathfrak{P}$ verschwinden, ein Primideal in $z$; wir sagen, der Punkt $\mathfrak{P}$ \emph{erzeuge} dies Primideal $\mathfrak{p}$. Ist $\omega$ eine ganze Function von $z$, welche in $\mathfrak{P}$ den Werth $\omega_{0}$ hat, so ist $\omega \equiv \omega_{0} \ (\text{mod.}\,\mathfrak{p})$.

\emph{Beweis.} Ist $\pi'_{0} = 0$, $\pi''_{0} = 0$, so ist auch $(\pi' + \pi'')_{0} = \pi'_{0} + \pi''_{0} = 0$, und wenn $\omega$ eine beliebige ganze Function von $z$, also $\omega_{0}$ endlich ist, so folgt aus $\pi_{0} = 0$ auch $(\omega\pi)_{0} = \omega_{0}\pi_{0} = 0$; also ist $\mathfrak{p}$ ein \emph{Ideal} in $z$ (§~7, I., II.). Das Ideal $\mathfrak{p}$ ist von $\mathfrak{o}$ verschieden, da es die Function „1“ nicht enthält.

Hat $\omega$ in $\mathfrak{P}$ den Werth $\omega_{0}$, so ist $(\omega - \omega_{0})_{0} = 0$, folglich $\omega \equiv \omega_{0} \ (\text{mod.}\,\mathfrak{p})$, also jede ganze Function von $z$ einer Constanten congruent nach dem Modul $\mathfrak{p}$. Daher ist (§~9, 7.) $\mathfrak{p}$ ein \emph{Primideal}.

4. \emph{Satz.} Dasselbe Primideal $\mathfrak{p}$ kann nicht durch zwei verschiedene Punkte erzeugt werden.

Denn zunächst ist der Werth einer jeden ganzen Function $\omega$ in einem das Ideal $\mathfrak{p}$ erzeugenden Punkt $\mathfrak{P}$ durch die Congruenz $\omega \equiv \omega_{0} \ (\text{mod.}\,\mathfrak{p})$ vollkommen bestimmt. Ist aber $\eta$ eine beliebige Function in $\Omega$, so lassen

\origpage{238}
sich nach §~12, 1. zwei ganze Functionen $\alpha$, $\beta$, die nicht beide durch $\mathfrak{p}$ theilbar sind, so bestimmen, dass
\[ \eta = \frac{\alpha}{\beta} \]
wird. Da nun die endlichen Werthe $\alpha_{0}$, $\beta_{0}$ nicht beide verschwinden, so folgt aus (V.)
\[ \eta_{0} = \frac{\alpha_{0}}{\beta_{0}}, \]
also ebenfalls durch $\mathfrak{p}$ vollkommen bestimmt.

Es ergiebt sich hieraus noch, dass zwei Punkte, in denen eine Variable $z$ endliche Werthe hat, dann und nur dann von einander verschieden sind, wenn eine \emph{ganze} Function von $z$ existirt, welche in beiden verschiedene Werthe hat.

5. \emph{Satz.} Ist $z$ irgend eine Variable in $\Omega$ und $\mathfrak{p}$ ein Primideal in $z$, so giebt es \emph{einen} (und nach 4. auch \emph{nur} einen) Punkt $\mathfrak{P}$, welcher dies Primideal erzeugt, und welcher der \emph{Nullpunkt} des Ideals $\mathfrak{p}$ genannt werden soll.

\emph{Beweis.} Es sei $\eta$ eine beliebige Function in $\Omega$, und $\varrho$ eine solche, deren Oberideal durch $\mathfrak{p}$ aber nicht durch $\mathfrak{p}^{2}$ theilbar ist. Es lassen sich dann nach §~12, 6. stets und nur auf eine Weise eine ganze Zahl $m$, eine von Null verschiedene endliche Constante $c$ und eine Function $\eta_{1}$, deren Unterideal nicht durch $\mathfrak{p}$ theilbar ist, so bestimmen, dass
\[ \eta = c\varrho^{m} + \eta_{1}\varrho^{m+1}. \]
Wir setzen
\[ \eta_{0} = 0, \quad c, \quad \infty, \]
je nachdem $m$ positiv, Null oder negativ ist. Dieser Werthbestimmung der Functionen des Körpers $\Omega$ entspricht ein Punkt $\mathfrak{P}$, da die Bedingungen (I.) bis (V.), wie man sofort übersieht, erfüllt sind${}^{*)}$.

\textbf{*)} Ist $\eta' = c'\varrho^{m'} + \eta'_{1}\varrho^{m'+1}$, so ist z.\ B.
\[ \frac{\eta}{\eta'} = \varrho^{m-m'}\left(\frac{c}{c'} + \varrho\eta''_{1}\right), \]
worin
\[ \eta''_{1} = \frac{c'\eta_{1} - c\eta'_{1}}{c'(c' + \varrho\eta'_{1})} \]
eine Function von derselben Beschaffenheit ist wie $\eta_{1}$ (noch einfacher ist der Beweis in den übrigen Fällen).

\origpage{239}

Jede Function, deren Oberideal durch $\mathfrak{p}$ theilbar ist, also ins Besondere jede Function in $\mathfrak{p}$ erhält nach dieser Festsetzung in $\mathfrak{P}$ den Werth Null, d.\ h.\ der so bestimmte Punkt $\mathfrak{P}$ erzeugt das Primideal $\mathfrak{p}$.

Jede Function, deren Unterideal durch $\mathfrak{p}$ theilbar ist, und nur eine solche hat in $\mathfrak{P}$ den Werth $\infty$, und daraus geht hervor, dass eine \emph{ganze} Function von $z$ in keinem Punkt, in welchem $z$ einen endlichen Werth hat, unendlich ist, und, da eine gebrochene Function von $z$ im Unterideal gewiss ein Primideal enthält, also mindestens in einem Punkt, in welchem $z$ endlich ist, unendlich sein muss, so ist auch umgekehrt jede Function, die in keinem Punkt, in welchem $z$ einen endlichen Werth hat, unendlich ist, eine ganze Function von $z$.

6. Aus 3., 4., 5. ergiebt sich nun das folgende Resultat. Um alle existirenden Punkte $\mathfrak{P}$ und jeden nur ein einziges mal zu erhalten, ergreife man eine beliebige Variable $z$ des Körpers $\Omega$; man bilde alle Primideale $\mathfrak{p}$ in $z$ und construire für jedes derselben den Nullpunkt, so sind alle diejenigen Punkte $\mathfrak{P}$ gefunden, in denen $z$ endlich bleibt; ist $\mathfrak{P}'$ ein von diesen verschiedener Punkt, so hat in ihm $z' = \frac{1}{z}$ den endlichen Werth Null; umgekehrt ist jeder Punkt $\mathfrak{P}'$, in dem $z'$ den Werth Null hat, von den Punkten $\mathfrak{P}$ verschieden. Das durch einen solchen Punkt $\mathfrak{P}'$ erzeugte Primideal $\mathfrak{p}'$ in $z'$ (welches aus allen in $\mathfrak{P}'$ verschwindenden ganzen Functionen von $z'$ besteht) geht in $z'$ auf, und umgekehrt ist der Nullpunkt eines jeden in $z'$ aufgehenden Primideals $\mathfrak{p}'$ in $z'$ ein Punkt $\mathfrak{P}'$, in welchem $z' = 0$ also $z = \infty$ ist. Mit diesen in endlicher Anzahl vorhandenen, den verschiedenen $\mathfrak{p}'$ entsprechenden Ergänzungspunkten und den vorher aus den Primidealen $\mathfrak{p}$ in $z$ abgeleiteten ist die Gesammtheit aller Punkte $\mathfrak{P}$ erschöpft, deren Inbegriff die \emph{Riemann}sche \emph{Fläche} $T$ bildet.

\subsection*{§~15. Die Ordnungszahlen.}

1. \emph{Definition.} Ist $\mathfrak{P}$ ein bestimmter Punkt, so betrachten wir die sämmtlichen in $\mathfrak{P}$ verschwindenden Functionen $\pi$ in $\Omega$, und ertheilen jeder derselben eine bestimmte \emph{Ordnungszahl} nach folgendem Gesichtspunkt.

Eine solche Function $\varrho$ hat die Ordnungszahl 1, oder heisst \emph{unendlich klein in der ersten Ordnung} oder $0^{1}$ in $\mathfrak{P}$, wenn \emph{alle} Quotienten $\frac{\pi}{\varrho}$ in $\mathfrak{P}$ endlich bleiben. Ist $\varrho'$ eine ebensolche Function wie $\varrho$, so ist $\frac{\varrho'}{\varrho}$ in $\mathfrak{P}$

\origpage{240}
weder 0 noch $\infty$, und umgekehrt, ist $\frac{\varrho'}{\varrho}$ in $\mathfrak{P}$ weder 0 noch $\infty$, so ist $\varrho'$ gleichfalls unendlich klein von der ersten Ordnung. Giebt es ferner für irgend eine Function $\pi$ einen ganzen positiven Exponenten $r$, so dass $\frac{\pi}{\varrho^{r}}$ in $\mathfrak{P}$ weder 0 noch $\infty$ wird, so gilt dasselbe von $\frac{\pi}{\varrho^{\prime r}}$, und $\pi$ erhält die Ordnungszahl $r$ oder heisst unendlich klein in der Ordnung $r$ im Punkte $\mathfrak{P}$. Wir werden auch sagen, $\pi$ ist $0^{r}$ in $\mathfrak{P}$ oder $\pi$ ist 0 in $\mathfrak{P}^{r}$.

Um die Frage nach der Existenz solcher Functionen $\varrho$ und solcher Ordnungszahlen $r$ zu entscheiden, ergreife man eine beliebige in $\mathfrak{P}$ endliche Variable $z$, bezeichne mit $\mathfrak{p}$ das durch $\mathfrak{P}$ erzeugte Primideal in $z$, und stelle jede Function $\pi$ (mit Ausnahme der ordnungslosen Constanten 0) nach §~12 als Quotienten von zwei relativen Primidealen in $z$ dar. Das Oberideal jeder dieser Functionen ist dann durch $\mathfrak{p}$ theilbar, und es giebt darunter auch solche, deren Oberideal nicht durch $\mathfrak{p}^{2}$ theilbar ist; diese haben die Ordnungszahl 1; für die übrigen Functionen $\pi$ ist die Ordnungszahl der Exponent der höchsten im Oberideal aufgehenden Potenz von $\mathfrak{p}$, was sich aus den Sätzen des §~12 ohne Weiteres ergiebt.

2. Hat eine Function $\eta$ den endlichen Werth $\eta_{0}$ in $\mathfrak{P}$, so sagen wir, $\eta$ habe diesen Werth $r$-mal in $\mathfrak{P}$ oder in $r$ mit $\mathfrak{P}$ zusammenfallenden Punkten, oder in $\mathfrak{P}^{r}$, wenn die Function $\eta - \eta_{0}$ in $\mathfrak{P}$ unendlich klein in der Ordnung $r$ ist. Ist aber $\eta_{0} = \infty$, so sagen wir, $\eta$ habe den Werth $\infty$ $r$-mal in $\mathfrak{P}$ oder in $r$ mit $\mathfrak{P}$ zusammenfallenden Punkten, oder $\eta$ sei $\infty^{r}$ in $\mathfrak{P}$ oder $\infty$ in $\mathfrak{P}^{r}$, wenn $\frac{1}{\eta}$ in $\mathfrak{P}^{r}$ verschwindet.

3. Wird eine Function $\eta$ in $\mathfrak{P}$ $\infty^{r}$, so legen wir derselben auch die Ordnungszahl $-r$ bei, wenn aber $\eta$ in $\mathfrak{P}$ weder 0 noch $\infty$ wird, so habe sie die Ordnungszahl 0. Hiernach kommt in einem beliebigen Punkt $\mathfrak{P}$ jeder Function des Körpers $\Omega$ eine ganz bestimmte \emph{Ordnungszahl} zu, mit Ausnahme der beiden Constanten 0 und $\infty$.

4. Ist $\varrho$ eine Function, welche in einem beliebigen Punkt $\mathfrak{P}$ die Ordnungszahl 1 besitzt, und $\eta$ eine Function mit der (positiven, negativen oder verschwindenden) Ordnungszahl $m$, so lässt sich nach dem Schlusssatz des §~12 für jedes beliebige positive $r$ eine Reihe von Constanten $c_{0}, c_{1}, \ldots c_{r-1}$, deren erste \emph{nicht} verschwindet, und eine in $\mathfrak{P}$ endliche Function $\sigma$ so bestimmen, dass
\[ \eta = c_{0}\varrho^{m} + c_{1}\varrho^{m+1} + \cdots + c_{r-1}\varrho^{m+r-1} + \sigma\varrho^{m+r} \]
wird.

\origpage{241}

5. Hieraus ergiebt sich unmittelbar, dass die Ordnungszahl eines Productes zweier oder mehrerer Functionen gleich ist der Summe der Ordnungszahlen der einzelnen Factoren.

Die Ordnungszahl eines Quotienten zweier Functionen ist gleich der Differenz der Ordnungszahlen des Zählers und Nenners.

Ist $\eta_{1}, \eta_{2}, \ldots \eta_{s}$ eine Reihe von Functionen und $m$ die algebraisch kleinste unter ihren Ordnungszahlen, so ist
\[ \begin{aligned} \eta_{1} &= e_{1}\varrho^{m} + \sigma_{1}\varrho^{m+1}, \\ \eta_{2} &= e_{2}\varrho^{m} + \sigma_{2}\varrho^{m+1}, \\ &\cdots\cdots\cdots\cdots \\ \eta_{s} &= e_{s}\varrho^{m} + \sigma_{s}\varrho^{m+1}, \end{aligned} \]
worin die Constanten $e_{1}, e_{2}, \ldots e_{s}$ jedenfalls nicht alle verschwinden. Sind daher $c_{1}, c_{2}, \ldots c_{s}$ Constanten, so ist die Ordnungszahl von
\[ \eta = c_{1}\eta_{1} + c_{2}\eta_{2} + \cdots + c_{s}\eta_{s}, \]
falls $c_{1}e_{1} + c_{2}e_{2} + \cdots + c_{s}e_{s}$ von Null verschieden ist, ebenfalls $m$, sonst grösser als $m$.

6. Complexe von Punkten, welche denselben Punkt auch mehrmals enthalten können, nennen wir \emph{Polygone} und bezeichnen dieselben mit $\mathfrak{A}, \mathfrak{B}, \mathfrak{C}, \ldots$

Es bedeute ferner $\mathfrak{A}\mathfrak{B}$ das aus den Punkten von $\mathfrak{A}$ und von $\mathfrak{B}$ zusammengesetzte Polygon, in der Weise, dass, wenn ein Punkt $\mathfrak{P}$ $r$-mal in $\mathfrak{A}$, $s$-mal in $\mathfrak{B}$ auftritt, er $(r+s)$-mal in $\mathfrak{A}\mathfrak{B}$ vorkommt. Daraus ergiebt sich die Bedeutung von $\mathfrak{P}^{r}$ und von $\mathfrak{A} = \mathfrak{P}^{r}\mathfrak{P}_{1}^{r_{1}}\mathfrak{P}_{2}^{r_{2}}\ldots$, und die Gesetze der Theilbarkeit der Polygone in vollkommener Uebereinstimmung mit denen der Theilbarkeit der ganzen Zahlen und der Ideale. Die Rolle der Primfactoren übernehmen dabei die Punkte; um aber auch die Einheit zu erhalten, muss man das gar keinen Punkt enthaltende Polygon $\mathfrak{O}$ (das Nulleck) zulassen.

Die Anzahl der Punkte eines Polygons heisst seine \emph{Ordnung}. Ein Polygon von der Ordnung $n$ wird auch kurz ein \emph{$n$-Eck} genannt.

Der \emph{grösste gemeinschaftliche Theiler} zweier Polygone $\mathfrak{A}, \mathfrak{B}$ ist dasjenige Polygon, welches jeden Punkt $\mathfrak{P}$ so oft enthält, als er in $\mathfrak{A}$ und $\mathfrak{B}$ \emph{mindestens} vorkommt. Ist dies $\mathfrak{O}$, so heissen $\mathfrak{A}, \mathfrak{B}$ \emph{relativ prim}.

Das \emph{kleinste gemeinschaftliche Vielfache} von $\mathfrak{A}$ und $\mathfrak{B}$ ist dasjenige Polygon, welches jeden Punkt so oft enthält, als er in $\mathfrak{A}$ und $\mathfrak{B}$ \emph{höchstens}

\origpage{242}
vorkommt. Sind $\mathfrak{A}, \mathfrak{B}$ relativ prim, so ist $\mathfrak{A}\mathfrak{B}$ ihr kleinstes gemeinschaftliches Vielfache.

Ist $\mathfrak{A} = \mathfrak{P}^{r}\mathfrak{P}_{1}^{r_{1}}\mathfrak{P}_{2}^{r_{2}}\ldots$ ein beliebiges Polygon, so giebt es stets Functionen $z$ in $\Omega$, welche in keinem der Punkte $\mathfrak{A}$ unendlich sind. Denn wenn $z$ in einigen Punkten von $\mathfrak{A}$ unendlich ist, so kann man eine Constante $c$ so wählen, dass $z-c$ in keinem der Punkte von $\mathfrak{A}$ den Werth 0 hat, und dann ist $\frac{1}{z-c}$ in allen Punkten des Polygons $\mathfrak{A}$ endlich. Legt man eine solche Variable $z$ zu Grunde, so ist der Inbegriff aller derjenigen ganzen Functionen von $z$, welche in den Punkten des Polygons $\mathfrak{A}$ (jeden nach seiner Vielfachheit gezählt) verschwinden, ein Ideal $\mathfrak{a} = \mathfrak{p}^{r}\mathfrak{p}_{1}^{r_{1}}\mathfrak{p}_{2}^{r_{2}}\ldots$, und man kann sagen, das Polygon $\mathfrak{A}$ \emph{erzeuge} das Ideal $\mathfrak{a}$, oder $\mathfrak{A}$ sei das \emph{Nullpolygon} des Ideals $\mathfrak{a}$. Der Idealbegriff fällt hiernach vollständig zusammen mit dem Begriff eines Systems ganzer Functionen, welche alle in denselben festen Punkten verschwinden. Das Ideal $\mathfrak{o}$ wird erzeugt durch das Nulleck $\mathfrak{O}$.

Das Product zweier oder mehrerer Ideale wird erzeugt durch das Product der Nullpolygone der Factoren, grösster gemeinschaftlicher Theiler und kleinstes gemeinschaftliches Vielfache zweier Ideale durch den grössten gemeinschaftlichen Theiler und das kleinste gemeinschaftliche Vielfache der entsprechenden Nullpolygone.

7. \emph{Satz.} Ist $z$ irgend eine Variable in $\Omega$ und $n$ der Grad des Körpers $\Omega$ in Bezug auf $z$, so nimmt $z$ jeden bestimmten Werth $c$ in genau $n$ Punkten an. --- Denn wenn $\mathfrak{o}$ das System aller ganzen Functionen von $z$ und $c$ eine endliche Constante bedeutet, so ist
\[ \mathfrak{o}(z-c) = \mathfrak{p}_{1}^{e_{1}}\mathfrak{p}_{2}^{e_{2}}\ldots, \quad e_{1} + e_{2} + \cdots = n \quad (\text{§~9, 7.}), \]
wenn $\mathfrak{p}_{1}, \mathfrak{p}_{2}, \ldots$ von einander verschiedene Primideale in $z$ bedeuten. Bezeichnet man mit $\mathfrak{P}_{1}, \mathfrak{P}_{2}, \ldots$ die Nullpunkte von $\mathfrak{p}_{1}, \mathfrak{p}_{2}, \ldots$, so hat nach 2. $z$ den Werth $c$ in $e_{1}$ Punkten $\mathfrak{P}_{1}$ (oder in $\mathfrak{P}_{1}^{e_{1}}$), in $e_{2}$ Punkten $\mathfrak{P}_{2}$ (oder in $\mathfrak{P}_{2}^{e_{2}}$) u.\ s.\ f., also in den $n$ Punkten des Polygons $\mathfrak{P}_{1}^{e_{1}}\mathfrak{P}_{2}^{e_{2}}\ldots$ Umgekehrt: ist $\mathfrak{P}$ ein Punkt, in welchem $z$ den Werth $c$ hat, und $\mathfrak{p}$ das durch $\mathfrak{P}$ erzeugte Primideal in $z$, so ist $z \equiv c \ (\text{mod.}\,\mathfrak{p})$, und folglich ist $\mathfrak{p}$ eines der Ideale $\mathfrak{p}_{1}, \mathfrak{p}_{2}, \ldots$, mithin $\mathfrak{P}$ einer der Punkte $\mathfrak{P}_{1}\mathfrak{P}_{2}\ldots$ Dasselbe Resultat gilt aber auch für $c = \infty$; denn weil $n$ auch der Grad von $\Omega$ in Bezug auf $\frac{1}{z}$ ist, so nimmt letztere Variable den Werth 0, folglich $z$ den Werth $\infty$ in genau $n$ Punkten an. Aus §~11 folgt, dass nur für eine endliche Anzahl

\origpage{243}
von Werthen der Constanten $c$ einer der Exponenten $e_{1}, e_{2}, \ldots$ grösser als 1 sein kann.

Die Zahl $n$, d.\ h.\ die Anzahl der Punkte, in welchen die Function $z$ je einen constanten Werth hat, soll die \emph{Ordnung} der Function $z$ genannt werden. Die Constanten und \emph{nur} diese haben die Ordnung Null. Für alle anderen Functionen in $\Omega$ ist die Ordnung eine positive ganze Zahl. Die Ordnung einer Variablen $z$ ist zugleich der Grad des Körpers $\Omega$ in Bezug auf $z$.

\subsection*{§~16. Conjugirte Punkte und conjugirte Werthe.}

1. \emph{Definition.} Ist $c$ ein bestimmter Zahlwerth, so entspricht demselben, wie in §~15 gezeigt, ein Polygon $\mathfrak{A}$ von $n$ (gleichen oder verschiedenen) Punkten $\mathfrak{P}', \mathfrak{P}'', \ldots \mathfrak{P}^{(n)}$, in welchen die Variable $n^{\text{ter}}$ Ordnung $z$ eben diesen Werth hat; diese $n$ Punkte sollen \emph{conjugirt nach $z$} heissen; durch einen von ihnen (und durch die Variable $z$) sind die übrigen bestimmt. Lässt man $c$ nach und nach alle Werthe annehmen, so bewegt sich das Polygon $\mathfrak{A} = \mathfrak{P}'\mathfrak{P}''\ldots\mathfrak{P}^{(n)}$ und zwar so, dass stets alle seine Punkte sich verändern. Man erhält hierbei also alle überhaupt existirenden Punkte und nur diejenigen (in endlicher Anzahl vorhandenen) mehrfach, in welchen $z - z_{0}$ oder $\frac{1}{z}$ in höherer als der ersten Ordnung verschwindet. Es ist daher das Product aller dieser Polygone
\[ \Pi\mathfrak{A} = T\mathfrak{Z}_{z}, \]
wo $T$ die einfache Gesammtheit aller Punkte, die \emph{Riemann}sche \emph{Fläche}, $\mathfrak{Z}_{z}$ ein bestimmtes endliches Polygon ist, welches das \emph{Verzweigungs-} oder \emph{Windungspolygon von $T$ in $z$} heisst. Jeder in $\mathfrak{Z}_{z}$ enthaltene Punkt $\mathfrak{Q}$ heisst ein \emph{Verzweigungs-} oder \emph{Windungspunkt von $T$ in $z$}, und zwar von der Ordnung $s$, wenn er genau $s$-mal in $\mathfrak{Z}_{z}$ vorkommt. Es ist $s = e-1$, wenn $z - z_{0}$ oder $\frac{1}{z}$ in $\mathfrak{Q}$ unendlich klein von der $e^{\text{ten}}$ Ordnung ist. Die Ordnung des Polygons $\mathfrak{Z}_{z}$ heisst die \emph{Verzweigungs-} oder \emph{Windungszahl} $\mathrm{w}_{z}$ der Fläche $T$ nach $z$. Diejenigen Punkte des Verzweigungspolygons, in welchen $z$ einen endlichen Werth hat, erzeugen zusammen das \emph{Verzweigungsideal} in $z$ (§~11).

Will man von dieser Definition der „absoluten“ \emph{Riemann}schen Fläche, welche ein zu dem Körper $\Omega$ gehöriger invarianter Begriff ist, zu der bekannten \emph{Riemann}schen Vorstellung übergehen, so hat man sich die Fläche

\origpage{244}
in einer $z$-Ebene ausgebreitet zu denken, welche sie dann überall mit Ausnahme der Verzweigungspunkte $n$-fach bedeckt.

2. \emph{Satz.} Ist
\[ z' = \frac{c+dz}{a+bz}, \]
worin $a, b, c, d$ Constanten bedeuten, deren Determinante $ad-bc$ von Null verschieden ist, so ist
\[ \mathfrak{Z}_{z} = \mathfrak{Z}_{z'}; \quad \mathrm{w}_{z} = \mathrm{w}_{z'}. \]
Denn wenn in einem Punkt $\mathfrak{P}$ $z - z_{0}$ oder $\frac{1}{z}$ unendlich klein in der $e^{\text{ten}}$ Ordnung ist, so ist in demselben Punkte auch
\[ z' - z'_{0} = \frac{(ad-bc)(z-z_{0})}{(a+bz)(a+bz_{0})}, \]
oder falls $z_{0}$ unendlich ist:
\[ z' - z'_{0} = \frac{-(ad-bc)}{b(a+bz)}, \]
oder falls $z'_{0} = \infty$, also $a+bz_{0} = 0$ ist:
\[ \frac{1}{z'} = \frac{a+bz}{c+dz} \]
unendlich klein in der $e^{\text{ten}}$ Ordnung.

Ist ins Besondere $z' = \frac{1}{z}$, so ist die Verzweigungszahl $\mathrm{w}_{z} = \mathrm{w}_{z'}$ gleich dem Grad der Discriminante $\Delta_{z}(\Omega)$ vermehrt um die Anzahl der verschwindenden Wurzeln von $\Delta_{z'}(\Omega) = 0$ (§~11).

3. \emph{Definition.} Die Werthe $\eta', \eta'', \ldots \eta^{(n)}$, welche eine beliebige Function $\eta$ in $\Omega$ in $n$ nach $z$ conjugirten Punkten $\mathfrak{P}', \mathfrak{P}'', \ldots \mathfrak{P}^{(n)}$ annimmt, heissen \emph{conjugirte Werthe von $\eta$ nach $z$}.

4. \emph{Satz.} Ist $N_{z}(\eta)$ die Norm einer beliebigen Function $\eta$ in Bezug auf $z$, so ist der Werth, welchen diese rationale Function von $z$ für $z = z_{0}$ besitzt, gleich dem Product $\eta'\eta''\ldots\eta^{(n)}$ der zu $z = z_{0}$ gehörigen conjugirten Werthe von $\eta$, wobei von dem Fall, dass dies Product unbestimmt wird, also einer dieser conjugirten Werthe 0, ein anderer $\infty$ ist, abzusehen ist. Beim Beweis dieses Satzes können wir annehmen, es sei $z_{0}$ endlich; denn ist $z_{0} = \infty$, so legen wir statt $z$ die Variable $z' = \frac{1}{z}$ zu Grunde, wobei die Norm ungeändert bleibt. Ferner können wir annehmen, die Werthe $\eta', \eta'', \ldots \eta^{(n)}$ seien alle endlich; denn ist einer von ihnen unendlich, so ist n.\ V.\ keiner derselben gleich 0, und wir betrachten statt $\eta$ die Function $\frac{1}{\eta}$.

\origpage{245}

Es sei nun unter diesen Voraussetzungen
\[ \mathfrak{o}(z-z_{0}) = \mathfrak{p}_{1}^{e_{1}}\mathfrak{p}_{2}^{e_{2}}\mathfrak{p}_{3}^{e_{3}}\ldots \]
und $\mathfrak{P}_{1}, \mathfrak{P}_{2}, \mathfrak{P}_{3}, \ldots$ die Nullpunkte der von einander verschiedenen Primideale $\mathfrak{p}_{1}, \mathfrak{p}_{2}, \mathfrak{p}_{3}, \ldots$ Wir construiren ein System ganzer Functionen $\lambda, \mu$ von $z$ nach folgender Regel: Es sei
\[ \begin{matrix} \lambda_{1} & \text{theilbar} & \text{durch} & \mathfrak{p}_{1}, & \text{nicht} & \text{durch} & \mathfrak{p}_{1}^{2}; \\ \lambda_{2} & \text{''} & \text{''} & \mathfrak{p}_{2}, & \text{''} & \text{''} & \mathfrak{p}_{2}^{2}; \\ \lambda_{3} & \text{''} & \text{''} & \mathfrak{p}_{3}, & \text{''} & \text{''} & \mathfrak{p}_{3}^{2}; \\ \cdots & \cdots & \cdots & \cdots & \cdots & \cdots & \cdots \end{matrix} \]
\[ \begin{matrix} \mu_{1} & \text{theilbar} & \text{durch} & \mathfrak{p}_{2}^{e_{2}}, \mathfrak{p}_{3}^{e_{3}}, \ldots, & \text{nicht} & \text{durch} & \mathfrak{p}_{1}, \mathfrak{p}_{2}^{e_{2}+1}, \mathfrak{p}_{3}^{e_{3}+1}, \ldots; \\ \mu_{2} & \text{''} & \text{''} & \mathfrak{p}_{1}^{e_{1}}, \mathfrak{p}_{3}^{e_{3}}, \ldots, & \text{''} & \text{''} & \mathfrak{p}_{2}, \mathfrak{p}_{1}^{e_{1}+1}, \mathfrak{p}_{3}^{e_{3}+1}, \ldots; \\ \mu_{3} & \text{''} & \text{''} & \mathfrak{p}_{1}^{e_{1}}, \mathfrak{p}_{2}^{e_{2}}, \ldots, & \text{''} & \text{''} & \mathfrak{p}_{3}, \mathfrak{p}_{1}^{e_{1}+1}, \mathfrak{p}_{2}^{e_{2}+1}, \ldots{}^{*)}. \\ \cdots & \cdots & \cdots & \cdots & \cdots & \cdots & \cdots \end{matrix} \]
Die $n$ Functionen
\[ \begin{matrix} \mu_{1}, & \mu_{1}\lambda_{1}, & \mu_{1}\lambda_{1}^{2}, & \ldots & \mu_{1}\lambda_{1}^{e_{1}-1}, \\ \mu_{2}, & \mu_{2}\lambda_{2}, & \mu_{2}\lambda_{2}^{2}, & \ldots & \mu_{2}\lambda_{2}^{e_{2}-1}, \\ \mu_{3}, & \mu_{3}\lambda_{3}, & \mu_{3}\lambda_{3}^{2}, & \ldots & \mu_{3}\lambda_{3}^{e_{3}-1}, \\ \cdots & \cdots & \cdots & \cdots & \cdots \end{matrix} \]
die wir mit $\eta_{1}, \eta_{2}, \ldots \eta_{n}$ bezeichnen, bilden dann eine Basis von $\Omega$; diese Behauptung ist in der nun zu beweisenden allgemeineren enthalten.

Wenn
\[ (z-z_{0})\zeta = x_{1}\eta_{1} + x_{2}\eta_{2} + \cdots + x_{n}\eta_{n} \]
mit ganzen rationalen Coefficienten $x_{1}, x_{2}, \ldots x_{n}$ ist, und $\zeta$ in den Punkten $\mathfrak{P}_{1}, \mathfrak{P}_{2}, \mathfrak{P}_{3}, \ldots$ endliche Werthe $\zeta', \zeta'', \zeta''', \ldots$ hat, so müssen die sämmtlichen Coefficienten $x_{1}, x_{2}, \ldots x_{n}$ durch $z-z_{0}$ theilbar sein. In der That ist z.\ B.\ im Punkte $\mathfrak{P}_{1}$ die linke Seite unendlich klein mindestens in der Ordnung $e_{1}$. Es muss also nach §~15, 5. auch
\[ x_{1}\eta_{1} + x_{2}\eta_{2} + \cdots + x_{e_{1}}\eta_{e_{1}} = \mu_{1}(x_{1} + x_{2}\lambda_{1} + \cdots + x_{e_{1}}\lambda_{1}^{e_{1}-1}) \]
in dieser Ordnung unendlich klein sein. Dies ist aber nur möglich, wenn $x_{1}, x_{2}, \ldots x_{e_{1}}$ in $\mathfrak{P}_{1}$ verschwinden, also durch $z-z_{0}$ theilbar sind, w.\ z.\ b.\ w.

\textbf{*)} Die Möglichkeit, solche Functionen zu bestimmen, ergiebt sich aus §~9, 3. Anmerkung, oder auch nach §~11, 2., wonach man z.\ B.\ setzen kann
\[ \lambda = \varrho - b, \quad \mu\lambda^{e} = \psi(\varrho). \]

\origpage{246}

Hiernach können wir setzen:
\[ \begin{aligned} \eta\mu_{1}\lambda_{1}^{r} &= \mu_{1}(x_{1}^{(0)} + x_{1}^{(1)}\lambda_{1} + \cdots + x_{1}^{(e_{1}-1)}\lambda_{1}^{e_{1}-1}) \\ &\quad + \mu_{2}(x_{2}^{(0)} + x_{2}^{(1)}\lambda_{2} + \cdots + x_{2}^{(e_{2}-1)}\lambda_{2}^{e_{2}-1}) \\ &\quad + \mu_{3}(x_{3}^{(0)} + x_{3}^{(1)}\lambda_{3} + \cdots + x_{3}^{(e_{3}-1)}\lambda_{3}^{e_{3}-1}) + \cdots, \end{aligned} \]
worin die $x_{1}^{(0)}, x_{1}^{(1)}, \ldots x_{2}^{(0)}, \ldots$ rationale Functionen von $z$ sind, die alle für $z = z_{0}$ endlich bleiben. In den Punkten $\mathfrak{P}_{2}, \mathfrak{P}_{3}, \ldots$ ist die linke Seite unendlich klein mindestens in der Ordnung $e_{2}, e_{3}, \ldots$ Dasselbe gilt für $\mathfrak{P}_{2}$ von $\mu_{1}, \mu_{3}, \ldots$ aber nicht von $\mu_{2}$, für $\mathfrak{P}_{3}$ von $\mu_{1}, \mu_{2}, \ldots$ aber nicht von $\mu_{3}, \ldots$; folglich ist für $z = z_{0}$
\[ \begin{gathered} x_{2}^{(0)} = 0, \quad x_{2}^{(1)} = 0, \quad \ldots \quad x_{2}^{(e_{2}-1)} = 0, \\ x_{3}^{(0)} = 0, \quad x_{3}^{(1)} = 0, \quad \ldots \quad x_{3}^{(e_{3}-1)} = 0, \\ \cdots\cdots\cdots\cdots \end{gathered} \]
In $\mathfrak{P}_{1}$ ist die linke Seite unendlich klein mindestens in der Ordnung $r$; daher wird, wenn $r < e_{1}$ ist, für $z = z_{0}$
\[ x_{1}^{(0)} = 0, \quad x_{1}^{(1)} = 0, \quad \ldots \quad x_{1}^{(r-1)} = 0, \quad x_{1}^{(r)} = \eta'. \]
Dieselbe Betrachtung lässt sich auf die Functionen $\eta\mu_{2}\lambda_{2}^{r}, \eta\mu_{3}\lambda_{3}^{r}, \ldots$ anwenden. Setzt man also
\[ \begin{aligned} \eta\eta_{1} &= x_{1,1}\eta_{1} + x_{1,2}\eta_{2} + \cdots + x_{1,n}\eta_{n}, \\ \eta\eta_{2} &= x_{2,1}\eta_{1} + x_{2,2}\eta_{2} + \cdots + x_{2,n}\eta_{n}, \\ &\cdots\cdots\cdots\cdots \\ \eta\eta_{n} &= x_{n,1}\eta_{1} + x_{n,2}\eta_{2} + \cdots + x_{n,n}\eta_{n}, \end{aligned} \]
so werden in der Determinante
\[ N(\eta) = \Sigma \pm x_{1,1}x_{2,2}\ldots x_{n,n} \]
sämmtliche links von der Diagonalreihe stehenden Glieder für $z = z_{0}$ verschwinden, während von den Diagonalgliedern $e_{1}$ gleich $\eta'$, $e_{2}$ gleich $\eta''$, $e_{3}$ gleich $\eta'''$, $\ldots$ werden. Es ist also für $z = z_{0}$
\[ N(\eta) = \eta^{\prime e_{1}}\eta^{\prime\prime e_{2}}\eta^{\prime\prime\prime e_{3}}\ldots \]
w.\ z.\ b.\ w.

5. Da nach der Definition der Spur (§~2)
\[ S(\eta) = x_{1,1} + x_{2,2} + \cdots + x_{n,n} \]
ist, so führen dieselben Betrachtungen zu dem Satze:

Es ist für $z = z_{0}$
\[ S(\eta) = e_{1}\eta' + e_{2}\eta'' + e_{3}\eta''' + \cdots, \]

\origpage{247}
welcher jedoch nur unter der Voraussetzung gilt, dass die Werthe $\eta', \eta'', \eta''', \ldots$ endlich sind.

Der Satz 4. ergiebt für ein beliebiges constantes (oder rational von $z$ abhängiges) $t$ für $z = z_{0}$
\[ N(t-\eta) = (t-\eta')^{e_{1}}(t-\eta'')^{e_{2}}(t-\eta''')^{e_{3}}\ldots, \]
und daraus durch Vergleichung der Coefficienten gleicher Potenzen von $t$ für jeden dieser Coefficienten einen Ausdruck durch die conjugirten Werthe (symmetrische Functionen).

6. Ist $\eta_{1}, \eta_{2}, \ldots \eta_{n}$ eine Basis von $\Omega$, so ergiebt sich durch Anwendung von 5. sofort der Werth der Discriminante dieses Systems für $z = z_{0}$
\[ \Delta_{z}(\eta_{1}, \eta_{2}, \ldots \eta_{n}) = (\Sigma \pm \eta'_{1}\eta''_{2}\ldots\eta_{n}^{(n)})^{2}, \]
wenn $\eta'_{\iota}, \eta''_{\iota}, \ldots \eta_{\iota}^{(n)}$ die sämmtlichen gleichen oder verschiedenen, aber als endlich vorausgesetzten, zu $z = z_{0}$ gehörigen conjugirten Werthe von $\eta_{\iota}$ bedeuten.

\subsection*{§~17. Darstellung der Functionen des Körpers $\Omega$ durch Polygonquotienten.}

Eine Function $\eta$ des Körpers $\Omega$ hat nur in einer endlichen Anzahl von Punkten eine von Null verschiedene Ordnungszahl; die Summe sämmtlicher Ordnungszahlen ist gleich 0, also die Summe der positiven gleich der Summe der negativen Ordnungszahlen, und zwar gleich der Ordnung der Function $\eta$ (§~15). Sind die Ordnungszahlen einer Function $\eta$ für jeden Punkt $\mathfrak{P}$ bekannt, so ist damit die Function $\eta$ bis auf einen constanten Factor bestimmt; denn hat $\eta'$ überall dieselbe Ordnungszahl wie $\eta$, so hat $\frac{\eta}{\eta'}$ (nach §~15, 5.) überall die Ordnungszahl Null und ist also (nach §~15, 7.) eine Constante.

Bilden wir also ein Polygon $\mathfrak{A}$, in welches wir jeden Punkt, in dem $\eta$ eine positive Ordnungszahl hat, so oft aufnehmen, als diese Ordnungszahl angiebt, und ein zweites Polygon $\mathfrak{B}$, in welches wir in entsprechender Weise die Punkte aufnehmen, in welchen $\eta$ eine negative Ordnungszahl hat, so sind die Polygone $\mathfrak{A}, \mathfrak{B}$ von gleicher Ordnung, und zwar von der Ordnung der Function $\eta$. Durch diese Polygone $\mathfrak{A}, \mathfrak{B}$ ist also die Function $\eta$ bis auf einen constanten Factor bestimmt. Wir setzen in symbolischer Bezeichnung

\origpage{248}
\[ \eta = \frac{\mathfrak{A}}{\mathfrak{B}}, \]
und nennen $\mathfrak{A}$ das \emph{Obereck,} $\mathfrak{B}$ das \emph{Untereck} der Function $\eta{}^{*)}$.

Nach dieser Festsetzung sind die beiden Polygone $\mathfrak{A}, \mathfrak{B}$ relativ prim; es ist aber zweckmässig, die Bezeichnung dahin auszudehnen, dass man auch gemeinschaftliche Factoren in $\mathfrak{A}, \mathfrak{B}$ zulässt, was durch die Bestimmung geschieht, dass
\[ \frac{\mathfrak{M}\mathfrak{A}}{\mathfrak{M}\mathfrak{B}} = \frac{\mathfrak{A}}{\mathfrak{B}} \]
sein soll, wenn $\mathfrak{M}$ ein beliebiges Polygon bedeutet. Setzen wir nach dieser verallgemeinerten Bezeichnung
\[ \eta = \frac{\mathfrak{A}}{\mathfrak{B}}, \]
so kann ein Punkt $\mathfrak{P}$, in welchem $\eta$ die Ordnungszahl $m$ besitzt, $m_{1}$-mal in $\mathfrak{A}$, $m_{2}$-mal in $\mathfrak{B}$ aufgenommen werden, wenn $m_{1} - m_{2} = m$ ist. Es ist auch jetzt noch die Ordnung von $\mathfrak{A}$ gleich der von $\mathfrak{B}$, aber nicht mehr gleich der Ordnung der Function $\eta$.

Aus dieser Definition ergiebt sich (nach §~15, 5.) unmittelbar der Satz: Ist
\[ \eta = \frac{\mathfrak{A}}{\mathfrak{B}}, \quad \eta' = \frac{\mathfrak{A}'}{\mathfrak{B}'}, \]
so ist
\[ \eta\eta' = \frac{\mathfrak{A}\mathfrak{A}'}{\mathfrak{B}\mathfrak{B}'}, \quad \frac{\eta}{\eta'} = \frac{\mathfrak{A}\mathfrak{B}'}{\mathfrak{B}\mathfrak{A}'}. \]
Nach §~14, 5. \emph{ist eine Function $\eta'$ dann und nur dann eine ganze Function von $\eta$, wenn jeder im Untereck von $\eta'$ aufgehende Punkt auch in dem von $\eta$ enthalten ist.}

\subsection*{§~18. Aequivalente Polygone und Polygonclassen.}

1. \emph{Definition.} Zwei Polygone $\mathfrak{A}, \mathfrak{A}'$ von gleichviel Punkten heissen \emph{äquivalent,} wenn eine Function $\eta$ in $\Omega$ existirt, welche (nach §~17) die Bezeichnung hat:

\textbf{*)} Alle Functionen der einfachen Schaar $(\eta)$ haben hiernach dieselbe Bezeichnung $\frac{\mathfrak{A}}{\mathfrak{B}}$, und es würde daher correcter sein, $(\eta) = \frac{\mathfrak{A}}{\mathfrak{B}}$ zu setzen; indessen führt diese Bezeichnung zu unnöthigen Weitläufigkeiten.

\origpage{249}
\[ \eta = \frac{\mathfrak{A}}{\mathfrak{A}'}. \]

2. \emph{Satz.} Ist $\mathfrak{A}$ äquivalent mit $\mathfrak{A}'$ und mit $\mathfrak{A}''$, so ist auch $\mathfrak{A}'$ mit $\mathfrak{A}''$ äquivalent; denn aus
\[ \eta' = \frac{\mathfrak{A}'}{\mathfrak{A}}; \quad \eta'' = \frac{\mathfrak{A}''}{\mathfrak{A}} \]
folgt:
\[ \frac{\eta'}{\eta''} = \frac{\mathfrak{A}'}{\mathfrak{A}''}. \]

3. \emph{Definition und Satz.} Alle mit einem gegebenen Polygon $\mathfrak{A}$ äquivalenten Polygone $\mathfrak{A}', \mathfrak{A}'', \ldots$ bilden eine \emph{Polygonclasse $A$}. Nach 2. kommt dann jedes beliebige Polygon in einer und \emph{nur} in einer Classe vor; denn sind $\mathfrak{A}, \mathfrak{B}$ zwei äquivalente Polygone, welche zu den Classen $A, B$ führen, so ist nach 2. jedes Polygon der Classe $B$ zugleich in $A$ enthalten und umgekehrt, und daher sind beide Classen identisch.

Alle Polygone einer Classe haben dieselbe Ordnung, welche die \emph{Ordnung der Classe} genannt werden soll.

4. Es können aber Polygone existiren, welche mit keinem andern äquivalent sind, und deren jedes daher für sich eine Classe bildet. Solche Polygone mögen \emph{isolirte} genannt sein.

5. Ist $\mathfrak{M}$ ein beliebiges Polygon, und $\mathfrak{A}$ äquivalent mit $\mathfrak{A}'$, so ist auch $\mathfrak{M}\mathfrak{A}$ äquivalent mit $\mathfrak{M}\mathfrak{A}'$; aber auch umgekehrt folgt aus der Aequivalenz von $\mathfrak{M}\mathfrak{A}$ mit $\mathfrak{M}\mathfrak{A}'$ die Aequivalenz von $\mathfrak{A}$ mit $\mathfrak{A}'$.

6. Ist $\mathfrak{A}$ mit $\mathfrak{A}'$, $\mathfrak{B}$ mit $\mathfrak{B}'$ äquivalent, so ist auch $\mathfrak{A}\mathfrak{B}$ mit $\mathfrak{A}'\mathfrak{B}'$ äquivalent. Die Classe $C$, welcher das Product $\mathfrak{A}\mathfrak{B}$ angehört, umfasst daher die sämmtlichen Producte je zweier Polygone der Classen $A, B$ von $\mathfrak{A}$ und $\mathfrak{B}$ (aber ausserdem unter Umständen noch unendlich viele andere Polygone) und soll als das \emph{Product} der beiden Classen $A, B$ bezeichnet sein:
\[ C = AB = BA. \]
Die Definition des Products von mehreren Classen und die Gültigkeit des Fundamentalsatzes der Multiplication ergiebt sich hieraus von selbst.

7. Sind $A, B, D$ drei Classen, welche der Bedingung
\[ DA = DB \]
genügen, so folgt $A = B$; denn sind $\mathfrak{A}, \mathfrak{B}, \mathfrak{D}$ drei Polygone der Classen $A, B, D$, so folgt aus der Voraussetzung, dass $\mathfrak{D}\mathfrak{B}$ mit $\mathfrak{D}\mathfrak{A}$, und folglich $\mathfrak{B}$ mit $\mathfrak{A}$ äquivalent ist.

\origpage{250}

8. Geht ein Polygon $\mathfrak{A}$ der Classe $A$ in einem Polygon $\mathfrak{C}$ der Classe $C$ auf, so gilt dasselbe von jedem Polygon $\mathfrak{A}'$ der Classe $A$; denn aus $\mathfrak{C} = \mathfrak{A}\mathfrak{B}$ folgt nach 5., dass $\mathfrak{C}' = \mathfrak{A}'\mathfrak{B}$ in $C$ enthalten ist, und wir können also, obschon nicht umgekehrt jedes Polygon der Classe $C$ durch ein Polygon der Classe $A$ theilbar zu sein braucht, sagen, \emph{die Classe $C$ sei durch die Classe $A$ theilbar}. Ist $\mathfrak{B}'$ irgend ein Polygon der Classe $B$ von $\mathfrak{B}$, so ist auch $\mathfrak{C}'' = \mathfrak{A}'\mathfrak{B}'$ in $C$ enthalten und folglich
\[ C = AB. \]
Ist also $C$ durch $A$ theilbar, so giebt es eine und (nach 7.) \emph{nur} eine Classe $B$, welche der Bedingung
\[ C = AB \]
genügt.

\subsection*{§~19. Die Polygonschaaren.}

1. Sind $\mathfrak{A}_{1}, \mathfrak{A}_{2}, \ldots \mathfrak{A}_{s}$ bestimmte, und zwar äquivalente Polygone, und $\mathfrak{A}$ ein beliebiges Polygon derselben Classe $A$, so existiren $s$ Functionen in $\Omega$
\[ \eta_{1} = \frac{\mathfrak{A}_{1}}{\mathfrak{A}}, \quad \eta_{2} = \frac{\mathfrak{A}_{2}}{\mathfrak{A}}, \quad \ldots \quad \eta_{s} = \frac{\mathfrak{A}_{s}}{\mathfrak{A}}. \]
Setzt man, wie in §~15, 5. für einen beliebigen Punkt $\mathfrak{P}$
\[ \begin{aligned} \eta_{1} &= e_{1}\varrho^{m} + \sigma_{1}\varrho^{m+1}, \\ \eta_{2} &= e_{2}\varrho^{m} + \sigma_{2}\varrho^{m+1}, \\ &\cdots\cdots\cdots\cdots \\ \eta_{s} &= e_{s}\varrho^{m} + \sigma_{s}\varrho^{m+1}, \end{aligned} \]
worin $\varrho$ in $\mathfrak{P}$ $0^{1}$ ist, $e_{1}, e_{2}, \ldots e_{s}$ Constanten, die nicht alle verschwinden, und $\sigma_{1}, \sigma_{2}, \ldots \sigma_{s}$ in $\mathfrak{P}$ endliche Functionen bedeuten, so folgt, dass jede Function $\eta$ der Schaar $(\eta_{1}, \eta_{2}, \ldots \eta_{s})$, d.\ h.\ jede Function von der Form
\[ \eta = c_{1}\eta_{1} + c_{2}\eta_{2} + \cdots + c_{s}\eta_{s} \]
in $\mathfrak{P}$ eine Ordnungszahl hat, die nicht kleiner als $m$ ist, und daraus nach §~17, dass die Function $\eta$ in die Form
\[ \eta = \frac{\mathfrak{A}'}{\mathfrak{A}} \]
gesetzt werden kann, wo $\mathfrak{A}'$ gleichfalls in der Classe $A$ enthalten ist.

Wählt man für $\mathfrak{A}$ ein beliebiges anderes Polygon $\mathfrak{B}$ der Classe $A$, und setzt

\origpage{251}
\[ \begin{gathered} \zeta = \frac{\mathfrak{A}}{\mathfrak{B}}, \\ \eta_{1}\zeta = \eta'_{1} = \frac{\mathfrak{A}_{1}}{\mathfrak{B}}, \quad \eta_{2}\zeta = \eta'_{2} = \frac{\mathfrak{A}_{2}}{\mathfrak{B}}, \quad \ldots \quad \eta_{s}\zeta = \eta'_{s} = \frac{\mathfrak{A}_{s}}{\mathfrak{B}}, \end{gathered} \]
so wird auch
\[ \eta\zeta = \eta' = c_{1}\eta'_{1} + c_{2}\eta'_{2} + \cdots + c_{s}\eta'_{s}, \]
und folglich
\[ \eta' = \frac{\mathfrak{A}'}{\mathfrak{B}}. \]
Jedes durch den Nenner $\mathfrak{A}$ und ein Constantensystem $c_{1}, c_{2}, \ldots c_{s}$ erzeugte Polygon wird daher auch durch jeden andern derselben Classe angehörigen Nenner $\mathfrak{B}$ erzeugt, und der Inbegriff der sämmtlichen Polygone $\mathfrak{A}'$, die den verschiedenen Werthen der Constanten $c_{1}, c_{2}, \ldots c_{s}$ entsprechen, ist nur abhängig von den Polygonen $\mathfrak{A}_{1}, \mathfrak{A}_{2}, \ldots \mathfrak{A}_{s}$. Dieser Inbegriff soll daher eine \emph{Polygonschaar} mit der Basis $\mathfrak{A}_{1}, \mathfrak{A}_{2}, \ldots \mathfrak{A}_{s}$ genannt und mit
\[ (\mathfrak{A}_{1}, \mathfrak{A}_{2}, \ldots \mathfrak{A}_{s}) \]
bezeichnet werden.

2. Haben die Polygone $\mathfrak{A}_{1}, \mathfrak{A}_{2}, \ldots \mathfrak{A}_{s}$ einen grössten gemeinschaftlichen Theiler $\mathfrak{M}$, so ist derselbe nach 1. auch Theiler eines jeden Polygons $\mathfrak{A}'$ der Schaar $(\mathfrak{A}_{1}, \mathfrak{A}_{2}, \ldots \mathfrak{A}_{s})$, und kann der \emph{Theiler der Schaar} genannt werden; aber es lässt sich in dieser Schaar ein Polygon $\mathfrak{A}' = \mathfrak{M}\mathfrak{B}$ der Art bestimmen, dass $\mathfrak{B}$ relativ prim zu einem beliebig gegebenen Polygon wird. Ist nämlich unter Beibehaltung der Bezeichnung von 1. ein Punkt $\mathfrak{P}$ genau $\mu$-mal in $\mathfrak{M}$ und $\nu$-mal in $\mathfrak{A}$ enthalten, so ist, wenn
\[ \eta = e\varrho^{m} + \sigma\varrho^{m+1} \]
gesetzt wird, $m$ niemals kleiner als $\mu - \nu$, und es ist $m = \mu - \nu$, wenn man die Constanten $c_{1}, c_{2}, \ldots c_{s}$ so wählt, dass
\[ e = c_{1}e_{1} + c_{2}e_{2} + \cdots + c_{s}e_{s} \]
von Null verschieden ist. Der Punkt $\mathfrak{P}$ ist daher mindestens $\mu$-mal in $\mathfrak{A}'$ enthalten, und unter der letzteren Voraussetzung auch nicht öfter als $\mu$-mal. Da man nun die Constanten $c_{1}, c_{2}, \ldots c_{s}$ immer so wählen kann, dass eine beliebige Anzahl von Ausdrücken der Form
\[ \Sigma c_{\iota}e_{\iota}, \quad \Sigma c_{\iota}e'_{\iota}, \quad \ldots, \]
in deren keinem die sämmtlichen Constanten $e_{\iota}, e'_{\iota}, \ldots$ verschwinden, von Null verschiedene Werthe haben, so folgt die Richtigkeit der aufgestellten Behauptung.

\origpage{252}

3. Sind die Functionen $\eta_{1}, \eta_{2}, \ldots \eta_{s}$ in 1. linear abhängig oder unabhängig, so gilt das Gleiche von den Functionen $\eta'_{1}, \eta'_{2}, \ldots \eta'_{s}$. Wir werden dem entsprechend auch die Polygone $\mathfrak{A}_{1}, \mathfrak{A}_{2}, \ldots \mathfrak{A}_{s}$ \emph{linear abhängig} oder \emph{unabhängig} und ihr System \emph{linear reductibel} oder \emph{irreductibel nennen.}

Da nach §~5, 4. jede Functionenschaar eine irreductible Basis besitzt, so folgt, dass auch jede Polygonschaar eine \emph{irreductible Basis} hat. Ist $s$ die Anzahl der Polygone einer solchen Basis, so heisst die Schaar \emph{eine $s$-fache,} oder $s$ die \emph{Dimension der Schaar}. Irgend $s$ Polygone einer solchen Schaar bilden eine irreductible Basis derselben oder nicht, je nachdem sie linear unabhängig oder abhängig sind. (Vgl.\ §~5, 4.).

4. Sind die Polygone $\mathfrak{A}_{1}, \mathfrak{A}_{2}, \ldots \mathfrak{A}_{s}$ linear abhängig oder unabhängig, so sind, wenn $\mathfrak{M}$ ein beliebiges Polygon bedeutet, auch $\mathfrak{M}\mathfrak{A}_{1}, \mathfrak{M}\mathfrak{A}_{2}, \ldots \mathfrak{M}\mathfrak{A}_{s}$ linear abhängig oder unabhängig und umgekehrt.

\subsection*{§~20. Erniedrigung der Dimension der Schaar durch Theilbarkeitsbedingungen.}

1. Es sei
\[ S = (\mathfrak{A}_{1}, \mathfrak{A}_{2}, \ldots \mathfrak{A}_{s}) \]
eine $s$-fache Schaar vom Theiler $\mathfrak{M}$. Es wird nach der Mannigfaltigkeit derjenigen Polygone $\mathfrak{A}'$ der Schaar $S$ gefragt, welche einen beliebig gegebenen Punkt wenigstens einmal öfter enthalten als der Theiler $\mathfrak{M}$ der Schaar.

Ist der Punkt $\mathfrak{P}$ $\mu$-mal in $\mathfrak{M}$ und $\nu$-mal in einem beliebigen mit $\mathfrak{A}_{1}, \mathfrak{A}_{2}, \ldots$ äquivalenten Polygon $\mathfrak{A}$ enthalten, so ist, wenn wir wie in §~19
\[ \begin{aligned} \frac{\mathfrak{A}_{1}}{\mathfrak{A}} &= \eta_{1} = e_{1}\varrho^{m} + \sigma_{1}\varrho^{m+1}, \\ \frac{\mathfrak{A}_{2}}{\mathfrak{A}} &= \eta_{2} = e_{2}\varrho^{m} + \sigma_{2}\varrho^{m+1}, \\ &\cdots\cdots\cdots\cdots \\ \frac{\mathfrak{A}_{s}}{\mathfrak{A}} &= \eta_{s} = e_{s}\varrho^{m} + \sigma_{s}\varrho^{m+1} \end{aligned} \]
setzen, $m = \mu - \nu$, und von den Constanten $e_{1}, e_{2}, \ldots e_{s}$ ist wenigstens eine, etwa $e_{s}$, von Null verschieden. Die gesuchten Polygone $\mathfrak{A}'$ sind dann durch die Gleichung charakterisirt
\[ \frac{\mathfrak{A}'}{\mathfrak{A}} = \eta' = c_{1}\eta_{1} + c_{2}\eta_{2} + \cdots + c_{s}\eta_{s}, \]
worin die Constanten $c_{1}, c_{2}, \ldots c_{s}$ an die Bedingung gebunden sind
\[ c_{1}e_{1} + c_{2}e_{2} + \cdots + c_{s}e_{s} = 0. \]

\origpage{253}
Hiernach können wir setzen
\[ \frac{\mathfrak{A}'}{\mathfrak{A}} = e_{s}\eta' = c_{1}(e_{s}\eta_{1} - e_{1}\eta_{s}) + \cdots + c_{s-1}(e_{s}\eta_{s-1} - e_{s-1}\eta_{s}). \]
Daraus aber ergiebt sich, wenn wir
\[ \begin{aligned} \eta'_{1} &= e_{s}\eta_{1} - e_{1}\eta_{s}, \\ \eta'_{2} &= e_{s}\eta_{2} - e_{2}\eta_{s}, \\ &\cdots\cdots\cdots\cdots \\ \eta'_{s-1} &= e_{s}\eta_{s-1} - e_{s-1}\eta_{s} \end{aligned} \]
setzen, dass die Functionen $\eta'$ eine $(s-1)$-fache Schaar $(\eta'_{1}, \eta'_{2}, \ldots \eta'_{s-1})$ bilden; denn die Functionen $\eta'_{1}, \eta'_{2}, \ldots \eta'_{s-1}$ sind linear unabhängig, wenn es, wie vorausgesetzt, die Functionen $\eta_{1}, \eta_{2}, \ldots \eta_{s}$ sind. Es bilden also auch die Polygone $\mathfrak{A}'$ eine $(s-1)$-fache Schaar
\[ S' = (\mathfrak{A}'_{1}, \mathfrak{A}'_{2}, \ldots \mathfrak{A}'_{s-1}), \]
wenn
\[ \frac{\mathfrak{A}'_{\iota}}{\mathfrak{A}} = e_{s}\eta_{\iota} - e_{\iota}\eta_{s} \]
gesetzt wird. Der Theiler dieser Schaar ist durch $\mathfrak{M}\mathfrak{P}$ theilbar, wenn auch nicht nothwendig damit identisch.

2. Hieraus ergiebt sich sofort, dass die Polygone einer Schaar $S$, welche durch ein beliebiges $r$-Eck $\mathfrak{R}$ theilbar sind, eine \emph{mindestens} $(s-r)$-fache Schaar bilden. Denn nehmen wir an, es sei dies bereits für ein $r$-Eck $\mathfrak{R}$ bewiesen, so folgt die Richtigkeit der Behauptung für ein $(r+1)$-Eck $\mathfrak{P}\mathfrak{R}$ unmittelbar aus 1, indem durch das Hinzutreten des Punktes $\mathfrak{P}$, wenn $\mathfrak{P}$ im Theiler der durch $\mathfrak{R}$ bereits reducirten Schaar enthalten ist, die Dimension nicht weiter geändert, sonst um 1 erniedrigt wird.

Hieraus folgt als Specialfall, dass es in einer $s$-fachen Schaar immer \emph{wenigstens ein} Polygon giebt, welches durch ein gegebenes $(s-1)$-Eck theilbar ist.

3. Man kann, wenn $r \leqq s$ ist, das $r$-Eck $\mathfrak{R}$ so wählen, dass die durch $\mathfrak{R}$ theilbaren Polygone der Schaar $S$ eine genau $(s-r)$-fache Schaar bilden. Zu diesem Ende wähle man einen Punkt $\mathfrak{P}$, welcher im Theiler von $S$ nicht enthalten ist; die durch $\mathfrak{P}$ theilbaren Polygone in $S$ bilden nach 1. eine $(s-1)$-fache Schaar $S'$; man wähle einen zweiten Punkt $\mathfrak{P}'$, der nicht im Theiler von $S'$ enthalten ist; die Polygone in $S'$, die durch $\mathfrak{P}'$, d.\ h.\ die Polygone in $S$, die durch $\mathfrak{P}\mathfrak{P}'$ theilbar sind, bilden eine $(s-2)$-fache

\origpage{254}
Schaar, u.\ s.\ f.; zugleich erhellt aus dieser Bildungsweise, dass man $\mathfrak{R}$ zu einem beliebig gegebenen Polygon relativ prim annehmen kann. Ist $r = s$, so wird hiernach in $S$ kein durch $\mathfrak{R}$ theilbares Polygon existiren.

\subsection*{§~21. Die Dimensionen der Polygonclassen.}

1. \emph{Die Polygone einer Classe bilden eine Schaar von endlicher Dimension, welche die Dimension der Classe heissen soll.}

\emph{Beweis.} Wählt man in einer Classe $A$, deren Ordnung $m$ sei, irgend $s$ Polygone $\mathfrak{A}_{1}, \mathfrak{A}_{2}, \ldots \mathfrak{A}_{s}$ aus, so gehören alle Polygone der Schaar $(\mathfrak{A}_{1}, \mathfrak{A}_{2}, \ldots \mathfrak{A}_{s})$ zugleich in die Classe $A$. Die Anzahl der linear unabhängigen Polygone, die in $A$ enthalten sind, kann daher gewiss nicht grösser sein als $m+1$, weil man sonst (nach §~20, 2.) in der Classe ein durch ein beliebiges $(m+1)$-Eck theilbares Polygon finden könnte, was widersinnig ist. Wenn daher $s$ die Maximalzahl der linear unabhängigen Polygone $\mathfrak{A}_{1}, \mathfrak{A}_{2}, \ldots \mathfrak{A}_{s}$ der Classe $A$ ist, so muss jedes Polygon dieser Classe in der Schaar $(\mathfrak{A}_{1}, \mathfrak{A}_{2}, \ldots \mathfrak{A}_{s})$ enthalten sein, und $s$ ist die \emph{Dimension der Classe.} Das System der Polygone $\mathfrak{A}_{1}, \mathfrak{A}_{2}, \ldots \mathfrak{A}_{s}$ soll eine \emph{Basis der Classe} genannt werden.

Die isolirten Polygone bilden Classen von der Dimension 1.

2. Giebt es in einer Classe $C$ $s$ und nicht mehr linear unabhängige, durch ein gegebenes Polygon $\mathfrak{A}$ der Classe $A$ theilbare Polygone
\[ \mathfrak{C}_{1} = \mathfrak{A}\mathfrak{B}_{1}, \quad \mathfrak{C}_{2} = \mathfrak{A}\mathfrak{B}_{2}, \quad \ldots \quad \mathfrak{C}_{s} = \mathfrak{A}\mathfrak{B}_{s}, \]
so ist $C$ durch $A$ theilbar, und es existiren in $C$ auch ebenso viele linear unabhängige Polygone
\[ \mathfrak{C}'_{1} = \mathfrak{A}'\mathfrak{B}_{1}, \quad \mathfrak{C}'_{2} = \mathfrak{A}'\mathfrak{B}_{2}, \quad \ldots \quad \mathfrak{C}'_{s} = \mathfrak{A}'\mathfrak{B}_{s}, \]
welche durch ein beliebiges mit $\mathfrak{A}$ äquivalentes Polygon $\mathfrak{A}'$ theilbar sind. (§~18, 8.; §~19, 4.). Diese Zahl $s$ hängt daher nur von den beiden Classen $A$, $C$ ab und kann füglich mit $(A, C)$ bezeichnet werden. Der Werth des Symbols $(A, C)$ ist gleich 0 zu setzen, wenn $C$ nicht durch $A$ theilbar ist. Die Dimension einer Classe $A$ wird hiernach mit $(O, A)$ bezeichnet, wo $O$ die aus dem Nulleck $\mathfrak{O}$ bestehende Classe bedeutet. Ist (nach §~18, 8.)
\[ C = AB, \]
so folgt:
\[ (A, C) = (A, AB) = (O, B); \tag{1.} \]

\origpage{255}
denn die Polygone $\mathfrak{B}_{1}, \mathfrak{B}_{2}, \ldots \mathfrak{B}_{s}$, die sämmtlich in $B$ enthalten sind, sind linear unabhängig, daher $(O, B)$ gewiss nicht kleiner als $s$. Ist umgekehrt $\mathfrak{B}$ ein beliebiges Polygon der Classe $B$, so ist $\mathfrak{A}\mathfrak{B}$ in $C$ enthalten, also auch in der Schaar $(\mathfrak{A}\mathfrak{B}_{1}, \mathfrak{A}\mathfrak{B}_{2}, \ldots \mathfrak{A}\mathfrak{B}_{s})$, mithin $\mathfrak{B}$ in der Schaar $(\mathfrak{B}_{1}, \mathfrak{B}_{2}, \ldots \mathfrak{B}_{s})$ enthalten, d.\ h.\ $(O, B) = s$.

Ist $a$ die Ordnung der Classe $A$, so ist nach §~20, 2. $(A, C) \geqq (O, C) - a$, und daraus folgt mittelst (1.) der allgemeine Satz
\[ (O, B) \geqq (O, AB) - a. \tag{2.} \]

3. Haben die sämmmtlichen\ednote{Printed ``sämmmtlichen'' with three m; read ``sämmtlichen''. An apparent misprint.} Basis-Polygone einer Classe $A$ den grössten gemeinschaftlichen Theiler $\mathfrak{M}$, so ist dieser auch Theiler sämmtlicher Polygone der Classe $A$. Ist $\mathfrak{M}$ gleich dem Nulleck $\mathfrak{O}$, so heisst die Classe eine \emph{eigentliche,} im entgegengesetzten Fall eine \emph{uneigentliche vom Theiler $\mathfrak{M}$}.

Unterdrückt man in sämmtlichen Polygonen einer uneigentlichen Classe $A$ den Theiler $\mathfrak{M}$, so erhält man eine eigentliche Classe $A'$ von niedrigerer Ordnung, aber von derselben Dimension. Diese Beziehung von $A$ zu $A'$ soll durch das Zeichen ausgedrückt sein
\[ A = \mathfrak{M}A'. \]

4. Dêr\ednote{Printed ``Dêr'' with a circumflex-like mark over the e; read ``Der''. Probably a stray or damaged piece of type.} Theiler $\mathfrak{M}$ einer uneigentlichen Classe $A$ ist stets ein isolirtes Polygon. Ist nämlich
\[ A = \mathfrak{M}A', \]
so kann man in der eigentlichen Classe $A'$ nach §~19, 2. ein Polygon $\mathfrak{A}'$ so wählen, dass es relativ prim zu $\mathfrak{M}$ ist. Ist also $\mathfrak{M}'$ äquivalent mit $\mathfrak{M}$, so ist $\mathfrak{M}'\mathfrak{A}'$ äquivalent mit $\mathfrak{M}\mathfrak{A}'$, also in $A$ enthalten, mithin durch $\mathfrak{M}$ theilbar. Es ist also auch $\mathfrak{M}'$ durch $\mathfrak{M}$ theilbar, und da $\mathfrak{M}$ und $\mathfrak{M}'$ von gleicher Ordnung sind,
\[ \mathfrak{M} = \mathfrak{M}'. \]
Hiernach bildet das einzige Polygon $\mathfrak{M}$ eine Classe $M$, und die Bezeichnung $\mathfrak{M}A'$ ist gleichbedeutend mit $MA'$ (§~18, 6.).

\subsection*{§~22. Die Normalbasen von $\mathfrak{o}$.}

1. Wir betrachten im Folgenden das System $\mathfrak{o}$ der ganzen Functionen $\omega$ einer beliebigen Variablen $z$ in $\Omega$ und zugleich das System $\mathfrak{o}'$ der ganzen Functionen $\omega'$ von $z' = \frac{1}{z}$. Aus der Definition der ganzen Functionen

\origpage{256}
erhellt sofort, dass die beiden Systeme $\mathfrak{o}, \mathfrak{o}'$ \emph{nur} die Constanten mit einander gemein haben, dass dagegen \emph{jede} Function $\omega$ durch Multiplication mit einer bestimmten positiven Potenz von $z'$ in eine Function $\omega'$ verwandelt werden kann. Ist $\omega z'^{r}$ in $\mathfrak{o}'$ enthalten, so gilt das Gleiche auch von $\omega z'^{r+1}, \omega z'^{r+2}, \ldots$ In der Reihe der Functionen
\[ \omega, \quad \frac{\omega}{z} = z'\omega, \quad \frac{\omega}{z^{2}} = z'^{2}\omega, \quad \ldots \]
werden also von einem bestimmten Gliede $\omega z'^{r}$ an alle folgenden Functionen in $\mathfrak{o}'$ enthalten sein, während alle vorangehenden nicht darin enthalten sind. Die kleinste Zahl $r$, für welche $z'^{r}\omega$ in $\mathfrak{o}'$ enthalten ist, soll der \emph{Exponent} der Function $\omega$ in Bezug auf $z$ genannt werden. Die Constanten, und nur diese, haben den Exponenten Null. Ist $\omega$ von Null verschieden, und $r$ sein Exponent, so ist $r+1$ der Exponent von $(z-c)\omega$; denn ist $\omega = z^{r}\omega'$, so ist
\[ \begin{aligned} \frac{(z-c)\omega}{z^{r+1}} &= (1-cz')\omega' \quad \text{in } \mathfrak{o}' \text{ enthalten,} \\ \frac{(z-c)\omega}{z^{r}} &= z\omega' - c\omega' \quad \text{nicht in } \mathfrak{o}' \text{ enthalten,} \end{aligned} \]
da zwar $c\omega'$, nicht aber $z\omega' = \frac{\omega}{z^{r-1}}$ in $\mathfrak{o}'$ enthalten ist. Daraus folgt allgemein:

\emph{Ist $x$ eine ganze rationale Function von $z$ vom Grade $s$, und $r$ der Exponent von $\omega$, so ist $(r+s)$ der Exponent von $x\omega$.}

2. Wir wählen nun ein Functionensystem $\lambda_{1}, \lambda_{2}, \ldots \lambda_{n}$ in $\mathfrak{o}$ nach folgender Regel aus:

Es sei $\lambda_{1}$ eine von Null verschiedene Constante, z.\ B.\ 1; $\lambda_{2}$ sei unter denjenigen Functionen in $\mathfrak{o}$, welche nicht einer Constanten nach dem Modul $\mathfrak{o}z$ congruent sind, eine von möglichst niedrigem Exponenten $r_{2}$ u.\ s.\ f.; allgemein sei $\lambda_{s}$ unter denjenigen Functionen in $\mathfrak{o}$, welche nicht congruent sind einer Function der Schaar $(\lambda_{1}, \lambda_{2}, \ldots \lambda_{s-1})\ (\text{mod.}\,\mathfrak{o}z)$, eine von möglichst niedrigem Exponenten $r_{s}$. Da $(\mathfrak{o}, \mathfrak{o}z) = N(z) = z^{n}$ vom $n^{\text{ten}}$ Grade ist, so giebt es in $\mathfrak{o}$ $n$ und nicht mehr nach dem Modul $\mathfrak{o}z$ linear unabhängige Functionen (§~6), und daher kann die Reihe der Functionen $\lambda_{1}, \lambda_{2}, \lambda_{3}, \ldots$ nicht mehr und nicht weniger als $n$ Glieder enthalten. Es ist dann (§~5)
\[ \mathfrak{o} \equiv (\lambda_{1}, \lambda_{2}, \ldots \lambda_{n})\ (\text{mod.}\,\mathfrak{o}z). \]
Die Exponenten $r_{1}, r_{2}, \ldots r_{n}$ der Functionen $\lambda_{1}, \lambda_{2}, \ldots \lambda_{n}$ genügen der Forderung
\[ r_{1} = 0, \quad 1 \leqq r_{2} \leqq r_{3} \cdots \leqq r_{n}. \]

\origpage{257}
Jede Function in $\mathfrak{o}$, deren Exponent $< r_{s}$, ist nach dem Modul $\mathfrak{o}z$ congruent einer Function aus der $(s-1)$-fachen Schaar
\[ (\lambda_{1}, \lambda_{2}, \ldots \lambda_{s-1}). \]
\emph{Diese Functionen $\lambda_{1}, \lambda_{2}, \ldots \lambda_{n}$ bilden eine Basis von $\mathfrak{o}$,} wie sich aus folgender Betrachtung ergiebt.

Wäre es nicht der Fall, so könnte man (§~3, 7.) eine lineare Function $z-c$ und ein System nicht alle verschwindender Constanten $a_{1}, a_{2}, \ldots a_{n}$ so bestimmen, dass
\[ a_{1}\lambda_{1} + a_{2}\lambda_{2} + \cdots + a_{n}\lambda_{n} = (z-c)\omega \]
wäre. Ist unter den Constanten $a$ die letzte nicht verschwindende $a_{s}$, so ist auch
\[ a_{1}\lambda_{1} + a_{2}\lambda_{2} + \cdots + a_{s}\lambda_{s} = (z-c)\omega, \]
und der Exponent von $\omega$ ist sicher kleiner als $r_{s}$ (weil $\frac{(z-c)\omega}{z^{r_{s}}}$ in $\mathfrak{o}'$ enthalten ist). Es ist also $\omega$, und mithin, da $a_{s}$ von 0 verschieden ist, auch $\lambda_{s}$ congruent einer Function der Schaar $(\lambda_{1}, \lambda_{2}, \ldots \lambda_{s-1})\ (\text{mod.}\,\mathfrak{o}z)$, was gegen die Voraussetzung ist.

Die Functionen $\lambda_{1}, \lambda_{2}, \ldots \lambda_{n}$ bilden daher eine Basis von $\mathfrak{o}$, und diese soll \emph{Normalbasis} genannt werden. Die charakteristischen Eigenschaften der Normalbasis sind:

I. Die Functionen $\lambda_{1}, \lambda_{2}, \ldots \lambda_{n}$ sind linear unabhängig nach dem Modul $\mathfrak{o}z$.

II. Jede Function in $\mathfrak{o}$, deren Exponent kleiner ist als der Exponent $r_{s}$ von $\lambda_{s}$, ist in der Form enthalten
\[ c_{1}\lambda_{1} + c_{2}\lambda_{2} + \cdots + c_{s-1}\lambda_{s-1} + z\omega_{s}, \]
worin $c_{1}, c_{2}, \ldots c_{s-1}$ Constanten, $\omega_{s}$ eine Function in $\mathfrak{o}$.

3. Die in $\mathfrak{o}'$ enthaltenen Functionen
\[ \lambda'_{1} = \frac{\lambda_{1}}{z^{r_{1}}}, \quad \lambda'_{2} = \frac{\lambda_{2}}{z^{r_{2}}}, \quad \ldots \quad \lambda'_{n} = \frac{\lambda_{n}}{z^{r_{n}}} \]
bilden eine \emph{Normalbasis} von $\mathfrak{o}'$.

Ist nämlich $\omega$ eine \emph{durch $z$ nicht theilbare} Function in $\mathfrak{o}$ vom Exponenten $r$, so ist der Exponent von $\omega' = \frac{\omega}{z^{r}}$ in Bezug auf $z'$ ebenfalls $r$; denn es ist zwar $\frac{\omega'}{z'^{r}} = \omega$, aber nicht $\frac{\omega'}{z'^{r-1}} = \frac{\omega}{z}$ in $\mathfrak{o}$ enthalten. Da die Functionen $\lambda_{1}, \lambda_{2}, \ldots \lambda_{n}$ alle durch $z$ nicht theilbar sind, so sind hiernach

\origpage{258}
die Exponenten von $\lambda'_{1}, \lambda'_{2}, \ldots \lambda'_{n}$ in Bezug auf $z'$ resp.\ $r_{1}, r_{2}, \ldots r_{n}$. Dies vorausgeschickt beweisen wir, dass das Functionensystem $\lambda'_{1}, \lambda'_{2}, \ldots \lambda'_{n}$ die Eigenschaften I., II. besitzt, wenn dort $\mathfrak{o}$, $z$ durch $\mathfrak{o}'$, $z'$ ersetzt werden.

Wäre die Bedingung I. nicht erfüllt, so liessen sich die Constanten $a_{1}, a_{2}, \ldots a_{s}$, deren letzte nicht verschwindet, so bestimmen, dass
\[ a_{1}\lambda'_{1} + a_{2}\lambda'_{2} + \cdots + a_{s}\lambda'_{s} = z'\omega', \]
also auch (durch Multiplication mit $z^{r_{s}}$)
\[ a_{1}z^{r_{s}-r_{1}}\lambda_{1} + a_{2}z^{r_{s}-r_{2}}\lambda_{2} + \cdots + a_{s}\lambda_{s} = \omega, \]
worin
\[ \omega = z^{r_{s}-1}\omega', \]
also eine Function in $\mathfrak{o}$, deren Exponent kleiner als $r_{s}$ wäre. Dies ist aber, da $a_{s}$ von Null verschieden, wegen der Voraussetzung über die $\lambda$ unmöglich, und folglich die Bedingung I. erfüllt; daraus folgt:
\[ \mathfrak{o}' \equiv (\lambda'_{1}, \lambda'_{2}, \ldots \lambda'_{n})\ (\text{mod.}\,\mathfrak{o}'z'). \]

Wäre die Bedingung II. nicht erfüllt, und $\lambda'$ eine Function in $\mathfrak{o}'$, deren Exponent $r < r_{s}$, die nicht in der Form enthalten ist
\[ a_{1}\lambda'_{1} + a_{2}\lambda'_{2} + \cdots + a_{s-1}\lambda'_{s-1} + z'\omega', \]
so könnte man $e \geqq s$ so wählen, dass
\[ \lambda' = a_{1}\lambda'_{1} + a_{2}\lambda'_{2} + \cdots + a_{e}\lambda'_{e} + z'\omega' \]
mit constanten Coefficienten, deren letzter $a_{e}$ nicht verschwindet. Es ist hiernach auch $r_{e} \geqq r_{s} > r$.

Demnach ist $\lambda = z^{r_{e}-1}\lambda'$ eine Function in $\mathfrak{o}$, und es ergiebt sich durch Multiplication mit $z^{r_{e}}$
\[ z\lambda = a_{1}z^{r_{e}-r_{1}}\lambda_{1} + a_{2}z^{r_{e}-r_{2}}\lambda_{2} + \cdots + a_{e}\lambda_{e} + z^{r_{e}-1}\omega'. \]
Es ist daher $\omega = z^{r_{e}-1}\omega'$ eine Function in $\mathfrak{o}$, deren Exponent (nach 1.) $\leqq r_{e}-1$, und welche der Congruenz genügt
\[ \omega \equiv a'_{1}\lambda_{1} + a'_{2}\lambda_{2} + \cdots + a'_{e}\lambda_{e}\ (\text{mod.}\,\mathfrak{o}z), \]
worin $a'_{e} = -a_{e}$ von Null verschieden ist. Hiernach müsste aber wegen der Eigenschaft II. der Functionen $\lambda$ der Exponent von $\omega \geqq r_{e}$ sein, woraus der Widerspruch erhellt.

Hiermit ist nachgewiesen, dass das Functionensystem $\lambda'_{1}, \lambda'_{2}, \ldots \lambda'_{n}$ eine Normalbasis von $\mathfrak{o}'$ bildet.

\origpage{259}

4. Wir bilden nun die Discriminante von $\Omega$ in Bezug auf die Variable $z$ und $z'$ mit Hülfe der beiden Normalbasen $\lambda$, $\lambda'$; es ist:
\[ \begin{aligned} \Delta_{z}(\Omega) &= \text{const.}\,\Delta(\lambda_{1}, \lambda_{2}, \ldots \lambda_{n}), \\ \Delta_{z'}(\Omega) &= \text{const.}\,\Delta(\lambda'_{1}, \lambda'_{2}, \ldots \lambda'_{n}). \end{aligned} \]
Setzt man aber für $\lambda'_{\iota}$ die Ausdrücke $z'^{r_{\iota}}\lambda_{\iota}$, so folgt aus dem Satze §~2, (13.)
\[ \Delta_{z'}(\Omega) = \text{const.}\,z'^{2(r_{1}+r_{2}+\cdots+r_{n})}\Delta_{z}(\Omega). \]
Ist $\Delta_{z}(\Omega)$ vom Grade $\delta$, so besitzt $\Delta_{z'}(\Omega)$ die Wurzel $z' = 0$ $(2(r_{1}+r_{2}+\cdots+r_{n})-\delta)$-mal, und daraus ergiebt sich nach §~16, 2. die \emph{Verzweigungszahl}
\[ \mathrm{w}_{z} = 2(r_{1}+r_{2}+\cdots+r_{n}), \]
welche hiernach stets eine \emph{gerade Zahl} ist.

\subsection*{§~23. Die Differentialquotienten.}

1. Da eine jede von Null verschiedene Function des Körpers $\Omega$ nur in einer endlichen Anzahl von Punkten den Werth Null hat, so folgt, dass eine Function in $\Omega$, von der sich unendlich viele Nullpunkte nachweisen lassen, nothwendig identisch Null ist, oder dass zwei Functionen in $\Omega$, welche in unendlich vielen Punkten denselben Werth haben, identisch sein müssen.

2. Sind $\alpha$, $\beta$ irgend zwei Variable des Körpers $\Omega$, so existirt in $\Omega$ eine mit $\left(\frac{d\alpha}{d\beta}\right)$ zu bezeichnende Function, welche in unendlich vielen Punkten $\mathfrak{P}$ der Bedingung genügt:
\[ \left(\frac{d\alpha}{d\beta}\right)_{0} = \left(\frac{\alpha-\alpha_{0}}{\beta-\beta_{0}}\right)_{0}, \]
welche der \emph{Differentialquotient} von $\alpha$ nach $\beta$ genannt wird. Ist nämlich $F(\alpha, \beta) = 0$ die zwischen $\alpha$, $\beta$ bestehende irreductible Gleichung, so ist, wenn wir zunächst diejenigen (in endlicher Zahl vorhandenen) Punkte ausschliessen, in welchen $\alpha_{0}$ oder $\beta_{0} = \infty$ oder $F'(\alpha_{0}) = 0$ oder $F'(\beta_{0}) = 0$ ist,
\[ \begin{gathered} 0 = F(\alpha, \beta) = F(\alpha_{0}, \beta_{0}) + (\alpha-\alpha_{0})F'(\alpha_{0}) + (\beta-\beta_{0})F'(\beta_{0}) \\ + \frac{1}{2}\{(\alpha-\alpha_{0})^{2}F''(\alpha_{0}, \alpha_{0}) + 2(\alpha-\alpha_{0})(\beta-\beta_{0})F''(\alpha_{0}, \beta_{0}) + (\beta-\beta_{0})^{2}F''(\beta_{0}, \beta_{0})\} + \cdots \end{gathered} \]
Von den beiden Quotienten $\left(\frac{\alpha-\alpha_{0}}{\beta-\beta_{0}}\right)_{0}$, $\left(\frac{\beta-\beta_{0}}{\alpha-\alpha_{0}}\right)_{0}$ ist gewiss der eine endlich; ist es der erstere, so ziehen wir aus der letzten Gleichung die folgende:

\origpage{260}
\[ \begin{gathered} 0 = \frac{\alpha-\alpha_{0}}{\beta-\beta_{0}}F'(\alpha_{0}) + F'(\beta_{0}) \\ + (\beta-\beta_{0})\frac{1}{2}\left\{\left(\frac{\alpha-\alpha_{0}}{\beta-\beta_{0}}\right)^{2}F''(\alpha_{0}, \alpha_{0}) + 2\left(\frac{\alpha-\alpha_{0}}{\beta-\beta_{0}}\right)F''(\alpha_{0}, \beta_{0}) + F''(\beta_{0}, \beta_{0})\right\} + \cdots, \end{gathered} \]
woraus für den Punkt $\mathfrak{P}$ folgt:
\[ \left(\frac{\alpha-\alpha_{0}}{\beta-\beta_{0}}\right)_{0} = -\frac{F'(\beta_{0})}{F'(\alpha_{0})} = -\left(\frac{F'(\beta)}{F'(\alpha)}\right)_{0}. \]
Wäre $\left(\frac{\alpha-\alpha_{0}}{\beta-\beta_{0}}\right)_{0}$ unendlich, so würden wir ebenso in Bezug auf $\frac{\beta-\beta_{0}}{\alpha-\alpha_{0}}$ schliessen.

Es hat also
\[ \left(\frac{d\alpha}{d\beta}\right) = -\frac{F'(\beta)}{F'(\alpha)} \tag{1.} \]
die verlangte Eigenschaft. Dies bleibt auch noch richtig, wenn von den beiden Functionen $\alpha$, $\beta$ eine constant ist; denn ist z.\ B.\ $\alpha$ constant, so ist $F(\alpha, \beta) = \alpha - \alpha_{0}$ von $\beta$ unabhängig, also $F'(\alpha) = 1$, $F'(\beta) = 0$.

3. Aus Vorstehendem folgt, dass, falls $\beta$ nicht constant ist, abgesehen von einer \emph{endlichen} Anzahl von Punkten $\left(\frac{\alpha-\alpha_{0}}{\beta-\beta_{0}}\right)_{0}$ ein endlicher Werth ist. Ist daher $\gamma$ eine dritte Variable in $\Omega$, so ist in unendlich vielen Punkten
\[ \left(\frac{\alpha-\alpha_{0}}{\beta-\beta_{0}}\right)_{0} = \left(\frac{\alpha-\alpha_{0}}{\gamma-\gamma_{0}}\right)_{0}\left(\frac{\gamma-\gamma_{0}}{\beta-\beta_{0}}\right)_{0}, \]
also auch
\[ \left(\frac{d\alpha}{d\beta}\right)_{0} = \left(\frac{d\alpha}{d\gamma}\right)_{0}\left(\frac{d\gamma}{d\beta}\right)_{0}. \]
Hiernach und nach 1. ist aber die Identität erfüllt:
\[ \left(\frac{d\alpha}{d\beta}\right) = \left(\frac{d\alpha}{d\gamma}\right)\left(\frac{d\gamma}{d\beta}\right){}^{*)}. \tag{2.} \]

4. In Folge dieses letzten Satzes können wir jeder der Functionen $\alpha, \beta, \gamma, \ldots$ des Körpers $\Omega$ eine Function $d\alpha, d\beta, d\gamma, \ldots$ (Differential) in

\textbf{*)} Man kann auch den Differentialquotienten durch die Gleichung
\[ \left(\frac{d\alpha}{d\beta}\right) = -\frac{F'(\beta)}{F'(\alpha)} \]
definiren und durch algebraische Division zum Beweis des Satzes
\[ \left(\frac{d\alpha}{d\beta}\right) = \left(\frac{d\alpha}{d\gamma}\right)\left(\frac{d\gamma}{d\beta}\right) \]
gelangen.

\origpage{261}
der Weise zuordnen, dass allgemein
\[ \frac{d\alpha}{d\beta} = \left(\frac{d\alpha}{d\beta}\right) \]
wird. Die Differentiale der Constanten, und \emph{nur} diese sind Null zu setzen; die übrigen sind völlig bestimmt, sobald eines derselben willkürlich angenommen ist. Besteht zwischen den Variablen $\alpha, \beta, \gamma, \ldots$ eine rationale Gleichung
\[ F(\alpha, \beta, \gamma, \ldots) = 0, \]
so folgt aus derselben
\[ F'(\alpha)\,d\alpha + F'(\beta)\,d\beta + F'(\gamma)\,d\gamma + \cdots = 0; \tag{3.} \]
denn auf dieselbe Weise wie in 2. schliesst man, dass diese Gleichung für unendlich viele Punkte befriedigt ist.

Unmittelbare Folgen des letzten Satzes sind die bekannten Regeln für die Differentiation von Summen, Differenzen, Producten und Quotienten:
\[ d(\alpha \pm \beta) = d\alpha \pm d\beta, \tag{4.} \]
\[ d(\alpha\beta) = \alpha\,d\beta + \beta\,d\alpha, \tag{5.} \]
\[ d\left(\frac{\alpha}{\beta}\right) = \frac{\beta\,d\alpha - \alpha\,d\beta}{\beta^{2}}. \tag{6.} \]

5. Ist $\omega$ eine \emph{ganze} Function von $z$, so wird im Allgemeinen $\frac{d\omega}{dz}$ keine ganze Function von $z$ sein. Es ist aber aus dem Ausdruck (§~3, 7.)
\[ \omega = x_{1}\omega_{1} + x_{2}\omega_{2} + \cdots + x_{n}\omega_{n} \]
ersichtlich, da die Differentialquotienten der ganzen rationalen Functionen $x_{1}, x_{2}, \ldots x_{n}$ wieder ganze rationale Functionen sind, dass die Unterideale der sämmtlichen Functionen $\frac{d\omega}{dz}$ in einem bestimmten Ideal aufgehen müssen, nämlich in dem kleinsten gemeinschaftlichen Vielfachen der Unterideale von $\frac{d\omega_{1}}{dz}, \frac{d\omega_{2}}{dz}, \ldots \frac{d\omega_{n}}{dz}$. Es soll untersucht werden, welches dies Ideal ist. Zu dem Ende sei $z-c$ eine beliebige lineare Function von $z$ und
\[ \mathfrak{o}(z-c) = \mathfrak{p}^{e}\mathfrak{p}_{1}^{e_{1}}\mathfrak{p}_{2}^{e_{2}}\ldots, \]
worin die Primideale $\mathfrak{p}, \mathfrak{p}_{1}, \mathfrak{p}_{2}, \ldots$ von einander verschieden sind. Es sei nun $\zeta$ dieselbe Function wie in §~11, 2., d.\ h.\ eine ganze Function von $z$, welche in den durch die Primideale $\mathfrak{p}, \mathfrak{p}_{1}, \mathfrak{p}_{2}, \ldots$ erzeugten Punkten $\mathfrak{P}, \mathfrak{P}_{1}, \mathfrak{P}_{2}, \ldots$ lauter verschiedene Werthe hat und jeden derselben nur einfach;

\origpage{262}
dann lässt sich $\omega$ in der Form darstellen
\[ \omega = y_{0} + y_{1}\zeta + \cdots + y_{n-1}\zeta^{n-1}, \]
worin die rationalen Functionen $y_{0}, y_{1}, \ldots y_{n-1}$ von $z$ zwar gebrochen sein können, aber den Factor $z-c$ gewiss nicht im Nenner enthalten. Daraus folgt, dass das Unterideal von $\frac{d\omega}{dz}$ durch keine höheren Potenzen der Ideale $\mathfrak{p}, \mathfrak{p}_{1}, \mathfrak{p}_{2}, \ldots$ theilbar sein kann, als das Unterideal von $\frac{d\zeta}{dz}$. Ist aber
\[ f(\zeta, z) = 0 \]
die zwischen $\zeta$ und $z$ bestehende irreductible Gleichung, so ist nach §~11, 2.
\[ \mathfrak{o}f'(\zeta) = \mathfrak{m}\mathfrak{p}^{e-1}\mathfrak{p}_{1}^{e_{1}-1}\mathfrak{p}_{2}^{e_{2}-1}\ldots \]
und $\mathfrak{m}$ relativ prim zu $\mathfrak{p}, \mathfrak{p}_{1}, \mathfrak{p}_{2}, \ldots$ Da aber
\[ \frac{d\zeta}{dz} = -\frac{f'(z)}{f'(\zeta)} \]
ist, so kann das Unterideal von $\frac{d\zeta}{dz}$, und mithin auch das von $\frac{d\omega}{dz}$ keinen der Factoren $\mathfrak{p}, \mathfrak{p}_{1}, \mathfrak{p}_{2}, \ldots$ öfter als $(e-1)$, $(e_{1}-1)$, $(e_{2}-1)$, \ldots\ mal enthalten. Da nun $z-c$ jede beliebige lineare Function sein kann, so folgt, dass $\frac{d\omega}{dz}$ kein anderes Unterideal haben kann, als ein solches, welches in dem Verzweigungsideal $\mathfrak{z} = \Pi\mathfrak{p}^{e-1}$ (§~11) aufgeht. Es ist also, wenn $\mathfrak{a}$ ein Ideal bedeutet:
\[ \mathfrak{z}\frac{d\omega}{dz} = \mathfrak{a}, \]
also nach §~11, (7.)
\[ \mathfrak{o}\frac{d\omega}{dz} = \mathfrak{e}\mathfrak{a}, \]
woraus hervorgeht, dass die Functionen $\frac{d\omega}{dz}$ sämmtlich dem zu $\mathfrak{o}$ complementären Modul $\mathfrak{e}$ angehören.

6. Ist die irreductible Gleichung $F(\omega, z) = 0$ zwischen $\omega$ und $z$ vom $n^{\text{ten}}$ Grade in Bezug auf $\omega$, also $1, \omega, \omega^{2}, \ldots \omega^{n-1}$ eine Basis von $\Omega$, so ist nach §~11, (10.)
\[ \mathfrak{o}F'(\omega) = \mathfrak{z}\mathfrak{k}, \]
und daher muss wegen
\[ \frac{d\omega}{dz} = -\frac{F'(z)}{F'(\omega)} \]

\origpage{263}
$\mathfrak{o}F'(z)$ durch das Ideal $\mathfrak{k}$ theilbar sein,
\[ \mathfrak{o}F'(z) = \mathfrak{k}\mathfrak{a}, \]
$\mathfrak{k}$ kann man daher das \emph{Ideal der Doppelpunkte} in Bezug auf $\omega$, $z$ nennen.

7. Ist $\mathfrak{P}$ ein Punkt, in welchem $z-c$ unendlich klein in der ersten Ordnung ist (also kein Verzweigungspunkt in $z$), so sind nach 5. die Functionen $\frac{d\omega}{dz}$ in $\mathfrak{P}$ alle endlich. Ist also $\eta$ irgend eine Function in $\Omega$, welche in $\mathfrak{P}$ endlich ist, so kann man diese als Quotienten zweier ganzen Functionen $\frac{\alpha}{\beta}$ darstellen, von denen $\beta$ in $\mathfrak{P}$ nicht verschwindet, und daher ist nach (6.) auch $\frac{d\eta}{dz}$ in $\mathfrak{P}$ endlich.

8. Es seien jetzt $\alpha$, $\beta$ irgend zwei Variable in $\Omega$; es soll das Verhalten von $\frac{d\alpha}{d\beta}$ in irgend einem Punkt $\mathfrak{P}$ untersucht werden.

Man wähle eine Variable $z$ in $\Omega$, welche in $\mathfrak{P}$ unendlich klein in der ersten Ordnung ist. Hat $\alpha$ in $\mathfrak{P}$ einen endlichen Werth $\alpha_{0}$, so kann man nach §~15, 1., 2. eine positive ganze Zahl $r$ und eine in $\mathfrak{P}$ endliche und von Null verschiedene Function $\alpha'$ so bestimmen, dass
\[ \alpha = \alpha_{0} + z^{r}\alpha' \]
wird. Dies gilt auch noch, wenn $\alpha$ in $\mathfrak{P}$ unendlich ist; nur ist dann $r$ eine negative ganze Zahl, und $\alpha_{0}$ ist durch eine beliebige endliche Constante, z.\ B.\ 0 zu ersetzen. Ebenso kann man
\[ \beta = \beta_{0} + z^{s}\beta' \]
setzen; $r$ und $s$ sind dann die Ordnungszahlen von $\alpha-\alpha_{0}$, $\beta-\beta_{0}$ im Punkte $\mathfrak{P}$, die sowohl positiv als negativ, aber nicht 0 sein können. Aus (2.) ergiebt sich dann:
\[ \frac{d\alpha}{d\beta} = z^{r-s}\frac{r\alpha' + z\frac{d\alpha'}{dz}}{s\beta' + z\frac{d\beta'}{dz}} \]
oder
\[ \frac{\beta-\beta_{0}}{\alpha-\alpha_{0}}\frac{d\alpha}{d\beta} = \frac{r + z\frac{d\alpha'}{\alpha'dz}}{s + z\frac{d\beta'}{\beta'dz}}. \]
Bezeichnet man nun wieder durch den Index 0 den Werth einer Function im Punkte $\mathfrak{P}$, so ist, da
\[ \left(\frac{d\alpha'}{\alpha'dz}\right)_{0}, \quad \left(\frac{d\beta'}{\beta'dz}\right)_{0} \]

\origpage{264}
nach 7. endlich sind,
\[ \left(\frac{\beta-\beta_{0}}{\alpha-\alpha_{0}}\frac{d\alpha}{d\beta}\right)_{0} = \frac{r}{s}, \tag{7.} \]
also endlich und von Null verschieden. Hieraus ergiebt sich, dass \emph{die Ordnungszahl des Differentialquotienten $\frac{d\alpha}{d\beta}$ gleich ist der Differenz der Ordnungszahlen von $\alpha-\alpha_{0}$ und $\beta-\beta_{0}$}. Ist $r \gtrless s$, so ist $\left(\frac{\alpha-\alpha_{0}}{\beta-\beta_{0}}\right)_{0}$ und mithin $\left(\frac{d\alpha}{d\beta}\right)_{0}$ Null oder unendlich. Ist dagegen $r = s$, so sind beide Werthe endlich und von 0 verschieden, und wir haben daher in allen Fällen
\[ \left(\frac{\alpha-\alpha_{0}}{\beta-\beta_{0}}\right)_{0} = \left(\frac{d\alpha}{d\beta}\right)_{0}. \tag{8.} \]
Hierin sind $\alpha_{0}$, $\beta_{0}$ die Werthe von $\alpha$, $\beta$ in $\mathfrak{P}$, wenn diese Werthe endlich sind, sonst beliebige Constanten, z.\ B.\ 0.

9. Sind $a$, $b$ die Ordnungszahlen von $\alpha-\alpha_{0}$, $\beta-\beta_{0}$ in $\mathfrak{P}$, so kommt, falls $a$, $b$ positiv sind, der Punkt $\mathfrak{P}$ $(a-1)$-mal resp.\ $(b-1)$-mal in den Verzweigungspolygonen $\mathfrak{Z}_{\alpha}$, $\mathfrak{Z}_{\beta}$ in $\alpha$, $\beta$ vor. Ist aber $a$ negativ, so enthält $\mathfrak{Z}_{\alpha}$ den Punkt $\mathfrak{P}$ $(-a-1)$-mal, und Entsprechendes gilt, wenn $b$ negativ ist (§~16, 1.). Bezeichnet man also mit $\mathfrak{A}$, $\mathfrak{B}$ die Unterecke von $\alpha$, $\beta$, so erhält man, weil die Ordnungszahl von $\frac{d\alpha}{d\beta}$ (wie eben bewiesen) immer gleich $a-b$ ist, für diese Function folgenden Ausdruck als Polygonquotienten
\[ \frac{d\alpha}{d\beta} = \frac{\mathfrak{Z}_{\alpha}\mathfrak{B}^{2}}{\mathfrak{Z}_{\beta}\mathfrak{A}^{2}}. \tag{9.} \]

\subsection*{§~24. Das Geschlecht des Körpers $\Omega$.}

1. Bezeichnet man mit $\mathrm{w}_{\alpha}$, $\mathrm{w}_{\beta}$ die Verzweigungszahlen, mit $n_{\alpha}$, $n_{\beta}$ die Ordnungen der Variablen $\alpha$, $\beta$, so folgt aus der Formel (9.) des vorigen §, da Zähler und Nenner von $\frac{d\alpha}{d\beta}$ gleichviel Punkte enthalten müssen, die wichtige Relation
\[ \mathrm{w}_{\alpha} - 2n_{\alpha} = \mathrm{w}_{\beta} - 2n_{\beta}; \]
wenn man also
\[ p = \frac{1}{2}\mathrm{w} - n + 1 \tag{1.} \]
setzt, welches nach §~22, 4. eine ganze Zahl ist, so ist diese von der Wahl der Variablen unabhängig und eine für den Körper $\Omega$ charakteristische Zahl,

\origpage{265}
welche das \emph{Geschlecht} des Körpers $\Omega$ genannt wird. Dass diese Zahl niemals negativ ist, ergiebt sich, wenn man für $\frac{1}{2}\mathrm{w}$ den Werth $r_{1}+r_{2}+\cdots+r_{n}$ aus §~22 einsetzt. Man erhält dann
\[ p = (r_{2}-1)+(r_{3}-1)+\cdots+(r_{n}-1), \tag{2.} \]
was, da $r_{2}, r_{3}, \ldots r_{n} \geqq 1$ sind, nicht negativ werden kann.

2. Es seien $\alpha$, $\beta$ zwei Functionen in $\Omega$ von den Ordnungen $m$, $n$, von der Beschaffenheit, dass \emph{alle} Functionen in $\Omega$ rational durch $\alpha$, $\beta$ darstellbar sind. Es ist dann
\[ \begin{aligned} F(\alpha, \beta) &= a_{0}\alpha^{n} + a_{1}\alpha^{n-1} + \cdots + a_{n-1}\alpha + a_{n} \\ &= b_{0}\beta^{m} + b_{1}\beta^{m-1} + \cdots + b_{m-1}\beta + b_{m} = 0 \end{aligned} \]
die zwischen $\alpha$, $\beta$ bestehende irreductible Gleichung, worin $a_{0}, a_{1}, \ldots a_{n}$ ganze rationale Functionen von $\beta$, ebenso $b_{0}, b_{1}, \ldots b_{m}$ ganze rationale Functionen von $\alpha$ sind.

Es sei ferner
\[ \alpha = \frac{\mathfrak{A}_{1}}{\mathfrak{A}}, \quad \beta = \frac{\mathfrak{B}_{1}}{\mathfrak{B}} \]
und $\mathfrak{A}_{1}$ relativ prim zu $\mathfrak{A}$, $\mathfrak{B}_{1}$ zu $\mathfrak{B}$, so dass $\mathfrak{A}$, $\mathfrak{A}_{1}$ von der Ordnung $m$, $\mathfrak{B}$, $\mathfrak{B}_{1}$ von der Ordnung $n$ sind. Nun ist
\[ \begin{aligned} F'(\alpha) &= na_{0}\alpha^{n-1} + (n-1)a_{1}\alpha^{n-2} + \cdots + a_{n-1}, \\ \alpha F'(\alpha) &= -a_{1}\alpha^{n-1} - 2a_{2}\alpha^{n-2} - \cdots - na_{n}, \end{aligned} \]
woraus hervorgeht, dass
\[ F'(\alpha) = \frac{\mathfrak{K}}{\mathfrak{A}^{n-2}\mathfrak{B}^{m}} \]
und ebenso
\[ F'(\beta) = \frac{\mathfrak{L}}{\mathfrak{A}^{n}\mathfrak{B}^{m-2}} \]
sein muss. Es ist nun nachzuweisen, dass das Polygon $\mathfrak{K}$ durch $\mathfrak{Z}_{\beta}$, $\mathfrak{L}$ durch $\mathfrak{Z}_{\alpha}$ theilbar ist.

Für $\mathfrak{K}$ ist dies leicht einzusehen unter der Voraussetzung, dass in sämmtlichen Punkten von $\mathfrak{Z}_{\beta}$ die Function $\beta$ einen endlichen, und $a_{0}$ einen von Null verschiedenen Werth hat; denn es ist
\[ \alpha' = a_{0}\alpha \]
eine ganze Function von $\beta$, und wenn man
\[ f(\alpha') = a^{n-1}F(\alpha, \beta) \]
\ednote{Printed $a^{n-1}$ without the subscript; $a_{0}^{n-1}$ is meant, as $\alpha' = a_{0}\alpha$ and the next line has $a_{0}^{n-2}$. Possibly a dropped subscript.}
setzt, so ist
\[ f'(\alpha') = a_{0}^{n-2}F'(\alpha). \]

\origpage{266}
Da nun nach §~11, 5. $\mathfrak{o}_{\beta}f'(\alpha')$ durch das von $\mathfrak{Z}_{\beta}$ erzeugte Verzweigungsideal in $\beta$ theilbar ist, so folgt hieraus die Richtigkeit der Behauptung. Analoges gilt für $F'(\beta)$.

Macht man nun für $\alpha$, $\beta$ beliebige lineare Substitutionen:
\[ \begin{gathered} \alpha = \frac{c+d\alpha'}{a+b\alpha'}; \quad \beta = \frac{c'+d'\beta'}{a'+b'\beta'}, \\ (a+b\alpha')(d-b\alpha) = ad-bc, \\ (a'+b'\beta')(d'-b'\beta) = a'd'-b'c', \end{gathered} \]
so ist nach §~16, 2.
\[ \mathfrak{Z}_{\alpha} = \mathfrak{Z}_{\alpha'}; \quad \mathfrak{Z}_{\beta} = \mathfrak{Z}_{\beta'}, \]
und die zwischen $\alpha'$, $\beta'$ bestehende irreductible Gleichung lautet:
\[ F_{1}(\alpha', \beta') = (a+b\alpha')^{n}(a'+b'\beta')^{m}F(\alpha, \beta) = 0. \]
Es lassen sich aber unter allen Umständen die Constanten $a$, $b$, $c$, $d$; $a'$, $b'$, $c'$, $d'$ so wählen, dass die oben angegebenen Voraussetzungen sowohl für $\alpha'$ als für $\beta'$ erfüllt sind.

Denn setzt man die Coefficienten $a'_{0}$, $b'_{0}$ von $\alpha'^{n}$, $\beta'^{m}$ in $F_{1}(\alpha', \beta')$ in die Form
\[ \begin{aligned} a'_{0} &= (a'+b'\beta')^{m}(a_{0}d^{n} + a_{1}d^{n-1}b + \cdots + a_{n}b^{n}) = \left(\frac{a'd'-b'c'}{d'-b'\beta}\right)^{m}(a_{0}d^{n} + a_{1}d^{n-1}b + \cdots + a_{n}b^{n}), \\ b'_{0} &= (a+b\alpha')^{n}(b_{0}d'^{m} + b_{1}d'^{m-1}b' + \cdots + b_{m}b'^{m}) = \left(\frac{ad-bc}{d-b\alpha}\right)^{n}(b_{0}d'^{m} + b_{1}d'^{m-1}b' + \cdots + b_{m}b'^{m}), \end{aligned} \]
so erkennt man leicht, dass nur für eine endliche Anzahl von Werthen der Verhältnisse $d : b$, $d' : b'$ die Functionen $a'_{0}$, $d'-b'\beta$ in einem Punkt von $\mathfrak{Z}_{\beta}$, $b'_{0}$, $d-b\alpha$ in einem Punkt von $\mathfrak{Z}_{\alpha}$ verschwinden können.

Setzen wir nun
\[ \alpha' = \frac{\mathfrak{A}'_{1}}{\mathfrak{A}'}, \quad \beta' = \frac{\mathfrak{B}'_{1}}{\mathfrak{B}'}, \]
so folgt (§~19, 1.)
\[ d-b\alpha = \frac{\mathfrak{A}_{2}}{\mathfrak{A}}, \quad a+b\alpha' = \frac{\mathfrak{A}'_{2}}{\mathfrak{A}'}, \]
also:
\[ \mathfrak{A}_{2}\mathfrak{A}'_{2} = \mathfrak{A}\mathfrak{A}'. \]
Ist aber, wie angenommen, $b$ von Null verschieden, so ist $\mathfrak{A}_{2}$ relativ prim zu $\mathfrak{A}$, weil in einem Punkt von $\mathfrak{A}$ die Ordnungszahl von $d-b\alpha$ dieselbe ist, wie die von $\alpha$ (§~15, 5.) und folglich
\[ \mathfrak{A}_{2} = \mathfrak{A}', \quad \mathfrak{A}'_{2} = \mathfrak{A}, \]

\origpage{267}
also:
\[ a+b\alpha' = \frac{\mathfrak{A}}{\mathfrak{A}'}, \]
und ebenso:
\[ a'+b'\beta' = \frac{\mathfrak{B}}{\mathfrak{B}'}. \]
Nun ist aber, da $F(\alpha, \beta) = 0$ ist:
\[ F'_{1}(\alpha') = (ad-bc)(a+b\alpha')^{n-2}(a'+b'\beta')^{m}F'(\alpha), \]
und wenn also, wie vorausgesetzt:
\[ F'_{1}(\alpha') = \frac{\mathfrak{R}\mathfrak{Z}_{\beta}}{\mathfrak{A}'^{n-2}\mathfrak{B}'^{m}}, \]
so folgt
\[ F'(\alpha) = \frac{\mathfrak{R}\mathfrak{Z}_{\beta}}{\mathfrak{A}^{n-2}\mathfrak{B}^{m}} \]
und in gleicher Weise
\[ F'(\beta) = \frac{\mathfrak{R}\mathfrak{Z}_{\alpha}}{\mathfrak{A}^{n}\mathfrak{B}^{m-2}}. \]
Dass das im Zähler dieser beiden Ausdrücke auftretende Polygon $\mathfrak{R}$ in beiden Ausdrücken dasselbe sein muss, ergiebt sich aus
\[ \frac{d\alpha}{d\beta} = -\frac{F'(\beta)}{F'(\alpha)} = \frac{\mathfrak{B}^{2}\mathfrak{Z}_{\alpha}}{\mathfrak{A}^{2}\mathfrak{Z}_{\beta}}. \]
Nun ist die Ordnung des Polygons $\mathfrak{A}^{n-2}\mathfrak{B}^{m}$
\[ m(n-2)+mn = 2m(n-1), \]
also die Ordnung von $\mathfrak{R}$
\[ 2r = 2m(n-1)-\mathrm{w}_{\beta} \]
stets eine gerade Zahl, und daraus ergiebt sich
\[ p = \frac{1}{2}\mathrm{w}_{\beta}-n+1 = (n-1)(m-1)-r. \tag{3.} \]
Das Polygon $\mathfrak{R}$ wird das \emph{Polygon der Doppelpunkte} in $(\alpha, \beta)$ genannt.

\subsection*{§~25. Die Differentiale in $\Omega$.}

Sind $z$, $z_{1}$ irgend zwei Variable in $\Omega$ von den Ordnungen $n$, $n_{1}$ und den Verzweigungszahlen $\mathrm{w}$, $\mathrm{w}_{1}$, ferner $\mathfrak{Z}$, $\mathfrak{Z}_{1}$ die Verzweigungspolygone, $\mathfrak{U}$, $\mathfrak{U}_{1}$ die Unterecke von $z$, $z_{1}$, so ist (§~23)
\[ \frac{dz}{dz_{1}} = \frac{\mathfrak{Z}\mathfrak{U}_{1}^{2}}{\mathfrak{Z}_{1}\mathfrak{U}^{2}}. \tag{1.} \]

\origpage{268}
Jede Function $\omega$ in $\Omega$ lässt sich in die Form setzen
\[ \omega = \frac{\mathfrak{U}^{2}\mathfrak{A}}{\mathfrak{Z}\mathfrak{B}}, \tag{2.} \]
worin $\mathfrak{A}$, $\mathfrak{B}$ Polygone bedeuten, deren Ordnungen $a$, $b$ der Bedingung genügen
\[ 2n+a = \mathrm{w}+b \]
oder (§~24)
\[ a = b+2p-2. \tag{3.} \]
Wenn man nun eine Function $\omega_{1}$ durch die Gleichung erklärt
\[ \omega\,dz = \omega_{1}\,dz_{1}, \]
so erhält nach (1.) $\omega_{1}$ die Bezeichnung
\[ \omega_{1} = \frac{\mathfrak{U}_{1}^{2}\mathfrak{A}}{\mathfrak{Z}_{1}\mathfrak{B}}. \]
Wir nennen in der Folge solche Ausdrücke, wie
\[ \omega\,dz = \omega_{1}\,dz_{1} \]
\emph{Differentiale in $\Omega$}, und bezeichnen dieselben in symbolischer Weise durch ein Zeichen wie $d\varpi$. Ein solches Differential ist hierdurch \emph{invariant}, d.\ h.\ unabhängig von der Wahl der Veränderlichen $z$ erklärt und ist durch die beiden Polygone $\mathfrak{A}$, $\mathfrak{B}$ vollständig bestimmt.

Wir können ohne Gefahr eines Missverständnisses die symbolische Bezeichnung
\[ d\varpi = \frac{\mathfrak{A}}{\mathfrak{B}}, \]
also beispielsweise auch
\[ dz = \frac{\mathfrak{Z}}{\mathfrak{U}^{2}} \]
anwenden. Diese Bezeichnung eines Differentials durch einen Polygonquotienten unterscheidet sich von der ähnlichen Bezeichnung der Functionen in $\Omega$ (§~17) dadurch, dass bei letzterer Zähler und Nenner von gleicher Ordnung sind, während bei den Differentialen die Ordnung des Zählers die des Nenners um $2p-2$ übertrifft. Wie bei der Bezeichnung in §~17, können auch hier gemeinschaftliche Theiler, welche $\mathfrak{A}$ und $\mathfrak{B}$ etwa enthalten, unterdrückt werden. Sind $\mathfrak{A}$ und $\mathfrak{B}$ relativ prim, so heisst $\mathfrak{A}$ das Obereck, $\mathfrak{B}$ das Untereck des Differentials $d\varpi$.

Unter den hier aufgestellten allgemeinen Begriff des Differentials in $\Omega$ fallen als specielle Fälle auch die in §~23, 4. erklärten Differentiale der Functionen des Körpers $\Omega$. Diese nennen wir \emph{eigentliche Differentiale,}

\origpage{269}
während die anderen, welche nicht als Differentiale von in $\Omega$ existirenden Functionen dargestellt werden können, \emph{uneigentliche oder Abelsche Differentiale} genannt werden.

Functionen von der Form (2.), die nach unserer jetzt getroffenen Festsetzung mit $\frac{d\varpi}{dz}$ bezeichnet werden können, nennen wir \emph{Differentialquotienten nach $z$} und unterscheiden gleichfalls zwischen \emph{eigentlichen} und \emph{uneigentlichen} Differentialquotienten, je nachdem $d\varpi$ ein eigentliches oder uneigentliches Differential ist${}^{*)}$.

Es entsteht nun die Aufgabe, den Umfang des Begriffs der Differentiale festzustellen, d.\ h.\ alle Polygone $\mathfrak{A}$, $\mathfrak{B}$ zu finden, welche Ober- und Untereck eines Differentials sein können. Wir schicken darüber die folgenden allgemeinen Bemerkungen voraus:

Die nothwendige und hinreichende Bedingung dafür, dass $\frac{\mathfrak{A}}{\mathfrak{B}}$ ein Differential sei, ist die, dass für eine beliebige Variable $z$
\[ \frac{\mathfrak{U}^{2}\mathfrak{A}}{\mathfrak{Z}\mathfrak{B}} \]
eine Function in $\Omega$ ist, also dass $\mathfrak{U}^{2}\mathfrak{A}$ mit $\mathfrak{Z}\mathfrak{B}$ äquivalent ist. Dies Verhältniss bleibt aber bestehen, wenn $\mathfrak{A}$, $\mathfrak{B}$ selbst durch äquivalente Polygone $\mathfrak{A}'$, $\mathfrak{B}'$ ersetzt werden. Halten wir $\mathfrak{B}$ fest, und ist $\frac{\mathfrak{A}}{\mathfrak{B}}$ ein Differential, so werden hiernach
\[ \frac{\mathfrak{A}'}{\mathfrak{B}}, \quad \frac{\mathfrak{A}''}{\mathfrak{B}}, \quad \ldots \]
dann und nur dann Differentiale darstellen, wenn die Polygone $\mathfrak{A}$, $\mathfrak{A}'$, $\mathfrak{A}''$, $\ldots$ alle derselben Classe $A$ angehören. Bilden die Polygone $\mathfrak{A}_{1}$, $\mathfrak{A}_{2}$, $\mathfrak{A}_{3}$, $\ldots$ eine Basis von $A$, ist also
\[ A = (\mathfrak{A}_{1}, \mathfrak{A}_{2}, \mathfrak{A}_{3}, \ldots), \]
so bilden die zugehörigen Differentialquotienten in Bezug auf eine beliebige Variable $z$, $\frac{d\varpi_{1}}{dz}$, $\frac{d\varpi_{2}}{dz}$, $\frac{d\varpi_{3}}{dz}$, $\ldots$ die Basis einer Functionenschaar von endlicher Dimension, und dem entsprechend werden wir auch $d\varpi_{1}$, $d\varpi_{2}$, $d\varpi_{3}$, $\ldots$ die Basis einer \emph{Schaar von Differentialen}
\[ (d\varpi_{1}, d\varpi_{2}, d\varpi_{3}, \ldots) \]

\textbf{*)} Der Quotient irgend zweier eigentlichen oder uneigentlichen Differentiale $\frac{d\varpi}{d\varpi'}$ hat stets die Bedeutung einer bestimmten Function in $\Omega$. Wir beschränken uns im Folgenden aber auf die Betrachtung solcher Quotienten, bei denen wenigstens der Nenner ein eigentliches Differential ist.

\origpage{270}
von derselben Dimension nennen. Dies besagt, dass jedes Differential $d\varpi$, dessen Untereck $\mathfrak{B}$ oder ein Theiler von $\mathfrak{B}$ ist, in der Form dargestellt werden kann
\[ d\varpi = c_{1}\,d\varpi_{1} + c_{2}\,d\varpi_{2} + c_{3}\,d\varpi_{3} + \cdots \]
mit constanten Coefficienten $c_{1}$, $c_{2}$, $c_{3}$, $\ldots$

\subsection*{§~26. Die Differentiale erster Gattung.}

Wir betrachten zunächst die einfachsten unter den Differentialen in $\Omega$, nämlich die, deren Untereck das Nulleck $\mathfrak{O}$ ist. Solche Differentiale (deren Existenz freilich erst noch nachzuweisen ist) heissen \emph{Differentiale erster Gattung}. Das Obereck $\mathfrak{W}$ eines solchen Differentials $dw$, dessen Ordnung $2p-2$ ist, wird als das \emph{Grundpolygon} von $dw$ bezeichnet und heisst ein \emph{vollständiges Polygon erster Gattung}, während jeder Theiler eines solchen ein \emph{Polygon erster Gattung} schlechtweg genannt wird. Ist $\mathfrak{W} = \mathfrak{A}\mathfrak{B}$, so heissen $\mathfrak{A}$, $\mathfrak{B}$ \emph{Ergänzungspolygone} von einander. Ein Polygon, welches nicht Theiler eines vollständigen Polygons erster Gattung ist, also ins Besondere jedes Polygon von mehr als $2p-2$ Punkten heisst ein \emph{Polygon zweiter Gattung}.

1. Nach dem oben Bemerkten bilden alle vollständigen Polygone erster Gattung eine Polygonclasse $W$, deren Dimension zu bestimmen ist; ergiebt sich diese Dimension $> 0$, so ist damit zugleich die Existenz der Polygone erster Gattung nachgewiesen. Diese Dimension ist aber dieselbe wie die Dimension der Schaar der Differentiale erster Gattung oder auch, für eine beliebige Variable $z$, der Schaar der \emph{Differentialquotienten erster Gattung}, wenn wir als Differentialquotienten erster Gattung nach $z$ die Functionen
\[ u = \frac{dw}{dz} \]
bezeichnen. Eine solche Function $u$ hat nach §~25, (2.) den Ausdruck
\[ u = \frac{\mathfrak{U}^{2}\mathfrak{W}}{\mathfrak{Z}}, \]
und man erkennt leicht aus der Betrachtung der Ordnungszahlen in den verschiedenen Punkten, dass ein solcher Differentialquotient erster Gattung durch folgende beiden Eigenschaften vollkommen definirt ist:

I. In jedem Punkt $\mathfrak{P}$, in welchem $z$ einen endlichen Werth $z_{0}$ hat, ist
\[ (u(z-z_{0}))_{0} = 0. \]

\origpage{271}

II. In einem Punkt $\mathfrak{P}$, in welchem $z$ unendlich ist, ist
\[ (zu)_{0} = 0. \]

Bedeutet wie in §~11, 4.
\[ r = (z-c)(z-c_{1})(z-c_{2})\ldots \]
das Product sämmtlicher von einander verschiedenen Linearfactoren der Discriminante $\Delta_{z}(\Omega)$, $\mathfrak{r}$ das Product sämmtlicher von einander verschiedenen in $r$ aufgehenden Primideale, so ist die Bedingung I. vollkommen gleichbedeutend mit der, dass $ru$ eine Function in $\mathfrak{r}$, oder dass $u$ eine Function des zu $\mathfrak{o}$ complementären Moduls $\mathfrak{e}$ sein muss (§~11, 4. (6.)). Um also die Gesammtheit der Functionen $u$ zu erhalten, hat man unter den Functionen in $\mathfrak{e}$ diejenigen aufzusuchen, welche der Bedingung II. genügen.

2. Zu diesem Zweck legen wir eine \emph{Normalbasis} $\lambda_{1}, \lambda_{2}, \ldots \lambda_{n}$ von $\mathfrak{o}$ zu Grunde (§~22) und bezeichnen die dazu complementäre Basis mit $\mu_{1}, \mu_{2}, \ldots \mu_{n}$, so dass jede der Bedingung I. genügende Function, also auch jeder Differentialquotient erster Gattung, in der Form enthalten ist
\[ u = y_{1}\mu_{1} + y_{2}\mu_{2} + \cdots + y_{n}\mu_{n}, \tag{1.} \]
worin $y_{1}, y_{2}, \ldots y_{n}$ \emph{ganze} rationale Functionen von $z$ sind. Aus den Grundeigenschaften der complementären Basis ergiebt sich aber (§~10, 3.)
\[ y_{s} = S(u\lambda_{s}); \quad \frac{y_{s}}{z^{r_{s}-1}} = S\left(uz\cdot\frac{\lambda_{s}}{z^{r_{s}}}\right). \]
Da nun $\frac{\lambda_{s}}{z^{r_{s}}}$ in $\mathfrak{o}'$ enthalten, also für $z = \infty$ endlich ist, und $uz$ nach II. in jedem solchen Punkt verschwindet, so folgt (§~16, 5.), dass $\frac{y_{s}}{z^{r_{s}-1}}$ für $z = \infty$ verschwinden muss, d.\ h.\ dass die ganze rationale Function $y_{s}$ den Grad $r_{s}-2$ nicht übersteigen kann.

Es muss daher, falls $r_{s} < 2$ ist, $y_{s}$ verschwinden, also ist unter allen Umständen (§~22, 2.)
\[ y_{1} = 0; \quad S(u) = 0 \]
(\emph{Abelsches Theorem für Differentiale erster Gattung}) und, falls $r_{s} \geqq 2$:
\[ y_{s} = c_{0} + c_{1}z + c_{2}z^{2} + \cdots + c_{r_{s}-2}z^{r_{s}-2}. \tag{2.} \]
Es ist noch zu zeigen, dass diese Bedingungen auch hinreichend sind, d.\ h.\ dass jede Function von der Form (1.), in welcher die $y_{s}$ den Ausdruck (2.) haben, der Forderung II. genügt, oder, was dasselbe ist, dass, wenn $r_{s} \geqq 2$ ist, $z^{r_{s}-1}\mu_{s}$ in allen Punkten, in welchen $z$ unendlich wird, verschwindet.

\origpage{272}
Dies ergiebt sich sofort durch die Betrachtung des Systems $\mathfrak{o}'$ der ganzen Functionen von $z' = \frac{1}{z}$, für welches nach §~22, 3. die Functionen
\[ \lambda'_{1} = \frac{\lambda_{1}}{z^{r_{1}}}, \quad \lambda'_{2} = \frac{\lambda_{2}}{z^{r_{2}}}, \quad \ldots \quad \lambda'_{n} = \frac{\lambda_{n}}{z^{r_{n}}} \]
eine Normalbasis bilden. Die hierzu complementäre Basis ist nach §~10, 5.
\[ \mu'_{1} = z^{r_{1}}\mu_{1}, \quad \mu'_{2} = z^{r_{2}}\mu_{2}, \quad \ldots \quad \mu'_{n} = z^{r_{n}}\mu_{n}, \]
und da (wegen der Eigenschaft I., auf $z'$, $\mu'$ angewandt)
\[ z'\mu'_{s} = 0 \quad \text{für} \quad z' = 0, \]
so folgt
\[ z^{r_{s}-1}\mu_{s} = 0 \quad \text{für} \quad z = \infty \]
w.\ z.\ b.\ w.

Da aber die Functionen $z^{h}\mu_{s}$ linear unabhängig sind (wegen der rationalen Unabhängigkeit der Functionen $\mu_{s}$), so ergiebt sich hieraus nach §~24, (2.) der \emph{Hauptsatz:}

\emph{Die Schaar der Differentiale erster Gattung ist von der Dimension}
\[ (r_{2}-1) + (r_{3}-1) + \cdots + (r_{n}-1) = p, \]
\emph{und demnach ist auch $p$ die Dimension der Classe $W$ der vollständigen Polygone erster Gattung.}

Als Basis der Schaar der Differentialquotienten erster Gattung nach $z$ kann man die $p$ Functionen $z^{h}\mu_{s}$ $(h \leqq r_{s}-2)$ wählen, und die Grundpolygone $\mathfrak{W}_{1}, \mathfrak{W}_{2}, \ldots \mathfrak{W}_{p}$ der zugehörigen Differentiale $dw$ bilden eine Basis der Classe $W$.

3. Wegen einer späteren Anwendung soll hier noch eine besondere Art von Differentialquotienten erster Gattung $u'$ betrachtet werden, nämlich die, bei welchen die Bedingung II. ersetzt ist durch die dieselbe einschliessende Bedingung.

III. In jedem Punkte $\mathfrak{P}$, in welchem $z$ unendlich ist, sei
\[ (z^{k}u')_{0} = 0, \]
wo $k$ eine gegebene positive ganze Zahl.

Die Functionen $u'$ lassen sich darstellen durch
\[ u' = \frac{\mathfrak{U}^{k+1}\mathfrak{W}'}{\mathfrak{Z}} \]
und bilden ebenfalls eine Schaar; desgleichen bilden die Polygone $\mathfrak{W}'$ eine Classe $W'$, deren Ordnung ist
\[ \mathrm{w}-n(k+1) = 2p-2-n(k-1). \]
Die Polygone $\mathfrak{W}'$ sind jedoch von der Wahl der Variablen $z$ \emph{nicht} unabhängig.

\origpage{273}
Die Dimension der Classe $W'$ lässt sich nach derselben Methode bestimmen, wie die der Classe $W$. Da nämlich die Bedingung I. erfüllt ist, so sind die Functionen $u'$ gleichfalls in der Form (1.) enthalten; jedoch muss jetzt
\[ \frac{y_{s}}{z^{r_{s}-k}} = S\left(u'z^{k}\frac{\lambda_{s}}{z^{r_{s}}}\right) \]
für $z = \infty$ verschwinden, und daher kann der Grad der ganzen rationalen Function $y_{s}$ die Zahl $r_{s}-k-1$ nicht übersteigen. Es verschwindet also $y_{s}$ identisch, sobald $r_{s} < k+1$; andernfalls ist
\[ y_{s} = c_{0} + c_{1}z + \cdots + c_{r_{s}-k-1}z^{r_{s}-k-1}. \tag{3.} \]
Hat umgekehrt $y_{s}$ diese Form, so wird durch die Function
\[ u' = \overset{s}{\Sigma} y_{s}\mu_{s} \]
der Bedingung III. genügt, denn es hat, wie in 2. bewiesen,
\[ z^{k}(z^{r_{s}-k-1}\mu_{s}) = z^{r_{s}-1}\mu_{s} \]
für $z = \infty$ den Werth 0.

Daraus ergiebt sich, dass die Dimension der Schaar der Functionen $u'$ und folglich auch der Classe $W'$
\[ = \overset{\iota}{\Sigma}(r_{\iota}-k) \]
ist, wobei jedoch in der Summe nur diejenigen Glieder beizubehalten sind, die einen positiven Werth haben. Sind alle $r_{\iota}-k \leqq 0$, so existiren die gesuchten Functionen überhaupt nicht.

\subsection*{§~27. Polygonclassen erster und zweiter Gattung.}

Ist $\mathfrak{A}$ ein Polygon erster Gattung, so sind alle mit $\mathfrak{A}$ äquivalenten Polygone gleichfalls von der ersten Gattung. Denn wenn $\mathfrak{A}$ und $\mathfrak{B}$ Ergänzungspolygone sind und
\[ \mathfrak{A}\mathfrak{B} = \mathfrak{W}, \]
so ist, wenn $A$, $B$ die Classen von $\mathfrak{A}$ und $\mathfrak{B}$ sind:
\[ AB = W, \]
und, wenn $\mathfrak{A}'$ mit $\mathfrak{A}$ äquivalent ist, auch $\mathfrak{A}'\mathfrak{B} = \mathfrak{W}'$ äquivalent mit $\mathfrak{W}$ (§~18, 5.).

Wir nennen daher solche Classen, welche Polygone erster Gattung enthalten, \emph{Polygonclassen erster Gattung}, die übrigen \emph{Polygonclassen zweiter}

\origpage{274}
\emph{Gattung}. Die Classe $W$ der vollständigen Polygone erster Gattung heisst die \emph{Hauptclasse}, und zwei Classen $A$, $B$, die der Bedingung genügen
\[ AB = W, \]
\emph{Ergänzungsclassen}.

Ist
\[ \eta = \frac{\mathfrak{A}'}{\mathfrak{A}} \]
eine Function in $\Omega$, und $\mathfrak{A}'$ relativ prim zu $\mathfrak{A}$, also die Classe $A$ von $\mathfrak{A}$ eine eigentliche, so nennen wir $\eta$ eine Function \emph{erster} oder \emph{zweiter Gattung}, je nachdem die Classe $A$ von der ersten oder von der zweiten Gattung ist.

Ist $A$ eine beliebige Classe erster Gattung und $q$ die Anzahl der von einander unabhängigen Polygone $\mathfrak{W}$, die durch irgend ein Polygon $\mathfrak{A}$ der Classe $A$ theilbar sind, so ist nach §~21, 2.
\[ q = (A, W) = (O, B) \]
d.\ h.\ gleich der Dimension der Ergänzungsclasse $B$ von $A$. Ebenso ist $(B, W)$ gleich der Dimension der Classe $A$. Ist $A$ eine Classe zweiter Gattung, so ist $(A, W) = 0$. Da $p$ die Dimension von $W$ ist, so ist nach §~20, 2., 3. jede Classe, deren Ordnung $\leqq p-1$ ist, von der ersten Gattung, und es giebt ins Besondere Classen $A$ von der Ordnung $p-k$ der Art, dass $(A, W) = (O, B) = k$ ist. Aus den gleichen Sätzen folgt, dass es Classen von der Ordnung $p$ giebt, welche von der zweiten Gattung sind.

\subsection*{§~28. Der Riemann-Rochsche Satz für eigentliche Classen.}

Der \emph{Riemann-Roch}sche Satz, der nach seiner gewöhnlichen Ausdrucksweise die Anzahl der willkürlichen Constanten kennen lehrt, welche eine Function enthält, die in einer gewissen Anzahl gegebener Punkte unendlich wird, enthält nach unserer Darstellungsweise eine Beziehung zwischen der Dimension und der Ordnung einer Classe, resp.\ einer Classe und ihrer Ergänzungsclasse. Indem wir uns zunächst auf \emph{eigentliche} Classen beschränken, schicken wir der Ableitung dieser fundamentalen Relation die folgenden Bemerkungen voraus.

1. In einer eigentlichen Classe $A$ kann man nach §~19, 2. stets zwei zu einander relativ prime Polygone $\mathfrak{A}$, $\mathfrak{A}'$ auswählen (eines derselben kann in der Classe beliebig angenommen werden). Setzt man also
\[ z = \frac{\mathfrak{A}'}{\mathfrak{A}} \]

\origpage{275}
und, wenn $\mathfrak{A}''$ ein beliebiges drittes Polygon der Classe $A$ bedeutet:
\[ \omega = \frac{\mathfrak{A}''}{\mathfrak{A}}, \quad \frac{\omega}{z} = \frac{\mathfrak{A}''}{\mathfrak{A}'}, \]
so ist nach §~17 $\omega$ eine \emph{ganze} Function von $z$, $\frac{\omega}{z}$ eine ganze Function von $\frac{1}{z}$. Es ist daher (§~22) der Exponent von $\omega \leqq 1$.

Ist umgekehrt $\omega$ eine ganze Function von $z$, deren Exponent $\leqq 1$ ist, so hat es die Form
\[ \omega = \frac{\mathfrak{A}''}{\mathfrak{A}}, \]
wo $\mathfrak{A}''$ ein Polygon der Classe $A$ ist. Wenn nämlich
\[ \omega = \frac{\mathfrak{A}''_{1}}{\mathfrak{A}_{1}}, \quad \frac{\omega}{z} = \frac{\mathfrak{A}''_{1}\mathfrak{A}}{\mathfrak{A}_{1}\mathfrak{A}'} \]
und $\mathfrak{A}''_{1}$ relativ prim zu $\mathfrak{A}_{1}$ angenommen wird, so kann zunächst, da $\omega$ eine ganze Function von $z$ sein soll, $\mathfrak{A}_{1}$ keinen Punkt enthalten, der nicht auch in $\mathfrak{A}$ enthalten wäre. Es kann aber auch $\mathfrak{A}_{1}$ keinen Punkt öfter als $\mathfrak{A}$ enthalten, weil sonst $\frac{\omega}{z}$ in einem solchen Punkt (der nicht in $\mathfrak{A}'$ vorkommen kann) unendlich, also keine ganze Function von $\frac{1}{z}$ wäre. Daher ist $\mathfrak{A}$ theilbar durch $\mathfrak{A}_{1}$, und $\omega$ kann in die Form $\frac{\mathfrak{A}''}{\mathfrak{A}}$ gesetzt werden.

2. Um also die Gesammtheit der Polygone der Classe $A$ zu erhalten, haben wir nur diejenigen ganzen Functionen von $z$ aufzusuchen, deren Exponent $\leqq 1$ ist.

Ist $n$ die Ordnung der Classe $A$, also auch die Ordnung der Variablen $z$, und bilden $\lambda_{1}, \lambda_{2}, \ldots \lambda_{n}$ eine Normalbasis von $\mathfrak{o}$ mit den Exponenten $r_{1}, r_{2}, \ldots r_{n}$, darunter $r_{s}$ der letzte, welcher $\leqq 1$ ist, so kann jede Function $\omega$, deren Exponent $\leqq 1$ ist, nach §~22, 2. in der Form dargestellt werden
\[ \omega = c_{1}\lambda_{1} + c_{2}\lambda_{2} + \cdots + c_{s}\lambda_{s} + z\omega_{1}. \]
Da der Exponent von $z\omega_{1}$ aber nicht grösser als 1 sein kann, so muss $\omega_{1}$ eine Constante sein, und daher
\[ \omega = c_{1}\lambda_{1} + c_{2}\lambda_{2} + \cdots + c_{s}\lambda_{s} + c_{s+1}z. \]
Umgekehrt genügt jede Function von dieser Form der gestellten Forderung. Es ist also $s+1$ die Dimension der Classe $A$, welche hiernach, in Uebereinstimmung mit §~21, 1., stets $\leqq n+1$ ist. Die obere Grenze $n+1$ kann aber nur in dem Falle $p = 0$ erreicht werden und wird auch wirklich erreicht, weil in diesem Falle $r_{2}, r_{3}, \ldots r_{n} = 1$ sind. Daraus ergiebt sich,

\origpage{276}
dass ein \emph{einzelner Punkt} $\mathfrak{P}$ \emph{nur}, falls $p = 0$ ist, zu einer eigentlichen Classe gehören kann.

3. Wenn von den Exponenten $r_{s+1}, r_{s+2}, \ldots r_{n}$ einer grösser als 2 ist, so ist sicher auch $r_{n} > 2$, und es sind nach §~26, 2., wenn $\mathfrak{Z}$ das Verzweigungspolygon in $z$ bedeutet,
\[ \mu_{n} = \frac{\mathfrak{A}^{2}\mathfrak{W}}{\mathfrak{Z}}, \quad \mu_{n}z = \frac{\mathfrak{A}^{2}\mathfrak{W}_{1}}{\mathfrak{Z}} = \frac{\mathfrak{A}\mathfrak{A}'\mathfrak{W}}{\mathfrak{Z}} \]
Differentialquotienten erster Gattung nach $z$, also
\[ \mathfrak{A}\mathfrak{W}_{1} = \mathfrak{A}'\mathfrak{W} \]
oder, da $\mathfrak{A}$, $\mathfrak{A}'$ relativ prim sind,
\[ \mathfrak{W} = \mathfrak{A}\mathfrak{B}, \quad \mathfrak{W}_{1} = \mathfrak{A}'\mathfrak{B}, \]
d.\ h.\ die Classe $A$ ist von der \emph{ersten} Gattung ($z$ eine Variable erster Gattung). Machen wir daher zunächst die Annahme, es sei $A$ eine Classe \emph{zweiter} Gattung, so folgt
\[ r_{s+1} = 2, \quad r_{s+2} = 2, \quad \ldots \quad r_{n} = 2 \]
und
\[ p = (r_{2}-1) + \cdots + (r_{s}-1) + (r_{s+1}-1) + \cdots + (r_{n}-1) = n-s. \]
Die Dimension $s+1$ der Classe $A$ ist daher
\[ (O, A) = n-p+1. \]

4. Machen wir zweitens die Annahme, es sei $A$ von der \emph{ersten} Gattung und wie in §~27
\[ q = (A, W), \]
so existiren $q$ linear unabhängige, durch $\mathfrak{A}$ theilbare vollständige Polygone erster Gattung, und die diesen entsprechenden Differentialquotienten erster Gattung nach $z$, deren es ebenfalls $q$ und nicht mehr linear unabhängige giebt, haben den Ausdruck
\[ v = \frac{\mathfrak{A}^{3}\mathfrak{B}}{\mathfrak{Z}}, \]
worin $\mathfrak{B}$ ein Polygon von $2p-2-n$ Punkten bedeutet; die Classe $B$ von $\mathfrak{B}$ ist die Ergänzungsclasse von $A$, und daher ihre Dimension gleich $q$ (§~27).

Diese Functionen $v$ haben aber die Eigenschaft, dass in den Eckpunkten von $\mathfrak{A}$, d.\ h.\ für $z = \infty$ nicht nur $zv$, sondern auch
\[ z^{2}v = \frac{\mathfrak{A}\mathfrak{A}'^{2}\mathfrak{B}}{\mathfrak{Z}} \]
verschwindet, und sind hierdurch und durch die Forderung, Differentialquotienten

\origpage{277}
erster Gattung zu sein, völlig bestimmt. Denn ist
\[ v = \frac{\mathfrak{A}^{2}\mathfrak{W}}{\mathfrak{Z}}, \quad vz^{2} = \frac{\mathfrak{A}'^{2}\mathfrak{W}}{\mathfrak{Z}}, \]
so muss, wenn $z^{2}v$ in allen Punkten von $\mathfrak{A}$ verschwinden soll, $\mathfrak{W}$ durch $\mathfrak{A}$ theilbar sein, da $\mathfrak{A}'$ relativ prim zu $\mathfrak{A}$ vorausgesetzt ist. Es ist daher nach §~26, 3:
\[ q = (r_{s+1}-2)+(r_{s+2}-2)+\cdots+(r_{n}-2), \]
andererseits
\[ p = (r_{s+1}-1)+(r_{s+2}-1)+\cdots+(r_{n}-1), \]
folglich:
\[ p-q = n-s, \quad s = n-p+q. \]
Hierin ist der \emph{Riemann-Roch}sche \emph{Satz} enthalten, dem wir, mit Rücksicht auf §~27, für diesen Fall folgenden Ausdruck geben können: \emph{Sind $A$, $B$ Ergänzungsclassen erster Gattung, von denen wenigstens die eine eine eigentliche ist, und $a$, $b$ ihre Ordnungen, also}
\[ a+b = 2p-2, \]
\emph{so ist}
\[ (O, A)-\frac{1}{2}a = (O, B)-\frac{1}{2}b. \]

5. Wir können, wenn wir den Fall $(A, W) = 0$ nicht ausschliessen, den \emph{Riemann-Roch}schen Satz für beide Fälle dahin zusammenfassen:

\emph{Ist $A$ eine eigentliche Classe von der Ordnung $n$, so ist ihre Dimension}
\[ (O, A) = n-p+1+(A, W). \]
Da die Dimension einer eigentlichen Classe (wenn sie nicht aus dem einzigen Nulleck besteht) mindestens $= 2$ sein muss, so folgt noch, wenn $(A, W) = 0$ ist,
\[ n \geqq p+1, \]
und daraus der von \emph{Riemann} herrührende Satz:

\emph{Jede Function, deren Ordnung $\leqq p$ ist, ist eine Function erster Gattung.}

6. Es lässt sich mit Hülfe dieser Sätze leicht beweisen, \emph{dass die Hauptclasse $W$ der vollständigen Polygone erster Gattung stets eine eigentliche ist.}

Ist nämlich $\mathfrak{M}$ der Theiler von $W$, so lässt sich nach §~19, 2. in $W$ ein Polygon $\mathfrak{A}\mathfrak{M}$ der Art finden, dass $\mathfrak{A}$ relativ prim zu $\mathfrak{M}$ ist. Die Classe $A$ von $\mathfrak{A}$ ist eine eigentliche (§~21, 3.), und zugleich ist $\mathfrak{A}\mathfrak{M}$ das einzige durch $\mathfrak{A}$ theilbare Polygon der Classe $W$ (weil jedes Polygon in $W$ den Theiler $\mathfrak{M}$ hat). Also ist
\[ (A, W) = 1. \]

\origpage{278}
Nun ist $p$ die Dimension von $W$, also auch die von $A$, und mithin nach dem \emph{Riemann-Roch}schen Satze die Ordnung von $A$ gleich $2p-2$, d.\ h.\ ebenso gross wie die von $W$. Mithin ist $\mathfrak{M} = \mathfrak{O}$.

\subsection*{§~29. Der Riemann-Rochsche Satz für uneigentliche Classen erster Gattung.}

Ist $A$ eine Classe erster Gattung vom Theiler $\mathfrak{M}$ und
\[ A = \mathfrak{M}A', \]
so ist $A'$ eine eigentliche Classe erster Gattung. Es sei $B$ die Ergänzungsclasse von $A$; $B'$ die von $A'$; $a$, $b$ die Ordnungen der Classen $A$, $B$; $m$ die Ordnung von $\mathfrak{M}$. Die gesammte Classe $B$ erhält man, wenn man in sämmtlichen durch $\mathfrak{M}$ theilbaren Polygonen der Classe $B'$ den Factor $\mathfrak{M}$ unterdrückt; denn ist
\[ \mathfrak{A}\mathfrak{B} = \mathfrak{A}'\mathfrak{M}\mathfrak{B} = \mathfrak{W}, \]
so gehört $\mathfrak{M}\mathfrak{B}$ in die Classe $B'$, und umgekehrt, wenn
\[ \mathfrak{A}'\mathfrak{B}' = \mathfrak{A}'\mathfrak{M}\mathfrak{B} = \mathfrak{W} \]
ist, so gehört $\mathfrak{B}$ in die Classe $B$.

Hieraus ergiebt sich aber nach §~21, 2.
\[ (O, B) \geqq (O, B')-m. \]
Nun ist $A'$ eine eigentliche Classe von derselben Dimension wie $A$ und von der Ordnung $a-m$, also (§~28, 5.)
\[ (O, A) = (O, A') = a-m-p+1+(A', W), \]
oder
\[ (O, A) = (O, B')-m+a-p+1; \]
daher
\[ (O, A) \leqq (O, B)+a-p+1 = (O, B)+\frac{1}{2}(a-b), \]
also
\[ (O, A)-\frac{1}{2}a \leqq (O, B)-\frac{1}{2}b. \]
Da aber die Classen $A$, $B$ mit einander vertauscht werden können, so folgt in gleicher Weise
\[ (O, B)-\frac{1}{2}b \leqq (O, A)-\frac{1}{2}a, \]
d.\ h.
\[ (O, A)-\frac{1}{2}a = (O, B)-\frac{1}{2}b, \]

\origpage{279}
\emph{wodurch der Riemann-Rochsche Satz in derselben Form wie in §~28, 4. für Polygonclassen erster Gattung allgemein nachgewiesen ist}${}^{*)}$.

\subsection*{§~30. Uneigentliche Classen zweiter Gattung.}

Es soll nun die Bedingung aufgesucht werden, unter der eine Polygonclasse zweiter Gattung $A$ von der Ordnung $n$ überhaupt eine uneigentliche sein kann, wobei sich die allgemeine Gültigkeit des \emph{Riemann-Roch}schen Satzes von selbst ergeben wird.

1. Jede Classe $A$ kann stets durch Multiplication mit einer andern Classe $N$ von der Ordnung $\nu$ in eine eigentliche Classe $AN$ verwandelt werden. Denn ist $\mathfrak{A}$ ein beliebiges Polygon in $A$, so wähle man eine Variable $z$, welche in sämmtlichen Punkten von $\mathfrak{A}$ endlich bleibt (§~15, 6.). Ist dann $\eta$ eine beliebige Function des durch $\mathfrak{A}$ erzeugten Ideals in $z$, so ist das Obereck von $\eta$ durch $\mathfrak{A}$ theilbar, also von der Form $\mathfrak{A}\mathfrak{N}$, und die Classe von $\mathfrak{A}\mathfrak{N}$ ist eine eigentliche.

2. Die Dimension der eigentlichen Classe $AN$ zweiter Gattung ist nach §~28, 3.
\[ (O, AN) = n+\nu-p+1, \]
und hieraus folgt nach §~21, 2.
\[ (O, A) \geqq n-p+1. \]
Ist nun der Theiler $\mathfrak{M}$ der Classe $A$ von der Ordnung $m$, und
\[ A = \mathfrak{M}A', \]
so ist $A'$ eine eigentliche Classe von derselben Dimension wie $A$, und mithin (§~28, 5.)
\[ (O, A) = (O, A') = n-m-p+1+(A', W), \]
also
\[ (A', W) \geqq m, \]
d.\ h.\ $A'$ muss gewiss von der ersten Gattung sein, wenn $A$ eine uneigentliche Classe ist. Ist also $B'$ die Ergänzungsclasse von $A'$, so ist auch
\[ (O, B') \geqq m. \]

\textbf{*)} Nach der Ausdrucksweise von \emph{Christoffel} (Ueber die canonische Form der \emph{Riemann}schen Integrale erster Gattung, Annali di Matematica pura ed applicata, Serie II, Tomo IX) ist
\[ (A, W)+a-p = (O, B)+a-p = (O, A)-1 \]
der „Ueberschuss“,
\[ (A, W)-1 = (O, B)-1 \]
der „Defect“ des Punktsystems $\mathfrak{A}$.

\origpage{280}
Wäre aber $(O, B') > m$, so würde sich nach §~20, 2. in $B'$ ein durch $\mathfrak{M}$ theilbares Polygon $\mathfrak{M}\mathfrak{B}$ finden lassen und es wäre
\[ \mathfrak{A}'\mathfrak{M}\mathfrak{B} = \mathfrak{A}\mathfrak{B} = \mathfrak{W}, \]
also $A$ von der ersten Gattung, gegen die Voraussetzung. Es ist also
\[ (A', W) = m \]
und folglich
\[ (O, A) = n-p+1, \]
worin wieder der \emph{Riemann-Roch}sche Satz für diesen Fall, genau in der Form von §~28, 3. enthalten ist.

3. Enthält die Classe $A$ nur ein einziges isolirtes Polygon, so ist $(O, A) = n-p+1 = 1$, mithin $n = p$, d.\ h.\ ein isolirtes Polygon zweiter Gattung hat stets die Ordnung $p$. Umgekehrt ist nach 2. jedes Polygon zweiter Gattung von der Ordnung $p$ ein isolirtes.

4. Unter Beibehaltung der Bezeichnung von 2. ist $(O, B') = m$, und daher lässt sich nach dem oft angewandten Satze (§~20, 2.) in $B'$ ein durch ein beliebiges $(m-1)$-Eck theilbares Polygon finden. Setzt man also, indem man einen beliebigen Punkt $\mathfrak{P}$ von $\mathfrak{M}$ absondert,
\[ \mathfrak{M} = \mathfrak{P}\mathfrak{M}', \]
so ist ein Polygon $\mathfrak{M}'\mathfrak{B}$ in $B'$ enthalten und also
\[ \mathfrak{A}'\mathfrak{M}'\mathfrak{B} = \mathfrak{W}. \]
Das Polygon $\mathfrak{A}'\mathfrak{M}' = \mathfrak{A}''$ und seine Classe $A''$ sind daher von der \emph{ersten Gattung}, und $A$ hat, wenn $P$ die Classe von $\mathfrak{P}$ bedeutet, die Form
\[ A = PA''. \]
Zugleich muss $(A'', W) = (O, B'') = 1$ sein, d.\ h.\ die Ergänzungsclasse $B''$ von $A''$ enthält nur ein einziges isolirtes Polygon $\mathfrak{B}''$, da sonst in $B''$ ein durch $\mathfrak{P}$ theilbares Polygon existiren würde, und also auch $A$ gegen die Voraussetzung von der ersten Gattung wäre.

5. Ist umgekehrt $A''$ eine Classe erster Gattung, für welche $(A'', W) = 1$, so dass die Ergänzungsclasse $B''$ von $A''$ aus einem isolirten Polygon $\mathfrak{B}''$ besteht; ist ferner $\mathfrak{P}$ ein in $\mathfrak{B}''$ nicht aufgehender Punkt, und seine Classe $P$, so ist $A = PA''$ eine uneigentliche Classe zweiter Gattung von der Ordnung $n$, in deren Theiler $\mathfrak{P}$ aufgeht.

Dass $A$ von der zweiten Gattung ist, ergiebt sich zunächst aus der Annahme, dass $\mathfrak{P}$ in $\mathfrak{B}''$ nicht aufgeht. Die Dimension von $A$ ist daher nach 2.
\[ (O, A) = n-p+1, \]

\origpage{281}
wo $n$ die Ordnung von $A$ bedeutet; andererseits ist die Dimension der Classe $A''$ nach §§~28 und 29:
\[ (O, A'') = n-p+(A'', W) = n-p+1; \]
also sind $A$ und $A''$ von derselben Dimension. Sämmtliche Polygone der Classe $A''$ gehen aber durch Multiplication mit $\mathfrak{P}$ in Polygone der Classe $A$ über, und wegen der Gleichheit der Dimensionen wird hierdurch auch die letzte Classe vollständig erschöpft. Es enthalten daher sämmtliche Polygone der Classe $A$ den Factor $\mathfrak{P}$, der sonach auch im Theiler von $A$ aufgeht.

6. In dem besonderen Fall, wo das Geschlecht $p$ des Körpers $\Omega$ den Werth 0 hat, kommen Polygone und Classen erster Gattung überhaupt nicht vor. Es existiren also in diesem Falle auch keine uneigentlichen Classen. Die Dimension einer jeden Classe ist um 1 grösser als ihre Ordnung. Insbesondere gehört also auch jeder Punkt $\mathfrak{P}$ zu einer eigentlichen Classe von der Dimension 2, und daher existiren in diesem Fall in $\Omega$ Functionen $z$, welche von der ersten Ordnung sind. Durch eine solche lässt sich jede andere Function des Körpers \emph{rational} ausdrücken, denn die zwischen $z$ und einer andern Variabeln des Körpers bestehende irreductible Gleichung ist in Bezug auf letztere vom ersten Grad (§~15, 7.).

\subsection*{§~31. Die Differentiale zweiter und dritter Gattung.}

1. Ist jetzt nach der in §~25 eingeführten Bezeichnung
\[ d\varpi = \frac{\mathfrak{A}}{\mathfrak{B}} \]
ein beliebiges Differential in $\Omega$, also, wenn $a$, $b$ die Ordnungen von $\mathfrak{A}$ und $\mathfrak{B}$ sind,
\[ a = b+2p-2, \]
und werden $\mathfrak{A}$, $\mathfrak{B}$ als relativ prim vorausgesetzt, so muss, wenn $\mathfrak{U}$, $\mathfrak{Z}$ Untereck und Verzweigungspolygon für eine beliebige Variable $z$ bedeuten, $\mathfrak{U}^{2}\mathfrak{A}$ mit $\mathfrak{Z}\mathfrak{B}$ äquivalent sein (§~25). Bezeichnet man also mit $U$, $Z$, $A$, $B$ die Classen der Polygone $\mathfrak{U}$, $\mathfrak{Z}$, $\mathfrak{A}$, $\mathfrak{B}$, so muss
\[ U^{2}A = ZB \]
sein. Andererseits ist aber, wenn $W$ die Hauptclasse erster Gattung ist,
\[ U^{2}W = Z, \]

\origpage{282}
woraus sich die Relation
\[ A = BW \]
ergiebt. Ist umgekehrt $\mathfrak{A}$ ein beliebiges Polygon der Classe $BW$, so folgt daraus die Aequivalenz von $\mathfrak{U}^{2}\mathfrak{A}$ mit $\mathfrak{Z}\mathfrak{B}$, also die Existenz eines Differentials von der Bezeichnung $\frac{\mathfrak{A}}{\mathfrak{B}}$. Daraus ergiebt sich, dass $\mathfrak{B}$ dann und nur dann Untereck eines Differentials $d\varpi$ sein kann, wenn in $BW$ ein zu $\mathfrak{B}$ relativ primes Polygon existirt, d.\ h.\ wenn der Theiler der Classe $BW$ relativ prim zu $\mathfrak{B}$ ist. Die Dimension der Classe $BW$ giebt dann zugleich die Dimension der zum Untereck $\mathfrak{B}$ gehörigen Schaar von Differentialen $d\varpi$ (§~25). Die Sätze §~30, 4., 5. ergeben daher, da $(W, W) = 1$ ist, das folgende Resultat.

\emph{a}) Besteht $\mathfrak{B}$ aus einem einzigen Punkt (ist $b = 1$), so ist die Classe $BW$ eine uneigentliche mit dem Theiler $\mathfrak{B}$; \emph{also kann die Ordnung $b$ des Unterecks eines Differentials $d\varpi$ nicht gleich Eins sein.}

\emph{b}) Ist $b \geqq 2$, so ist $BW$ stets eine eigentliche Classe zweiter Gattung und daher ihre Dimension
\[ b+p-1. \]
\emph{Untereck eines Differentials kann also jedes beliebige Polygon von mehr als einem Punkt sein, und es existiren unter den zu einem Untereck von der Ordnung $b$ gehörigen Differentialen $b+p-1$ linear unabhängige.}

2. Wir suchen jetzt unter der Voraussetzung, dass $b \geqq 2$ ist, für die Classe $A$ eine Basis der Art auf, dass jedes Element $\mathfrak{A}_{r}$ dieser Basis ein Differential $d\varpi_{r}$ von möglichst einfacher Beschaffenheit liefert, nämlich ein solches, dessen Untereck eine Potenz eines einzelnen Punktes oder das Product aus nur zwei verschiedenen Punkten ist.

Angenommen, es sei für die Classe $BW$ eine solche Basis bereits gefunden
\[ \mathfrak{A}_{1}, \mathfrak{A}_{2}, \mathfrak{A}_{3}, \ldots \mathfrak{A}_{b+p-1}, \tag{1.} \]
so bilden wir daraus, wenn $P$ die Classe eines beliebigen Punktes $\mathfrak{P}$ bedeutet, eine ebensolche Basis für die Classe $BPW$ von der Dimension $b+p$, nämlich
\[ \mathfrak{P}\mathfrak{A}_{1}, \mathfrak{P}\mathfrak{A}_{2}, \ldots \mathfrak{P}\mathfrak{A}_{b+p-1}, \mathfrak{A}'. \tag{2.} \]
Die ersten $b+p-1$ dieser Polygone gehören wirklich der Classe $BPW$ an und sind von einander unabhängig, weil es die Polygone (1.) sind; zugleich

\origpage{283}
sind die aus ihnen gebildeten Differentiale
\[ d\varpi_{r} = \frac{\mathfrak{P}\mathfrak{A}_{r}}{\mathfrak{P}\mathfrak{B}} = \frac{\mathfrak{A}_{r}}{\mathfrak{B}} \]
mit den aus (1.) gebildeten identisch. Es kommt also nur noch auf die Bildung von $\mathfrak{A}'$ an, wobei zwei Fälle zu unterscheiden sind.

\emph{a}) Geht $\mathfrak{P}$ in $\mathfrak{B}$ auf und ist $\mathfrak{B} = \mathfrak{M}\mathfrak{P}^{m}$, $\mathfrak{M}$ nicht durch $\mathfrak{P}$ theilbar, so ist $P^{m+1}W$ eine eigentliche Classe (weil $m+1 \geqq 2$, §~30, 4.), in welcher folglich ein durch $\mathfrak{P}$ nicht theilbares Polygon $\mathfrak{N}$ existirt; setzt man nun $\mathfrak{A}' = \mathfrak{M}\mathfrak{N}$, so gehört $\mathfrak{A}'$ der Classe $BPW$ an und ist durch $\mathfrak{P}$ nicht theilbar, folglich auch nicht in der Schaar $(\mathfrak{P}\mathfrak{A}_{1}, \mathfrak{P}\mathfrak{A}_{2}, \ldots \mathfrak{P}\mathfrak{A}_{b+p-1})$, deren Theiler $\mathfrak{P}$ ist, enthalten; mithin sind die Polygone (2.) unabhängig von einander, und da ihre Anzahl $b+p$ ist, so bilden sie eine Basis der Classe $BPW$. Das aus $\mathfrak{A}'$ gebildete Differential
\[ d\varpi' = \frac{\mathfrak{A}'}{\mathfrak{P}\mathfrak{B}} = \frac{\mathfrak{N}}{\mathfrak{P}^{m+1}} \]
hat die geforderte Form, da sein Untereck eine Potenz eines einzelnen Punktes ist.

\emph{b}) Geht $\mathfrak{P}$ nicht in $\mathfrak{B}$ auf, so wähle man ein für allemal einen in $\mathfrak{B}$ aufgehenden Punkt $\mathfrak{P}_{1}$ und setze $\mathfrak{B} = \mathfrak{M}\mathfrak{P}_{1}$ (gleichgültig ob $\mathfrak{M}$ durch $\mathfrak{P}_{1}$ theilbar ist oder nicht). Man wähle sodann in der eigentlichen Classe $PP_{1}W$ ein durch $\mathfrak{P}$ und $\mathfrak{P}_{1}$ nicht theilbares Polygon $\mathfrak{N}$, so gehört $\mathfrak{A}' = \mathfrak{M}\mathfrak{N}$ wieder in die Classe $PBW$, und da $\mathfrak{A}'$ nicht durch $\mathfrak{P}$ theilbar ist, so folgt wie oben, dass die Polygone (2.) eine Basis von $BPW$ bilden. Zugleich ist
\[ d\varpi' = \frac{\mathfrak{A}'}{\mathfrak{P}\mathfrak{B}} = \frac{\mathfrak{N}}{\mathfrak{P}\mathfrak{P}_{1}}, \]
also von der verlangten Form.

Es bleibt noch übrig, den Anfang dieser Operation zu beschreiben. Ist $b = 0$, also $\mathfrak{B} = \mathfrak{O}$, so ist
\[ BW = W = (\mathfrak{W}_{1}, \mathfrak{W}_{2}, \ldots \mathfrak{W}_{p}) \]
(die Hauptclasse erster Gattung).

Ist $b = 2$, so wähle man aus der eigentlichen Classe $BW$ ein Polygon $\mathfrak{N}$, welches relativ prim zu $\mathfrak{B}$ ist; dann ist
\[ BW = (\mathfrak{B}\mathfrak{W}_{1}, \mathfrak{B}\mathfrak{W}_{2}, \ldots \mathfrak{B}\mathfrak{W}_{p}, \mathfrak{N}). \]
Geht man von dieser Basis aus, um in der oben beschriebenen Weise eine Basis (1.) zu bestimmen, die dem beliebig gegebenen Polygon
\[ \mathfrak{B} = \mathfrak{P}_{1}^{m_{1}}\mathfrak{P}_{2}^{m_{2}}\mathfrak{P}_{3}^{m_{3}}\ldots \]

\origpage{284}
entspricht, und bestimmt die beiden Polygone $\mathfrak{A}'_{r}$, $\mathfrak{B}'_{r}$ aus der Bedingung
\[ d\varpi_{r} = \frac{\mathfrak{A}_{r}}{\mathfrak{B}} = \frac{\mathfrak{A}'_{r}}{\mathfrak{B}'_{r}}, \]
so dass sie keinen gemeinschaftlichen Theiler haben, so sind die Polygone $\mathfrak{B}'_{r}$, die als Unterecke der Differentiale $d\varpi_{r}$ auftreten, folgende:

\emph{a}) $p$-mal tritt der Nenner $\mathfrak{O}$ auf, und die zugehörigen Differentiale $d\varpi_{r}$ sind die \emph{Differentiale erster Gattung}.

\emph{b}) Je einmal treten die Unterecke $\mathfrak{P}_{1}^{2}, \mathfrak{P}_{1}^{3}, \ldots \mathfrak{P}_{1}^{m_{1}}$ (wenn $m_{1} \geqq 2$), $\mathfrak{P}_{2}^{2}, \mathfrak{P}_{2}^{3}, \ldots \mathfrak{P}_{2}^{m_{2}}$; $\mathfrak{P}_{3}^{2}, \mathfrak{P}_{3}^{3}, \ldots \mathfrak{P}_{3}^{m_{3}}$, \ldots\ auf.

Die zu den Unterecken $\mathfrak{P}^{r}$ gehörigen Differentiale $d\varpi_{r}$ werden, wenn eine genauere Unterscheidung nöthig ist, mit $dt_{(\mathfrak{P}^{r-1})}$ bezeichnet und heissen \emph{Differentiale zweiter Gattung}.

\emph{c}) Endlich treten die Producte $\mathfrak{P}_{1}\mathfrak{P}_{2}$, $\mathfrak{P}_{1}\mathfrak{P}_{3}$, \ldots\ (bei festgehaltenem $\mathfrak{P}_{1}$) je einmal auf. Die zugehörigen Differentiale $d\varpi_{r}$ werden mit $d\pi_{(\mathfrak{P}_{1}, \mathfrak{P}_{r})}$ bezeichnet und heissen \emph{Differentiale dritter Gattung}.

Jedes Differential $d\varpi$, dessen Untereck $\mathfrak{B}$ ist, kann in der Form dargestellt werden
\[ d\varpi = \overset{r}{\Sigma} c_{r}\,d\varpi_{r} \tag{3.} \]
mit constanten Coefficienten $c_{r}$, welche die \emph{Normalform} des Differentials $d\varpi$ genannt wird. Hat man jedes der einzelnen Differentiale $d\varpi_{r}$ auf eine bestimmte Art gewählt, so lässt sich die Normalform auch \emph{nur auf eine einzige Weise} herstellen, was unmittelbar aus der linearen Unabhängigkeit der Differentiale $d\varpi_{r}$ folgt.

\subsection*{§~32. Die Residuen.}

1. Ist $d\varpi$ ein beliebiges Differential in $\Omega$ und $\mathfrak{P}$ ein Punkt, der $m$-mal im Untereck $\mathfrak{B}$ desselben vorkommt ($m \geqq 0$), so wähle man eine Variable $z$ so, dass sie in $\mathfrak{P}$ $\infty^{1}$ wird. Es lässt sich dann (nach §~15, 4.), und zwar nur auf eine Weise, setzen
\[ \frac{d\varpi}{dz} = a_{m-2}z^{m-2}+a_{m-3}z^{m-3}+\cdots+a_{1}z+a_{0}+a_{-1}z^{-1}+\eta z^{-2}, \tag{1.} \]
worin die $a$ Constanten, $\eta$ eine Function in $\Omega$, die in $\mathfrak{P}$ endlich ist. Der Coefficient $-a_{-1}$ von $-z^{-1}$ in diesem Ausdruck heisst das \emph{Residuum des Differentials $d\varpi$ in Bezug auf den Punkt $\mathfrak{P}$}. Aus dieser Definition ergeben sich die folgenden Sätze:

\origpage{285}

2. Das Residuum in Bezug auf einen Punkt $\mathfrak{P}$ kann nur dann von Null verschieden sein, wenn $m > 0$, d.\ h.\ wenn der Punkt $\mathfrak{P}$ im Untereck von $d\varpi$ wirklich vorkommt, und ist daher für die Differentiale erster Gattung immer gleich 0.

3. Das Residuum einer Summe von Differentialen ist gleich der Summe der Residuen der einzelnen Differentiale.

4. Das Residuum eines \emph{eigentlichen} Differentials ist stets gleich 0. Ist nämlich $\sigma$ eine Function in $\Omega$, und wenn die $b$ Constanten, $\sigma'$ eine in $\mathfrak{P}$ endliche Function bedeuten,
\[ \sigma = b_{m}z^{m}+b_{m-1}z^{m-1}+\cdots+b_{1}z+\sigma', \]
so ergiebt sich durch Differentiation dieses Ausdruckes nach $z$, da $\frac{d\sigma'}{dz}$ in $\mathfrak{P}$ unendlich klein von mindestens zweiter Ordnung ist (§~23, 10.)\ednote{Printed "§ 23, 10.", but §~23 has only articles 1.–9. and formulas (1.)–(9.). The fact used is presumably §~23, 8.: the order of $\frac{d\alpha}{d\beta}$ is the difference of the orders of $\alpha-\alpha_{0}$ and $\beta-\beta_{0}$ (following formula (7.)). Here $z$ is $\infty^{1}$ in $\mathfrak{P}$ and $\sigma'$ is finite there, so $\frac{d\sigma'}{dz}$ has order at least $1-(-1) = 2$.}, dass in dem Ausdruck für $\frac{d\sigma}{dz}$ ein Glied mit $z^{-1}$ gar nicht vorkommt, womit die Behauptung erwiesen ist.

5. Das Residuum eines Differentials $d\varpi$ ist unabhängig von der Wahl der Veränderlichen $z$. Ist nämlich $z_{1}$ eine zweite Veränderliche von derselben Beschaffenheit wie $z$, also, wenn $a$ constant, $\zeta$ in $\mathfrak{P}$ endlich ist:
\[ z = az_{1}+\zeta, \tag{2.} \]
so ergiebt sich, wenn zur Abkürzung
\[ \alpha = \frac{a_{m-2}z^{m-1}}{m-1}+\frac{a_{m-3}z^{m-2}}{m-2}+\cdots+a_{0}z \]
gesetzt wird:
\[ \frac{d\varpi}{dz_{1}} = \frac{d\varpi}{dz}\frac{dz}{dz_{1}} = \frac{d\alpha}{dz_{1}}+a_{-1}z^{-1}\frac{dz}{dz_{1}}-\eta\frac{dz^{-1}}{dz_{1}}. \]
Nun ist, wenn $\zeta'$, $\zeta''$ in $\mathfrak{P}$ endliche Functionen sind, wie sich nach §~23 und §~15, 4. leicht ergiebt:
\[ z^{-1}\frac{dz}{dz_{1}} = z_{1}^{-1}+z_{1}^{-2}\zeta', \quad \frac{dz^{-1}}{dz_{1}} = z_{1}^{-2}\zeta'', \]
und daraus folgt nach 3., 4. die Richtigkeit der aufgestellten Behauptung${}^{*)}$.

6. \emph{Die Summe der Residuen eines jeden Differentials $d\varpi$ in Bezug auf alle Punkte $\mathfrak{P}$ ist stets gleich Null.}

\textbf{*)} Man kann bei der Definition des Residuums auch eine Veränderliche $r$ zu Grunde legen, die in $\mathfrak{P}$ unendlich \emph{klein} in der ersten Ordnung ist. Ist dann
\[ \frac{d\varpi}{dr} = a_{m}r^{-m}+\cdots+a_{1}r^{-1}+\eta \]
und $\eta$ in $\mathfrak{P}$ endlich, so ist $a_{1}$ das Residuum von $d\varpi$ in Bezug auf $\mathfrak{P}$.

\origpage{286}

Beim Beweise dieses wichtigen Satzes können wir uns auf die Betrachtung der Residuen beschränken, welche zu den sämmtlichen im Untereck $\mathfrak{B}$ von $d\varpi$ aufgehenden von einander verschiedenen Punkten gehören; wir fügen jedoch zu diesen noch so viele von einander verschiedene willkürliche Punkte mit verschwindenden Residuen hinzu, bis wir ein aus lauter einfachen Punkten bestehendes einer eigentlichen Classe angehöriges Polygon $\mathfrak{P}_{1}\mathfrak{P}_{2}\ldots\mathfrak{P}_{n}$ erhalten. Dann wählen wir eine Variable $z$ von der Ordnung $n$, deren Untereck eben dies Polygon ist, welche also in jedem der Punkte $\mathfrak{P}_{1}, \mathfrak{P}_{2}, \ldots \mathfrak{P}_{n}$ und nur in diesen $\infty^{1}$ wird. Unter diesen finden sich dann sämmtliche von einander verschiedene in $\mathfrak{B}$ aufgehende Punkte. Es ergiebt sich unter dieser Voraussetzung für $\iota = 1, 2, \ldots n$
\[ \frac{d\varpi}{dz} = a_{m-2}^{(\iota)}z^{m-2}+a_{m-3}^{(\iota)}z^{m-3}+\cdots+a_{0}^{(\iota)}+a_{-1}^{(\iota)}z^{-1}+\eta^{(\iota)}z^{-2}, \tag{3.} \]
wo $\eta^{(\iota)}$ eine in $\mathfrak{P}_{\iota}$ endliche Function bedeutet. Lassen wir für die Constanten $a^{(\iota)}$ auch den Werth 0 zu, so kann der Exponent $m$ unabhängig von $\iota$ angenommen werden ($m$ ist dann, wenn nicht alle $a_{m-2}^{(\iota)}$ verschwinden, der Exponent der höchsten Potenz eines einzelnen Punktes, welche in $\mathfrak{B}$ vorkommt). Der zu beweisende Satz besteht dann darin, dass $\overset{\iota}{\Sigma} a_{-1}^{(\iota)} = 0$ ist. Um ihn zu beweisen, bilden wir die Spur der Function $\frac{d\varpi}{dz}$ für die Variable $z$ (§~2) und bedienen uns dabei einer Erweiterung des Verfahrens §~16, 4. Wir wählen ein Functionensystem $\varrho_{1}, \varrho_{2}, \ldots \varrho_{n}$ in $\Omega$ folgendermassen: Es sei
\[ \begin{matrix}
\varrho_{1} = 0^{m} & \text{in}\ \mathfrak{P}_{2}, \mathfrak{P}_{3}, \ldots \mathfrak{P}_{n}, & \text{endlich} & \text{und} & \text{von} & \text{Null} & \text{verschieden} & \text{in} & \mathfrak{P}_{1}, \\
\varrho_{2} = 0^{m} & \text{in}\ \mathfrak{P}_{1}, \mathfrak{P}_{3}, \ldots \mathfrak{P}_{n}, & \text{''} & \text{''} & \text{''} & \text{''} & \text{''} & \text{''} & \mathfrak{P}_{2}, \\
\cdots & \cdots & \cdots & \cdots & \cdots & \cdots & \cdots & \cdots & \cdots \\
\varrho_{n} = 0^{m} & \text{in}\ \mathfrak{P}_{1}, \mathfrak{P}_{2}, \ldots \mathfrak{P}_{n-1}, & \text{''} & \text{''} & \text{''} & \text{''} & \text{''} & \text{''} & \mathfrak{P}_{n}.
\end{matrix} \]
Sind nun $x_{1}, x_{2}, \ldots x_{n}$ rationale Functionen von $z$, und ist
\[ \eta = x_{1}\varrho_{1}+x_{2}\varrho_{2}+\cdots+x_{n}\varrho_{n} \]
eine Function in $\Omega$, welche für $z = \infty$, d.\ h.\ in $\mathfrak{P}_{1}, \mathfrak{P}_{2}, \ldots \mathfrak{P}_{n}$ endlich ist, so müssen $x_{1}, x_{2}, \ldots x_{n}$ ebenfalls für $z = \infty$ endlich sein. Sind nämlich die $x_{1}, x_{2}, \ldots x_{n}$ für $z = \infty$ \emph{nicht} alle endlich, so existirt ein positiver Exponent $r$ von der Beschaffenheit, dass die Producte $x_{1}z^{-r}, x_{2}z^{-r}, \ldots x_{n}z^{-r}$ für $z = \infty$ alle endlich sind, und mindestens eines von ihnen, etwa $x_{1}z^{-r}$ von Null verschieden; dann enthält aber die Gleichung
\[ \eta z^{-r} = x_{1}z^{-r}\varrho_{1}+\cdots+x_{n}z^{-r}\varrho_{n} \]
den Widerspruch, dass im Punkte $\mathfrak{P}_{1}$ die linke Seite und alle Glieder der rechten Seite mit Ausnahme des ersten verschwinden.

\origpage{287}

Hieraus ergiebt sich zugleich, wenn man $\eta = 0$ setzt, dass die Functionen $\varrho_{1}, \varrho_{2}, \ldots \varrho_{n}$ eine Basis von $\Omega$ bilden. Setzt man daher, indem man mit $x_{\iota,\iota'}$ rationale Functionen von $z$ bezeichnet,
\[ \frac{d\varpi}{dz}\varrho_{\iota} = x_{\iota,1}\varrho_{1}+x_{\iota,2}\varrho_{2}+\cdots+x_{\iota,n}\varrho_{n}, \qquad {\scriptstyle (\iota = 1,\, 2,\, \ldots\, n)} \tag{4.} \]
so ist (§~2)
\[ S\left(\frac{d\varpi}{dz}\right) = x_{1,1}+x_{2,2}+\cdots+x_{n,n}. \tag{5.} \]
Nun ist, wie aus (3.) hervorgeht, $z^{-m+2}\frac{d\varpi}{dz}\varrho_{\iota}$ für $z = \infty$ endlich und daraus ergiebt sich nach der soeben bewiesenen Eigenschaft der Functionen $\varrho$, dass auch
\[ z^{-m+2}x_{\iota,\iota'} \]
für $z = \infty$ endlich sind. Nun sind z.\ B.\ in dem Punkt $\mathfrak{P}_{2}$ die Functionen $\varrho_{1}, \varrho_{3}, \ldots \varrho_{n}$ unendlich klein in der $m^{\text{ten}}$ Ordnung, während $\varrho_{2}$ dort endlich und von Null verschieden ist. Daher werden in $\mathfrak{P}_{2}$ die Functionen
\[ z\frac{d\varpi}{dz}\varrho_{1}, \quad zx_{1,1}\varrho_{1}, \quad zx_{1,3}\varrho_{3}, \quad \ldots \quad zx_{1,n}\varrho_{n} \]
alle verschwinden, und es muss mithin auch $zx_{1,2}$ für $z = \infty$ verschwinden. Das Gleiche folgt für $zx_{1,3}, \ldots zx_{1,n}$ und allgemein für $zx_{\iota,\iota'}$, sobald $\iota$, $\iota'$ von einander verschieden sind. Daher wird $z^{2}x_{\iota,\iota'}$ für $z = \infty$ endlich sein.

Setzt man nun, indem man $x_{\iota}$ eine neue rationale Function bedeuten lässt,
\[ x_{\iota,\iota} = a_{m-2}^{(\iota)}z^{m-2}+a_{m-3}^{(\iota)}z^{m-3}+\cdots+a_{-1}^{(\iota)}z^{-1}+x_{\iota}z^{-2}, \tag{6.} \]
so folgt aus (3.)
\[ x_{\iota,\iota}-\frac{d\varpi}{dz} = z^{-2}(x_{\iota}-\eta^{(\iota)}), \]
und aus (4.)
\[ (\eta^{(\iota)}-x_{\iota})\varrho_{\iota} = z^{2}x_{\iota,1}\varrho_{1}+\cdots+z^{2}x_{\iota,\iota-1}\varrho_{\iota-1}+z^{2}x_{\iota,\iota+1}\varrho_{\iota+1}+\cdots+z^{2}x_{\iota,n}\varrho_{n}. \]
Da nun in $\mathfrak{P}_{\iota}$ $\eta^{(\iota)}$ endlich und $\varrho_{\iota}$ von Null verschieden, ferner alle Glieder der rechten Seite Null sind, so folgt, dass auch $x_{\iota}$ im Punkte $\mathfrak{P}_{\iota}$, und mithin, da es rational ist, für $z = \infty$ endlich ist. Aus (5.) und (6.) ergiebt sich dann
\[ S\left(\frac{d\varpi}{dz}\right) = \overset{\iota}{\Sigma} a_{m-2}^{(\iota)}z^{m-2}+\overset{\iota}{\Sigma} a_{m-3}^{(\iota)}z^{m-3}+\cdots+\overset{\iota}{\Sigma} a_{-1}^{(\iota)}z^{-1}+\overset{\iota}{\Sigma} x_{\iota}z^{-2}. \tag{7.} \]
Nun ist aber andererseits, wenn wieder $\mathfrak{U}$ das Untereck, $\mathfrak{Z}$ das Verzweigungspolygon von $z$ ist:
\[ \frac{d\varpi}{dz} = \frac{\mathfrak{U}^{2}\mathfrak{A}}{\mathfrak{Z}\mathfrak{B}}, \]

\origpage{288}
und $\mathfrak{B}$ enthält keinen Punkt, der nicht auch in $\mathfrak{U}$ enthalten ist. Daraus ergiebt sich wie in §~26, dass $\frac{d\varpi}{dz}$, als Function von $z$ aufgefasst, eine Function des zu $\mathfrak{o}$ complementären Moduls $\mathfrak{e}$ ist, und mithin ist
\[ S\left(\frac{d\varpi}{dz}\right) \]
eine \emph{ganze} rationale Function von $z$ (§~11, 4.). Beachtet man dies, so folgt aus (7.) $\overset{\iota}{\Sigma} x_{\iota} = 0$ und ferner der zu beweisende Satz
\[ \overset{(\iota)}{\Sigma} a_{-1}^{(\iota)} = 0. \]

Wir können diesem Satze auch den folgenden Ausdruck geben: Das Residuum eines Differentials zweiter Gattung $dt_{(\mathfrak{P}^{r})}$ in Bezug auf den Punkt $\mathfrak{P}$ ist Null.

Die Residuen eines Integrals dritter Gattung $d\pi_{(\mathfrak{P}_{1}, \mathfrak{P}_{2})}$ in Bezug auf $\mathfrak{P}_{1}$, $\mathfrak{P}_{2}$ sind einander gleich und entgegengesetzt, und sicher von Null verschieden, da sonst $d\pi$ ein Differential erster Gattung sein würde.

Aus diesen Bemerkungen ergiebt sich noch mittelst 4., dass ein eigentliches Differential $d\sigma$, in der Normalform dargestellt, kein Differential dritter Gattung enthalten kann. Es verdient ferner erwähnt zu werden, dass die Residuen des \emph{logarithmischen} Differentials $\frac{d\sigma}{\sigma}$ ganze Zahlen, nämlich die Ordnungszahlen der Function $\sigma$ sind (zufolge §~23).

\subsection*{§~33. Relationen zwischen Differentialen erster und zweiter Gattung.}

1. Es sei $\sigma$ eine Function in $\Omega$ mit dem Untereck
\[ \mathfrak{B}' = \mathfrak{P}_{1}^{m_{1}-1}\mathfrak{P}_{2}^{m_{2}-1}\ldots \qquad {\scriptstyle (m_{1},\, m_{2},\, \ldots\, \geqq 2)} \]
und dem Verzweigungspolygon (§~16)
\[ \mathfrak{S} = \mathfrak{S}'\mathfrak{P}_{1}^{m_{1}-2}\mathfrak{P}_{2}^{m_{2}-2}\ldots, \]
worin $\mathfrak{S}'$ durch die als verschieden vorausgesetzten Punkte $\mathfrak{P}_{1}, \mathfrak{P}_{2}, \ldots$ nicht theilbar ist. Demnach ist in der symbolischen Bezeichnung von §~25 das eigentliche Differential
\[ d\sigma = \frac{\mathfrak{S}}{\mathfrak{B}'^{2}} = \frac{\mathfrak{S}'}{\mathfrak{P}_{1}^{m_{1}}\mathfrak{P}_{2}^{m_{2}}\ldots}, \]
woraus zunächst hervorgeht, \emph{dass ein eigentliches Differential niemals von der ersten Gattung sein kann.}

\origpage{289}

2. Das eigentliche Differential $d\sigma$, welches in seiner Darstellung durch die Normalform nur Differentiale erster und zweiter Gattung enthalten kann, gehört zu der Schaar derjenigen Differentiale, deren Untereck
\[ \mathfrak{B} = \mathfrak{P}_{1}^{m_{1}}\mathfrak{P}_{2}^{m_{2}}\cdots = \mathfrak{B}'\mathfrak{P}_{1}\mathfrak{P}_{2}\ldots \]
ist. Umgekehrt wird man also auch in einer solchen Schaar, vorausgesetzt dass $m_{1}, m_{2}, \ldots \geqq 2$ sind, und dass $\mathfrak{B}'$ zu einer \emph{eigentlichen Polygonclasse} gehört, stets mindestens ein eigentliches Differential $d\sigma$ finden. Denn dazu ist nach 1. nur erforderlich, dass in $\Omega$ eine Function $\sigma$ mit dem Untereck $\mathfrak{B}'$ existirt.

3. Hieraus ergiebt sich nun der folgende wichtige Satz. \emph{Alle Differentiale zweiter Gattung lassen sich linear mit constanten Coefficienten darstellen durch $p$ besondere passend gewählte Differentiale zweiter Gattung, durch Differentiale erster Gattung und durch eigentliche Differentiale.}

Um dies einzusehen, wähle man ein beliebiges Polygon zweiter Gattung $\mathfrak{A}$ von der Ordnung $p$. Ist nun $\mathfrak{P}$ ein beliebiger Punkt, $r$ ein positiver Exponent, so ist das Polygon $\mathfrak{A}\mathfrak{P}^{r}$ gleichfalls von der zweiten Gattung, und folglich kann der Theiler $\mathfrak{M}$ der zugehörigen Classe nicht durch $\mathfrak{P}$ theilbar sein, weil sonst $\mathfrak{A}\mathfrak{P}^{r-1}$, also auch $\mathfrak{A}$ ein Polygon erster Gattung wäre (§~30, 4.). Setzt man daher
\[ \mathfrak{A}\mathfrak{P}^{r} = \mathfrak{M}\mathfrak{B}', \]
so wird $\mathfrak{P}$ nicht in $\mathfrak{M}$ aufgehen, und folglich enthält $\mathfrak{B}'$ den Factor $\mathfrak{P}$ genau $r$-mal öfter als $\mathfrak{A}$. Zugleich gehört $\mathfrak{B}'$ in eine eigentliche Classe. Ist nun
\[ \mathfrak{B}' = \mathfrak{P}^{m+r}\mathfrak{P}'^{m'}\mathfrak{P}''^{m''}\ldots, \]
so gehen die Punktpotenzen $\mathfrak{P}^{m}$, $\mathfrak{P}'^{m'}$, $\mathfrak{P}''^{m''}$, \ldots\ alle in $\mathfrak{A}$ auf. Setzen wir also
\[ \mathfrak{B} = \mathfrak{P}^{m+r+1}\mathfrak{P}'^{m'+1}\mathfrak{P}''^{m''+1}\ldots = \mathfrak{B}'\mathfrak{P}\mathfrak{P}'\mathfrak{P}''\ldots, \]
so existirt nach 2. in der zu dem Untereck $\mathfrak{B}$ gehörigen Differentialschaar gewiss ein eigentliches Differential $d\sigma$. Die Darstellung desselben durch die Normalform enthält sicher das Differential
\[ dt_{(\mathfrak{P}^{m+r})} \tag{1.} \]
und ausserdem alle oder einige der Differentiale
\[ \left\{ \begin{matrix}
dt_{(\mathfrak{P})}, & dt_{(\mathfrak{P}^{2})}, & \ldots & dt_{(\mathfrak{P}^{m})}\ldots dt_{(\mathfrak{P}^{m+r-1})}, \\
dt_{(\mathfrak{P}')}, & dt_{(\mathfrak{P}'^{2})}, & \ldots & dt_{(\mathfrak{P}'^{m'})}, \\
dt_{(\mathfrak{P}'')}, & dt_{(\mathfrak{P}''^{2})}, & \ldots & dt_{(\mathfrak{P}''^{m''})} \\
\cdots & \cdots & \cdots & \cdots
\end{matrix} \right. \tag{2.} \]

\origpage{290}
nebst Differentialen erster Gattung. Es lässt sich also das Differential (1.) linear und mit constanten Coefficienten durch (2.), durch Differentiale erster Gattung und durch $d\sigma$ ausdrücken.

Ist daher das $p$-Eck zweiter Gattung
\[ \mathfrak{A} = \mathfrak{P}_{1}^{m_{1}}\mathfrak{P}_{2}^{m_{2}}\ldots \]
so erkennt man durch wiederholte Anwendung des hier beschriebenen Verfahrens, dass alle Differentiale zweiter Gattung in der Weise, wie unser Satz es ausspricht, darstellbar sind durch die $p$ Differentiale
\[ \left\{ \begin{matrix}
dt_{(\mathfrak{P}_{1})}\ldots dt_{(\mathfrak{P}_{1}^{m_{1}})}, \\
dt_{(\mathfrak{P}_{2})}\ldots dt_{(\mathfrak{P}_{2}^{m_{2}})}, \\
\cdots
\end{matrix} \right. \tag{3.} \]

Braunschweig und Königsberg i.\ Pr., im October 1880.

\end{document}
